%&preformat-disser
\RequirePackage[l2tabu,orthodox]{nag} % Раскомментировав, можно в логе получать рекомендации относительно правильного использования пакетов и предупреждения об устаревших и нерекомендуемых пакетах
% Формат А4, 14pt (ГОСТ Р 7.0.11-2011, 5.3.6)
\documentclass[a4paper,14pt,oneside,openany]{memoir}

\input{common/setup}            % общие настройки шаблона
%%% Проверка используемого TeX-движка %%%
\newif\ifxetexorluatex   % определяем новый условный оператор (http://tex.stackexchange.com/a/47579)
\ifxetex
    \xetexorluatextrue
\else
    \ifluatex
        \xetexorluatextrue
    \else
        \xetexorluatexfalse
    \fi
\fi

\newif\ifsynopsis           % Условие, проверяющее, что документ --- автореферат

\usepackage{etoolbox}[2015/08/02]   % Для продвинутой проверки разных условий
\providebool{presentation}

\usepackage{comment}    % Позволяет убирать блоки текста (добавляет
                        % окружение comment и команду \excludecomment)

%%% Поля и разметка страницы %%%
\usepackage{pdflscape}  % Для включения альбомных страниц
\usepackage{geometry}   % Для последующего задания полей

%%% Математические пакеты %%%
\usepackage{amsthm,amsmath,amscd}   % Математические дополнения от AMS
\usepackage{amsfonts,amssymb}       % Математические дополнения от AMS
\usepackage{mathtools}              % Добавляет окружение multlined
\usepackage{xfrac}                  % Красивые дроби
\usepackage[
    locale = DE,
    list-separator       = {;\,},
    list-final-separator = {;\,},
    list-pair-separator  = {;\,},
    list-units           = single,
    range-units          = single,
    range-phrase={\text{\ensuremath{-}}},
    % quotient-mode        = fraction, % красивые дроби могут не соответствовать ГОСТ
    fraction-function    = \sfrac,
    separate-uncertainty,
    ]{siunitx}[=v2]                 % Размерности SI
\sisetup{inter-unit-product = \ensuremath{{}\cdot{}}}

% Кириллица в нумерации subequations
% Для правильной работы требуется выполнение сразу после загрузки пакетов
\patchcmd{\subequations}{\def\theequation{\theparentequation\alph{equation}}}
{\def\theequation{\theparentequation\asbuk{equation}}}
{\typeout{subequations patched}}{\typeout{subequations not patched}}

%%%% Установки для размера шрифта 14 pt %%%%
%% Формирование переменных и констант для сравнения (один раз для всех подключаемых файлов)%%
%% должно располагаться до вызова пакета fontspec или polyglossia, потому что они сбивают его работу
\newlength{\curtextsize}
\newlength{\bigtextsize}
\setlength{\bigtextsize}{13.9pt}

\makeatletter
%\show\f@size    % неплохо для отслеживания, но вызывает стопорение процесса,
                 % если документ компилируется без команды  -interaction=nonstopmode
\setlength{\curtextsize}{\f@size pt}
\makeatother

%%% Кодировки и шрифты %%%
\ifxetexorluatex
    \ifpresentation
        \providecommand*\autodot{} % quick fix for polyglossia 1.50
    \fi
    \PassOptionsToPackage{no-math}{fontspec}    % https://tex.stackexchange.com/a/26295/104425
    \usepackage{polyglossia}[2014/05/21]        % Поддержка многоязычности
                                        % (fontspec подгружается автоматически)
\else
   %%% Решение проблемы копирования текста в буфер кракозябрами
    \ifnumequal{\value{usealtfont}}{0}{}{
        \input glyphtounicode.tex
        \input glyphtounicode-cmr.tex %from pdfx package
        \pdfgentounicode=1
    }
    \usepackage{cmap}   % Улучшенный поиск русских слов в полученном pdf-файле
    \ifnumequal{\value{usealtfont}}{2}{}{
        \defaulthyphenchar=127  % Если стоит до fontenc, то переносы
                                % не впишутся в выделяемый текст при
                                % копировании его в буфер обмена
    }
    \usepackage{textcomp}
    \usepackage[T1,T2A]{fontenc}                    % Поддержка русских букв
    \ifnumequal{\value{usealtfont}}{1}{% Используется pscyr, при наличии
        \IfFileExists{pscyr.sty}{\usepackage{pscyr}}{}  % Подключение pscyr
    }{}
    \usepackage[utf8]{inputenc}[2014/04/30]         % Кодировка utf8
    \usepackage[english, russian]{babel}[2014/03/24]% Языки: русский, английский
    \makeatletter\AtBeginDocument{\let\@elt\relax}\makeatother % babel 3.40 fix
    \ifnumequal{\value{usealtfont}}{2}{
        % http://dxdy.ru/post1238763.html#p1238763
        \usepackage[scaled=0.914]{XCharter}[2017/12/19] % Подключение русифицированных шрифтов XCharter
        \usepackage[charter, vvarbb, scaled=1.048]{newtxmath}[2017/12/14]
        \ifpresentation
        \else
            \setDisplayskipStretch{-0.078}
        \fi
    }{}
\fi

%%% Оформление абзацев %%%
\ifpresentation
\else
    \indentafterchapter     % Красная строка после заголовков типа chapter
    \usepackage{indentfirst}
\fi

%%% Цвета %%%
\ifpresentation
\else
    \usepackage[dvipsnames, table, hyperref]{xcolor} % Совместимо с tikz
\fi

%%% Таблицы %%%
\usepackage{longtable,ltcaption} % Длинные таблицы
\usepackage{multirow,makecell}   % Улучшенное форматирование таблиц
\usepackage{tabu, tabulary}      % таблицы с автоматически подбирающейся
                                 % шириной столбцов (tabu обязательно
                                 % до hyperref вызывать)
\makeatletter
%https://github.com/tabu-issues-for-future-maintainer/tabu/issues/26
\@ifpackagelater{longtable}{2020/02/07}{
\def\tabuendlongtrial{%
    \LT@echunk  \global\setbox\LT@gbox \hbox{\unhbox\LT@gbox}\kern\wd\LT@gbox
                \LT@get@widths
}%
}{}
\makeatother

\usepackage{threeparttable}      % автоматический подгон ширины подписи таблицы

%%% Общее форматирование
\usepackage{soulutf8}% Поддержка переносоустойчивых подчёркиваний и зачёркиваний
\usepackage{icomma}  % Запятая в десятичных дробях

%%% Оптимизация расстановки переносов и длины последней строки абзаца
\IfFileExists{impnattypo.sty}{% проверка установленности пакета impnattypo
    \ifluatex
        \ifnumequal{\value{draft}}{1}{% Черновик
            \usepackage[hyphenation, lastparline, nosingleletter, homeoarchy,
            rivers, draft]{impnattypo}
        }{% Чистовик
            \usepackage[hyphenation, lastparline, nosingleletter]{impnattypo}
        }
    \else
        \usepackage[hyphenation, lastparline]{impnattypo}
    \fi
}{}

%% Векторная графика

\usepackage{tikz}                   % Продвинутый пакет векторной графики
\usetikzlibrary{chains}             % Для примера tikz рисунка
\usetikzlibrary{shapes.geometric}   % Для примера tikz рисунка
\usetikzlibrary{shapes.symbols}     % Для примера tikz рисунка
\usetikzlibrary{arrows}             % Для примера tikz рисунка
\usetikzlibrary{decorations.pathmorphing}

%%% Гиперссылки %%%
\ifxetexorluatex
    \let\CYRDZE\relax
\fi
\usepackage{hyperref}[2012/11/06]

%%% Изображения %%%
\usepackage{graphicx}[2014/04/25]   % Подключаем пакет работы с графикой
\usepackage{caption}                % Подписи рисунков и таблиц
\usepackage{subcaption}             % Подписи подрисунков и подтаблиц
\usepackage{pdfpages}               % Добавление внешних pdf файлов

%%% Счётчики %%%
\usepackage{aliascnt}
\usepackage[figure,table]{totalcount}   % Счётчик рисунков и таблиц
\usepackage{totcount}   % Пакет создания счётчиков на основе последнего номера
                        % подсчитываемого элемента (может требовать дважды
                        % компилировать документ)
\usepackage{totpages}   % Счётчик страниц, совместимый с hyperref (ссылается
                        % на номер последней страницы). Желательно ставить
                        % последним пакетом в преамбуле

%%% Продвинутое управление групповыми ссылками (пока только формулами) %%%
\ifpresentation
\else
    \usepackage[russian]{cleveref} % cleveref имеет сложности со считыванием
    % языка из babel. Такое решение русификации вывода выбрано вместо
    % определения в documentclass из опасности что-то лишнее передать во все
    % остальные пакеты, включая библиографию.

    % Добавление возможности использования пробелов в \labelcref
    % https://tex.stackexchange.com/a/340502/104425
    \usepackage{kvsetkeys}
    \makeatletter
    \let\org@@cref\@cref
    \renewcommand*{\@cref}[2]{%
        \edef\process@me{%
            \noexpand\org@@cref{#1}{\zap@space#2 \@empty}%
        }\process@me
    }
    \makeatother
\fi

\usepackage{placeins} % для \FloatBarrier

\ifnumequal{\value{draft}}{1}{% Черновик
    \usepackage[firstpage]{draftwatermark}
    \SetWatermarkText{DRAFT}
    \SetWatermarkFontSize{14pt}
    \SetWatermarkScale{15}
    \SetWatermarkAngle{45}
}{}

%%% Цитата, не приводимая в автореферате:
% возможно, актуальна только для biblatex
%\newcommand{\citeinsynopsis}[1]{\ifsynopsis\else ~\cite{#1} \fi}

% если текущий процесс запущен библиотекой tikz-external, то прекомпиляция должна быть включена
\ifdefined\tikzexternalrealjob
    \setcounter{imgprecompile}{1}
\fi

\ifnumequal{\value{imgprecompile}}{1}{% Только если у нас включена предкомпиляция
    \usetikzlibrary{external}   % подключение возможности предкомпиляции
    \tikzexternalize[prefix=images/cache/,optimize command away=\includepdf] % activate! % здесь можно указать отдельную папку для скомпилированных файлов
    \ifxetex
        \tikzset{external/up to date check={diff}}
    \fi
}{}
         % Пакеты общие для диссертации и автореферата
\synopsisfalse                      % Этот документ --- не автореферат
\input{Dissertation/dispackages}    % Пакеты для диссертации
\usepackage{fr-longtable}    %ради \endlasthead

% Листинги с исходным кодом программ
\usepackage{fancyvrb}
\usepackage{listings}
\lccode`\~=0\relax %Без этого хака из-за особенностей пакета listings перестают работать конструкции с \MakeLowercase и т. п. в (xe|lua)latex

% Русская традиция начертания греческих букв
\usepackage{upgreek} % прямые греческие ради русской традиции

%%% Микротипографика
%\ifnumequal{\value{draft}}{0}{% Только если у нас режим чистовика
%    \usepackage[final, babel, shrink=45]{microtype}[2016/05/14] % улучшает представление букв и слов в строках, может помочь при наличии отдельно висящих слов
%}{}

% Отметка о версии черновика на каждой странице
% Чтобы работало надо в своей локальной копии по инструкции
% https://www.ctan.org/pkg/gitinfo2 создать небходимые файлы в папке
% ./git/hooks
% If you’re familiar with tweaking git, you can probably work it out for
% yourself. If not, I suggest you follow these steps:
% 1. First, you need a git repository and working tree. For this example,
% let’s suppose that the root of the working tree is in ~/compsci
% 2. Copy the file post-xxx-sample.txt (which is in the same folder of
% your TEX distribution as this pdf) into the git hooks directory in your
% working copy. In our example case, you should end up with a file called
% ~/compsci/.git/hooks/post-checkout
% 3. If you’re using a unix-like system, don’t forget to make the file executable.
% Just how you do this is outside the scope of this manual, but one
% possible way is with commands such as this:
% chmod g+x post-checkout.
% 4. Test your setup with “git checkout master” (or another suitable branch
% name). This should generate copies of gitHeadInfo.gin in the directories
% you intended.
% 5. Now make two more copies of this file in the same directory (hooks),
% calling them post-commit and post-merge, and you’re done. As before,
% users of unix-like systems should ensure these files are marked as
% executable.
\ifnumequal{\value{draft}}{1}{% Черновик
   \IfFileExists{.git/gitHeadInfo.gin}{
      \usepackage[mark,pcount]{gitinfo2}
      \renewcommand{\gitMark}{rev.\gitAbbrevHash\quad\gitCommitterEmail\quad\gitAuthorIsoDate}
      \renewcommand{\gitMarkFormat}{\rmfamily\color{Gray}\small\bfseries}
   }{}
}{}

\usepackage{hhline}   % Пакеты для специфических пользовательских задач

\input{Dissertation/setup}      % Упрощённые настройки шаблона

% Новые переменные, которые могут использоваться во всём проекте
% ГОСТ 7.0.11-2011
% 9.2 Оформление текста автореферата диссертации
% 9.2.1 Общая характеристика работы включает в себя следующие основные структурные
% элементы:
% актуальность темы исследования;
\newcommand{\actualityTXT}{Актуальность темы.}
% степень ее разработанности;
\newcommand{\progressTXT}{Степень разработанности темы.}
% цели и задачи;
\newcommand{\aimTXT}{Целью}
\newcommand{\tasksTXT}{задачи}
% научную новизну;
\newcommand{\noveltyTXT}{Научная новизна:}
% теоретическую и практическую значимость работы;
%\newcommand{\influenceTXT}{Теоретическая и практическая значимость}
% или чаще используют просто
\newcommand{\influenceTXT}{Практическая значимость}
% методологию и методы исследования;
\newcommand{\methodsTXT}{Методология и методы исследования.}
% положения, выносимые на защиту;
\newcommand{\defpositionsTXT}{Основные положения, выносимые на~защиту:}
% степень достоверности и апробацию результатов.
\newcommand{\reliabilityTXT}{Достоверность}
\newcommand{\probationTXT}{Апробация работы.}

\newcommand{\contributionTXT}{Личный вклад.}
\newcommand{\publicationsTXT}{Публикации.}


%%% Заголовки библиографии:

% для автореферата:
\newcommand{\bibtitleauthor}{Публикации автора по теме диссертации}

% для стиля библиографии `\insertbiblioauthorgrouped`
\newcommand{\bibtitleauthorvak}{В изданиях из списка ВАК РФ}
\newcommand{\bibtitleauthorscopus}{В изданиях, входящих в международную базу цитирования Scopus}
\newcommand{\bibtitleauthorwos}{В изданиях, входящих в международную базу цитирования Web of Science}
\newcommand{\bibtitleauthorother}{В прочих изданиях}
\newcommand{\bibtitleauthorconf}{В сборниках трудов конференций}
\newcommand{\bibtitleauthorpatent}{Зарегистрированные патенты}
\newcommand{\bibtitleauthorprogram}{Зарегистрированные программы для ЭВМ}

% для стиля библиографии `\insertbiblioauthorimportant`:
\newcommand{\bibtitleauthorimportant}{Наиболее значимые \protect\MakeLowercase\bibtitleauthor}

% для списка литературы в диссертации и списка чужих работ в автореферате:
\newcommand{\bibtitlefull}{Bibliography} % (ГОСТ Р 7.0.11-2011, 4)
         % Новые переменные, для всего проекта

%%% Основные сведения %%%
\newcommand{\thesisAuthorLastName}{Савченко}
\newcommand{\thesisAuthorOtherNames}{Максим Алексеевич}
\newcommand{\thesisAuthorInitials}{М.\,А.}
\newcommand{\thesisAuthor}             % Диссертация, ФИО автора
{%
    \texorpdfstring{% \texorpdfstring takes two arguments and uses the first for (La)TeX and the second for pdf
        \thesisAuthorLastName~\thesisAuthorOtherNames% так будет отображаться на титульном листе или в тексте, где будет использоваться переменная
    }{%
        \thesisAuthorLastName, \thesisAuthorOtherNames% эта запись для свойств pdf-файла. В таком виде, если pdf будет обработан программами для сбора библиографических сведений, будет правильно представлена фамилия.
    }
}
\newcommand{\thesisAuthorShort}        % Диссертация, ФИО автора инициалами
{\thesisAuthorInitials~\thesisAuthorLastName}
%\newcommand{\thesisUdk}                % Диссертация, УДК
%{\fixme{xxx.xxx}}
\newcommand{\thesisTitle}              % Диссертация, название
{Влияние дополнительной информационной асимметрии на решения неантагонистических игр}
\newcommand{\thesisSpecialtyNumber}    % Диссертация, специальность, номер
{01.01.09}
\newcommand{\thesisSpecialtyTitle}     % Диссертация, специальность, название (название взято с сайта ВАК для примера)
{Дискретная математика и математическая кибернетика}
%% \newcommand{\thesisSpecialtyTwoNumber} % Диссертация, вторая специальность, номер
%% {\fixme{XX.XX.XX}}
%% \newcommand{\thesisSpecialtyTwoTitle}  % Диссертация, вторая специальность, название
%% {\fixme{Теория и~методика физического воспитания, спортивной тренировки,
%% оздоровительной и~адаптивной физической культуры}}
\newcommand{\thesisDegree}             % Диссертация, ученая степень
{кандидата физико-математических наук}
\newcommand{\thesisDegreeShort}        % Диссертация, ученая степень, краткая запись
{канд. физ.-мат. наук}
\newcommand{\thesisCity}               % Диссертация, город написания диссертации
{Москва}
\newcommand{\thesisYear}               % Диссертация, год написания диссертации
{\the\year}
\newcommand{\thesisOrganization}       % Диссертация, организация
{Федеральное государственное бюджетное образовательное учреждение высшего образования <<Московский государственный университет имени М. В. Ломоносова>>}
\newcommand{\thesisOrganizationShort}  % Диссертация, краткое название организации для доклада
{МГУ}

\newcommand{\thesisInOrganization}     % Диссертация, организация в предложном падеже: Работа выполнена в ...
{факультете вычислительной математики и кибернетики (ВМК) Московского государственного университета имени М. В. Ломоносова}

%% \newcommand{\supervisorDead}{}           % Рисовать рамку вокруг фамилии
\newcommand{\supervisorFio}              % Научный руководитель, ФИО
{Васин Александр Алексеевич}
\newcommand{\supervisorRegalia}          % Научный руководитель, регалии
{доктор физико-математических наук, профессор}
\newcommand{\supervisorFioShort}         % Научный руководитель, ФИО
{А.\,А.~Васин}
\newcommand{\supervisorRegaliaShort}     % Научный руководитель, регалии
{д-р~физ.-мат.~наук,~проф.}

%% \newcommand{\supervisorTwoDead}{}        % Рисовать рамку вокруг фамилии
%% \newcommand{\supervisorTwoFio}           % Второй научный руководитель, ФИО
%% {\fixme{Фамилия Имя Отчество}}
%% \newcommand{\supervisorTwoRegalia}       % Второй научный руководитель, регалии
%% {\fixme{уч. степень, уч. звание}}
%% \newcommand{\supervisorTwoFioShort}      % Второй научный руководитель, ФИО
%% {\fixme{И.\,О.~Фамилия}}
%% \newcommand{\supervisorTwoRegaliaShort}  % Второй научный руководитель, регалии
%% {\fixme{уч.~ст.,~уч.~зв.}}

\newcommand{\opponentOneFio}           % Оппонент 1, ФИО
{\fixme{Фамилия Имя Отчество}}
\newcommand{\opponentOneRegalia}       % Оппонент 1, регалии
{\fixme{доктор физико-математических наук, профессор}}
\newcommand{\opponentOneJobPlace}      % Оппонент 1, место работы
{\fixme{Не очень длинное название для места работы}}
\newcommand{\opponentOneJobPost}       % Оппонент 1, должность
{\fixme{старший научный сотрудник}}

\newcommand{\opponentTwoFio}           % Оппонент 2, ФИО
{\fixme{Фамилия Имя Отчество}}
\newcommand{\opponentTwoRegalia}       % Оппонент 2, регалии
{\fixme{кандидат физико-математических наук}}
\newcommand{\opponentTwoJobPlace}      % Оппонент 2, место работы
{\fixme{Основное место работы c длинным длинным длинным длинным названием}}
\newcommand{\opponentTwoJobPost}       % Оппонент 2, должность
{\fixme{старший научный сотрудник}}

%% \newcommand{\opponentThreeFio}         % Оппонент 3, ФИО
%% {\fixme{Фамилия Имя Отчество}}
%% \newcommand{\opponentThreeRegalia}     % Оппонент 3, регалии
%% {\fixme{кандидат физико-математических наук}}
%% \newcommand{\opponentThreeJobPlace}    % Оппонент 3, место работы
%% {\fixme{Основное место работы c длинным длинным длинным длинным названием}}
%% \newcommand{\opponentThreeJobPost}     % Оппонент 3, должность
%% {\fixme{старший научный сотрудник}}

\newcommand{\leadingOrganizationTitle} % Ведущая организация, дополнительные строки. Удалить, чтобы не отображать в автореферате
{\fixme{Федеральное государственное бюджетное образовательное учреждение высшего
профессионального образования с~длинным длинным длинным длинным названием}}

\newcommand{\defenseDate}              % Защита, дата
{\fixme{DD mmmmmmmm YYYY~г.~в~XX часов}}
\newcommand{\defenseCouncilNumber}     % Защита, номер диссертационного совета
{\fixme{Д\,123.456.78}}
\newcommand{\defenseCouncilTitle}      % Защита, учреждение диссертационного совета
{\fixme{Название учреждения}}
\newcommand{\defenseCouncilAddress}    % Защита, адрес учреждение диссертационного совета
{\fixme{Адрес}}
\newcommand{\defenseCouncilPhone}      % Телефон для справок
{\fixme{+7~(0000)~00-00-00}}

\newcommand{\defenseSecretaryFio}      % Секретарь диссертационного совета, ФИО
{\fixme{Фамилия Имя Отчество}}
\newcommand{\defenseSecretaryRegalia}  % Секретарь диссертационного совета, регалии
{\fixme{д-р~физ.-мат. наук}}            % Для сокращений есть ГОСТы, например: ГОСТ Р 7.0.12-2011 + http://base.garant.ru/179724/#block_30000

\newcommand{\synopsisLibrary}          % Автореферат, название библиотеки
{\fixme{Название библиотеки}}
\newcommand{\synopsisDate}             % Автореферат, дата рассылки
{\fixme{DD mmmmmmmm}\the\year~года}

% To avoid conflict with beamer class use \providecommand
\providecommand{\keywords}%            % Ключевые слова для метаданных PDF диссертации и автореферата
{}
             % Основные сведения
\input{common/fonts}            % Определение шрифтов (частичное)
\input{common/styles}           % Стили общие для диссертации и автореферата
\theoremstyle{definition}
\newtheorem{definition}{Определение}[section]

%%% Переопределение именований, если иначе не сработает %%%
%\gappto\captionsrussian{
%    \renewcommand{\chaptername}{Глава}
%    \renewcommand{\appendixname}{Приложение} % (ГОСТ Р 7.0.11-2011, 5.7)
%}

%%% Изображения %%%
\graphicspath{{images/}{Dissertation/images/}}         % Пути к изображениям

%%% Интервалы %%%
%% По ГОСТ Р 7.0.11-2011, пункту 5.3.6 требуется полуторный интервал
%% Реализация средствами класса (на основе setspace) ближе к типографской классике.
%% И правит сразу и в таблицах (если со звёздочкой)
%\DoubleSpacing*     % Двойной интервал
\OnehalfSpacing*    % Полуторный интервал
%\setSpacing{1.42}   % Полуторный интервал, подобный Ворду (возможно, стоит включать вместе с предыдущей строкой)

%%% Макет страницы %%%
% Выставляем значения полей (ГОСТ 7.0.11-2011, 5.3.7)
\geometry{a4paper, top=2cm, bottom=2cm, left=2.5cm, right=1cm, nofoot, nomarginpar} %, heightrounded, showframe
\setlength{\topskip}{0pt}   %размер дополнительного верхнего поля
\setlength{\footskip}{12.3pt} % снимет warning, согласно https://tex.stackexchange.com/a/334346

%%% Выравнивание и переносы %%%
%% http://tex.stackexchange.com/questions/241343/what-is-the-meaning-of-fussy-sloppy-emergencystretch-tolerance-hbadness
%% http://www.latex-community.org/forum/viewtopic.php?p=70342#p70342
\tolerance 1414
\hbadness 1414
\emergencystretch 1.5em % В случае проблем регулировать в первую очередь
\hfuzz 0.3pt
\vfuzz \hfuzz
%\raggedbottom
%\sloppy                 % Избавляемся от переполнений
\clubpenalty=10000      % Запрещаем разрыв страницы после первой строки абзаца
\widowpenalty=10000     % Запрещаем разрыв страницы после последней строки абзаца
\brokenpenalty=4991     % Ограничение на разрыв страницы, если строка заканчивается переносом

%%% Блок управления параметрами для выравнивания заголовков в тексте %%%
\newlength{\otstuplen}
\setlength{\otstuplen}{\theotstup\parindent}
\ifnumequal{\value{headingalign}}{0}{% выравнивание заголовков в тексте
    \newcommand{\hdngalign}{\centering}                % по центру
    \newcommand{\hdngaligni}{}% по центру
    \setlength{\otstuplen}{0pt}
}{%
    \newcommand{\hdngalign}{}                 % по левому краю
    \newcommand{\hdngaligni}{\hspace{\otstuplen}}      % по левому краю
} % В обоих случаях вроде бы без переноса, как и надо (ГОСТ Р 7.0.11-2011, 5.3.5)

%%% Оглавление %%%
\renewcommand{\cftchapterdotsep}{\cftdotsep}                % отбивка точками до номера страницы начала главы/раздела

%% Переносить слова в заголовке не допускается (ГОСТ Р 7.0.11-2011, 5.3.5). Заголовки в оглавлении должны точно повторять заголовки в тексте (ГОСТ Р 7.0.11-2011, 5.2.3). Прямого указания на запрет переносов в оглавлении нет, но по той же логике невнесения искажений в смысл, лучше в оглавлении не переносить:
\setrmarg{2.55em plus1fil}                             %To have the (sectional) titles in the ToC, etc., typeset ragged right with no hyphenation
\renewcommand{\cftchapterpagefont}{\normalfont}        % нежирные номера страниц у глав в оглавлении
\renewcommand{\cftchapterleader}{\cftdotfill{\cftchapterdotsep}}% нежирные точки до номеров страниц у глав в оглавлении
%\renewcommand{\cftchapterfont}{}                       % нежирные названия глав в оглавлении

\ifnumgreater{\value{headingdelim}}{0}{%
    \renewcommand\cftchapteraftersnum{.\space}       % добавляет точку с пробелом после номера раздела в оглавлении
}{}
\ifnumgreater{\value{headingdelim}}{1}{%
    \renewcommand\cftsectionaftersnum{.\space}       % добавляет точку с пробелом после номера подраздела в оглавлении
    \renewcommand\cftsubsectionaftersnum{.\space}    % добавляет точку с пробелом после номера подподраздела в оглавлении
    \renewcommand\cftsubsubsectionaftersnum{.\space} % добавляет точку с пробелом после номера подподподраздела в оглавлении
    \AfterEndPreamble{% без этого polyglossia сама всё переопределяет
        \setsecnumformat{\csname the#1\endcsname.\space}
    }
}{%
    \AfterEndPreamble{% без этого polyglossia сама всё переопределяет
        \setsecnumformat{\csname the#1\endcsname\quad}
    }
}

\renewcommand*{\cftappendixname}{\appendixname\space} % Слово Приложение в оглавлении

%%% Колонтитулы %%%
% Порядковый номер страницы печатают на середине верхнего поля страницы (ГОСТ Р 7.0.11-2011, 5.3.8)
\makeevenhead{plain}{}{\thepage}{}
\makeoddhead{plain}{}{\thepage}{}
\makeevenfoot{plain}{}{}{}
\makeoddfoot{plain}{}{}{}
\pagestyle{plain}

%%% добавить Стр. над номерами страниц в оглавлении
%%% http://tex.stackexchange.com/a/306950
\newif\ifendTOC

\newcommand*{\tocheader}{
\ifnumequal{\value{pgnum}}{1}{%
    \ifendTOC\else\hbox to \linewidth%
      {\noindent{}~\hfill{Стр.}}\par%
      \ifnumless{\value{page}}{3}{}{%
        \vspace{0.5\onelineskip}
      }
      \afterpage{\tocheader}
    \fi%
}{}%
}%

%%% Оформление заголовков глав, разделов, подразделов %%%
%% Работа должна быть выполнена ... размером шрифта 12-14 пунктов (ГОСТ Р 7.0.11-2011, 5.3.8). То есть не должно быть надписей шрифтом более 14. Так и поставим.
%% Эти установки будут давать одинаковый результат независимо от выбора базовым шрифтом 12 пт или 14 пт
\newcommand{\basegostsectionfont}{\fontsize{14pt}{16pt}\selectfont\bfseries}

\makechapterstyle{thesisgost}{%
    \chapterstyle{default}
    \setlength{\beforechapskip}{0pt}
    \setlength{\midchapskip}{0pt}
    \setlength{\afterchapskip}{\theintvl\curtextsize}
    \renewcommand*{\chapnamefont}{\basegostsectionfont}
    \renewcommand*{\chapnumfont}{\basegostsectionfont}
    \renewcommand*{\chaptitlefont}{\basegostsectionfont}
    \renewcommand*{\chapterheadstart}{}
    \ifnumgreater{\value{headingdelim}}{0}{%
        \renewcommand*{\afterchapternum}{.\space}   % добавляет точку с пробелом после номера раздела
    }{%
        \renewcommand*{\afterchapternum}{\quad}     % добавляет \quad после номера раздела
    }
    \renewcommand*{\printchapternum}{\hdngaligni\hdngalign\chapnumfont \thechapter}
    \renewcommand*{\printchaptername}{}
    \renewcommand*{\printchapternonum}{\hdngaligni\hdngalign}
}

\makeatletter
\makechapterstyle{thesisgostchapname}{%
    \chapterstyle{thesisgost}
    \renewcommand*{\printchapternum}{\chapnumfont \thechapter}
    \renewcommand*{\printchaptername}{\hdngaligni\hdngalign\chapnamefont \@chapapp} %
}
\makeatother

\chapterstyle{thesisgost}

\setsecheadstyle{\basegostsectionfont\hdngalign}
\setsecindent{\otstuplen}

\setsubsecheadstyle{\basegostsectionfont\hdngalign}
\setsubsecindent{\otstuplen}

\setsubsubsecheadstyle{\basegostsectionfont\hdngalign}
\setsubsubsecindent{\otstuplen}

\sethangfrom{\noindent #1} %все заголовки подразделов центрируются с учетом номера, как block

\ifnumequal{\value{chapstyle}}{1}{%
    \chapterstyle{thesisgostchapname}
    \renewcommand*{\cftchaptername}{\chaptername\space} % будет вписано слово Глава перед каждым номером раздела в оглавлении
}{}%

%%% Интервалы между заголовками
\setbeforesecskip{\theintvl\curtextsize}% Заголовки отделяют от текста сверху и снизу тремя интервалами (ГОСТ Р 7.0.11-2011, 5.3.5).
\setaftersecskip{\theintvl\curtextsize}
\setbeforesubsecskip{\theintvl\curtextsize}
\setaftersubsecskip{\theintvl\curtextsize}
\setbeforesubsubsecskip{\theintvl\curtextsize}
\setaftersubsubsecskip{\theintvl\curtextsize}

%%% Вертикальные интервалы глав (\chapter) в оглавлении как и у заголовков
% раскомментировать следующие 2
% \setlength{\cftbeforechapterskip}{0pt plus 0pt}   % ИЛИ эти 2 строки из учебника
% \renewcommand*{\insertchapterspace}{}
% или эту
% \renewcommand*{\cftbeforechapterskip}{0em}


%%% Блок дополнительного управления размерами заголовков
\ifnumequal{\value{headingsize}}{1}{% Пропорциональные заголовки и базовый шрифт 14 пт
    \renewcommand{\basegostsectionfont}{\large\bfseries}
    \renewcommand*{\chapnamefont}{\Large\bfseries}
    \renewcommand*{\chapnumfont}{\Large\bfseries}
    \renewcommand*{\chaptitlefont}{\Large\bfseries}
}{}

%%% Счётчики %%%

%% Упрощённые настройки шаблона диссертации: нумерация формул, таблиц, рисунков
\ifnumequal{\value{contnumeq}}{1}{%
    \counterwithout{equation}{chapter} % Убираем связанность номера формулы с номером главы/раздела
}{}
\ifnumequal{\value{contnumfig}}{1}{%
    \counterwithout{figure}{chapter}   % Убираем связанность номера рисунка с номером главы/раздела
}{}
\ifnumequal{\value{contnumtab}}{1}{%
    \counterwithout{table}{chapter}    % Убираем связанность номера таблицы с номером главы/раздела
}{}

\AfterEndPreamble{
%% регистрируем счётчики в системе totcounter
    \regtotcounter{totalcount@figure}
    \regtotcounter{totalcount@table}       % Если иным способом поставить в преамбуле то ошибка в числе таблиц
    \regtotcounter{TotPages}               % Если иным способом поставить в преамбуле то ошибка в числе страниц
    \newtotcounter{totalappendix}
    \newtotcounter{totalchapter}
}
  % Стили для диссертации
% для вертикального центрирования ячеек в tabulary
\def\zz{\ifx\[$\else\aftergroup\zzz\fi}
%$ \] % <-- чиним подсветку синтаксиса в некоторых редакторах
\def\zzz{\setbox0\lastbox
\dimen0\dimexpr\extrarowheight + \ht0-\dp0\relax
\setbox0\hbox{\raise-.5\dimen0\box0}%
\ht0=\dimexpr\ht0+\extrarowheight\relax
\dp0=\dimexpr\dp0+\extrarowheight\relax
\box0
}

\lstdefinelanguage{Renhanced}%
{keywords={abbreviate,abline,abs,acos,acosh,action,add1,add,%
        aggregate,alias,Alias,alist,all,anova,any,aov,aperm,append,apply,%
        approx,approxfun,apropos,Arg,args,array,arrows,as,asin,asinh,%
        atan,atan2,atanh,attach,attr,attributes,autoload,autoloader,ave,%
        axis,backsolve,barplot,basename,besselI,besselJ,besselK,besselY,%
        beta,binomial,body,box,boxplot,break,browser,bug,builtins,bxp,by,%
        c,C,call,Call,case,cat,category,cbind,ceiling,character,char,%
        charmatch,check,chol,chol2inv,choose,chull,class,close,cm,codes,%
        coef,coefficients,co,col,colnames,colors,colours,commandArgs,%
        comment,complete,complex,conflicts,Conj,contents,contour,%
        contrasts,contr,control,helmert,contrib,convolve,cooks,coords,%
        distance,coplot,cor,cos,cosh,count,fields,cov,covratio,wt,CRAN,%
        create,crossprod,cummax,cummin,cumprod,cumsum,curve,cut,cycle,D,%
        data,dataentry,date,dbeta,dbinom,dcauchy,dchisq,de,debug,%
        debugger,Defunct,default,delay,delete,deltat,demo,de,density,%
        deparse,dependencies,Deprecated,deriv,description,detach,%
        dev2bitmap,dev,cur,deviance,off,prev,,dexp,df,dfbetas,dffits,%
        dgamma,dgeom,dget,dhyper,diag,diff,digamma,dim,dimnames,dir,%
        dirname,dlnorm,dlogis,dnbinom,dnchisq,dnorm,do,dotplot,double,%
        download,dpois,dput,drop,drop1,dsignrank,dt,dummy,dump,dunif,%
        duplicated,dweibull,dwilcox,dyn,edit,eff,effects,eigen,else,%
        emacs,end,environment,env,erase,eval,equal,evalq,example,exists,%
        exit,exp,expand,expression,External,extract,extractAIC,factor,%
        fail,family,fft,file,filled,find,fitted,fivenum,fix,floor,for,%
        For,formals,format,formatC,formula,Fortran,forwardsolve,frame,%
        frequency,ftable,ftable2table,function,gamma,Gamma,gammaCody,%
        gaussian,gc,gcinfo,gctorture,get,getenv,geterrmessage,getOption,%
        getwd,gl,glm,globalenv,gnome,GNOME,graphics,gray,grep,grey,grid,%
        gsub,hasTsp,hat,heat,help,hist,home,hsv,httpclient,I,identify,if,%
        ifelse,Im,image,\%in\%,index,influence,measures,inherits,install,%
        installed,integer,interaction,interactive,Internal,intersect,%
        inverse,invisible,IQR,is,jitter,kappa,kronecker,labels,lapply,%
        layout,lbeta,lchoose,lcm,legend,length,levels,lgamma,library,%
        licence,license,lines,list,lm,load,local,locator,log,log10,log1p,%
        log2,logical,loglin,lower,lowess,ls,lsfit,lsf,ls,machine,Machine,%
        mad,mahalanobis,make,link,margin,match,Math,matlines,mat,matplot,%
        matpoints,matrix,max,mean,median,memory,menu,merge,methods,min,%
        missing,Mod,mode,model,response,mosaicplot,mtext,mvfft,na,nan,%
        names,omit,nargs,nchar,ncol,NCOL,new,next,NextMethod,nextn,%
        nlevels,nlm,noquote,NotYetImplemented,NotYetUsed,nrow,NROW,null,%
        numeric,\%o\%,objects,offset,old,on,Ops,optim,optimise,optimize,%
        options,or,order,ordered,outer,package,packages,page,pairlist,%
        pairs,palette,panel,par,parent,parse,paste,path,pbeta,pbinom,%
        pcauchy,pchisq,pentagamma,persp,pexp,pf,pgamma,pgeom,phyper,pico,%
        pictex,piechart,Platform,plnorm,plogis,plot,pmatch,pmax,pmin,%
        pnbinom,pnchisq,pnorm,points,poisson,poly,polygon,polyroot,pos,%
        postscript,power,ppoints,ppois,predict,preplot,pretty,Primitive,%
        print,prmatrix,proc,prod,profile,proj,prompt,prop,provide,%
        psignrank,ps,pt,ptukey,punif,pweibull,pwilcox,q,qbeta,qbinom,%
        qcauchy,qchisq,qexp,qf,qgamma,qgeom,qhyper,qlnorm,qlogis,qnbinom,%
        qnchisq,qnorm,qpois,qqline,qqnorm,qqplot,qr,Q,qty,qy,qsignrank,%
        qt,qtukey,quantile,quasi,quit,qunif,quote,qweibull,qwilcox,%
        rainbow,range,rank,rbeta,rbind,rbinom,rcauchy,rchisq,Re,read,csv,%
        csv2,fwf,readline,socket,real,Recall,rect,reformulate,regexpr,%
        relevel,remove,rep,repeat,replace,replications,report,require,%
        resid,residuals,restart,return,rev,rexp,rf,rgamma,rgb,rgeom,R,%
        rhyper,rle,rlnorm,rlogis,rm,rnbinom,RNGkind,rnorm,round,row,%
        rownames,rowsum,rpois,rsignrank,rstandard,rstudent,rt,rug,runif,%
        rweibull,rwilcox,sample,sapply,save,scale,scan,scan,screen,sd,se,%
        search,searchpaths,segments,seq,sequence,setdiff,setequal,set,%
        setwd,show,sign,signif,sin,single,sinh,sink,solve,sort,source,%
        spline,splinefun,split,sqrt,stars,start,stat,stem,step,stop,%
        storage,strstrheight,stripplot,strsplit,structure,strwidth,sub,%
        subset,substitute,substr,substring,sum,summary,sunflowerplot,svd,%
        sweep,switch,symbol,symbols,symnum,sys,status,system,t,table,%
        tabulate,tan,tanh,tapply,tempfile,terms,terrain,tetragamma,text,%
        time,title,topo,trace,traceback,transform,tri,trigamma,trunc,try,%
        ts,tsp,typeof,unclass,undebug,undoc,union,unique,uniroot,unix,%
        unlink,unlist,unname,untrace,update,upper,url,UseMethod,var,%
        variable,vector,Version,vi,warning,warnings,weighted,weights,%
        which,while,window,write,\%x\%,x11,X11,xedit,xemacs,xinch,xor,%
        xpdrows,xy,xyinch,yinch,zapsmall,zip},%
    otherkeywords={!,!=,~,$,*,\%,\&,\%/\%,\%*\%,\%\%,<-,<<-},%$
    alsoother={._$},%$
    sensitive,%
    morecomment=[l]\#,%
    morestring=[d]",%
    morestring=[d]'% 2001 Robert Denham
}%

%решаем проблему с кириллицей в комментариях (в pdflatex) https://tex.stackexchange.com/a/103712
\lstset{extendedchars=true,keepspaces=true,literate={Ö}{{\"O}}1
    {Ä}{{\"A}}1
    {Ü}{{\"U}}1
    {ß}{{\ss}}1
    {ü}{{\"u}}1
    {ä}{{\"a}}1
    {ö}{{\"o}}1
    {~}{{\textasciitilde}}1
    {а}{{\selectfont\char224}}1
    {б}{{\selectfont\char225}}1
    {в}{{\selectfont\char226}}1
    {г}{{\selectfont\char227}}1
    {д}{{\selectfont\char228}}1
    {е}{{\selectfont\char229}}1
    {ё}{{\"e}}1
    {ж}{{\selectfont\char230}}1
    {з}{{\selectfont\char231}}1
    {и}{{\selectfont\char232}}1
    {й}{{\selectfont\char233}}1
    {к}{{\selectfont\char234}}1
    {л}{{\selectfont\char235}}1
    {м}{{\selectfont\char236}}1
    {н}{{\selectfont\char237}}1
    {о}{{\selectfont\char238}}1
    {п}{{\selectfont\char239}}1
    {р}{{\selectfont\char240}}1
    {с}{{\selectfont\char241}}1
    {т}{{\selectfont\char242}}1
    {у}{{\selectfont\char243}}1
    {ф}{{\selectfont\char244}}1
    {х}{{\selectfont\char245}}1
    {ц}{{\selectfont\char246}}1
    {ч}{{\selectfont\char247}}1
    {ш}{{\selectfont\char248}}1
    {щ}{{\selectfont\char249}}1
    {ъ}{{\selectfont\char250}}1
    {ы}{{\selectfont\char251}}1
    {ь}{{\selectfont\char252}}1
    {э}{{\selectfont\char253}}1
    {ю}{{\selectfont\char254}}1
    {я}{{\selectfont\char255}}1
    {А}{{\selectfont\char192}}1
    {Б}{{\selectfont\char193}}1
    {В}{{\selectfont\char194}}1
    {Г}{{\selectfont\char195}}1
    {Д}{{\selectfont\char196}}1
    {Е}{{\selectfont\char197}}1
    {Ё}{{\"E}}1
    {Ж}{{\selectfont\char198}}1
    {З}{{\selectfont\char199}}1
    {И}{{\selectfont\char200}}1
    {Й}{{\selectfont\char201}}1
    {К}{{\selectfont\char202}}1
    {Л}{{\selectfont\char203}}1
    {М}{{\selectfont\char204}}1
    {Н}{{\selectfont\char205}}1
    {О}{{\selectfont\char206}}1
    {П}{{\selectfont\char207}}1
    {Р}{{\selectfont\char208}}1
    {С}{{\selectfont\char209}}1
    {Т}{{\selectfont\char210}}1
    {У}{{\selectfont\char211}}1
    {Ф}{{\selectfont\char212}}1
    {Х}{{\selectfont\char213}}1
    {Ц}{{\selectfont\char214}}1
    {Ч}{{\selectfont\char215}}1
    {Ш}{{\selectfont\char216}}1
    {Щ}{{\selectfont\char217}}1
    {Ъ}{{\selectfont\char218}}1
    {Ы}{{\selectfont\char219}}1
    {Ь}{{\selectfont\char220}}1
    {Э}{{\selectfont\char221}}1
    {Ю}{{\selectfont\char222}}1
    {Я}{{\selectfont\char223}}1
    {і}{{\selectfont\char105}}1
    {ї}{{\selectfont\char168}}1
    {є}{{\selectfont\char185}}1
    {ґ}{{\selectfont\char160}}1
    {І}{{\selectfont\char73}}1
    {Ї}{{\selectfont\char136}}1
    {Є}{{\selectfont\char153}}1
    {Ґ}{{\selectfont\char128}}1
}

% Ширина текста минус ширина надписи 999
\newlength{\twless}
\newlength{\lmarg}
\setlength{\lmarg}{\widthof{999}}   % ширина надписи 999
\setlength{\twless}{\textwidth-\lmarg}

\lstset{ %
%    language=R,                     %  Язык указать здесь, если во всех листингах преимущественно один язык, в результате часть настроек может пойти только для этого языка
    numbers=left,                   % where to put the line-numbers
    numberstyle=\fontsize{12pt}{14pt}\selectfont\color{Gray},  % the style that is used for the line-numbers
    firstnumber=1,                  % в этой и следующей строках задаётся поведение нумерации 5, 10, 15...
    stepnumber=5,                   % the step between two line-numbers. If it's 1, each line will be numbered
    numbersep=5pt,                  % how far the line-numbers are from the code
    backgroundcolor=\color{white},  % choose the background color. You must add \usepackage{color}
    showspaces=false,               % show spaces adding particular underscores
    showstringspaces=false,         % underline spaces within strings
    showtabs=false,                 % show tabs within strings adding particular underscores
    frame=leftline,                 % adds a frame of different types around the code
    rulecolor=\color{black},        % if not set, the frame-color may be changed on line-breaks within not-black text (e.g. commens (green here))
    tabsize=2,                      % sets default tabsize to 2 spaces
    captionpos=t,                   % sets the caption-position to top
    breaklines=true,                % sets automatic line breaking
    breakatwhitespace=false,        % sets if automatic breaks should only happen at whitespace
%    title=\lstname,                 % show the filename of files included with \lstinputlisting;
    % also try caption instead of title
    basicstyle=\fontsize{12pt}{14pt}\selectfont\ttfamily,% the size of the fonts that are used for the code
%    keywordstyle=\color{blue},      % keyword style
    commentstyle=\color{ForestGreen}\emph,% comment style
    stringstyle=\color{Mahogany},   % string literal style
    escapeinside={\%*}{*)},         % if you want to add a comment within your code
    morekeywords={*,...},           % if you want to add more keywords to the set
    inputencoding=utf8,             % кодировка кода
    xleftmargin={\lmarg},           % Чтобы весь код и полоска с номерами строк была смещена влево, так чтобы цифры не вылезали за пределы текста слева
}

%http://tex.stackexchange.com/questions/26872/smaller-frame-with-listings
% Окружение, чтобы листинг был компактнее обведен рамкой, если она задается, а не на всю ширину текста
\makeatletter
\newenvironment{SmallListing}[1][]
{\lstset{#1}\VerbatimEnvironment\begin{VerbatimOut}{VerbEnv.tmp}}
{\end{VerbatimOut}\settowidth\@tempdima{%
        \lstinputlisting{VerbEnv.tmp}}
    \minipage{\@tempdima}\lstinputlisting{VerbEnv.tmp}\endminipage}
\makeatother

\DefineVerbatimEnvironment% с шрифтом 12 пт
{Verb}{Verbatim}
{fontsize=\fontsize{12pt}{14pt}\selectfont}

\newfloat[chapter]{ListingEnv}{lol}{Listing}

\renewcommand{\lstlistingname}{Listing}

%Общие счётчики окружений листингов
%http://tex.stackexchange.com/questions/145546/how-to-make-figure-and-listing-share-their-counter
% Если смешивать плавающие и не плавающие окружения, то могут быть проблемы с нумерацией
\makeatletter
\AfterEndPreamble{% https://tex.stackexchange.com/a/252682
    \let\c@ListingEnv\relax % drop existing counter "ListingEnv"
    \newaliascnt{ListingEnv}{lstlisting} % команда требует пакет aliascnt
    \let\ftype@lstlisting\ftype@ListingEnv % give the floats the same precedence
}
\makeatother

% значок С++ — используйте команду \cpp
\newcommand{\cpp}{%
    C\nolinebreak\hspace{-.05em}%
    \raisebox{.2ex}{+}\nolinebreak\hspace{-.10em}%
    \raisebox{.2ex}{+}%
}

%%%  Чересстрочное форматирование таблиц
%% http://tex.stackexchange.com/questions/278362/apply-italic-formatting-to-every-other-row
\newcounter{rowcnt}
\newcommand\altshape{\ifnumodd{\value{rowcnt}}{\color{red}}{\vspace*{-1ex}\itshape}}
% \AtBeginEnvironment{tabular}{\setcounter{rowcnt}{1}}
% \AtEndEnvironment{tabular}{\setcounter{rowcnt}{0}}

%%% Ради примера во второй главе
\let\originalepsilon\epsilon
\let\originalphi\phi
\let\originalkappa\kappa
\let\originalle\le
\let\originalleq\leq
\let\originalge\ge
\let\originalgeq\geq
\let\originalemptyset\emptyset
\let\originaltan\tan
\let\originalcot\cot
\let\originalcsc\csc

%%% Русская традиция начертания математических знаков
\renewcommand{\le}{\ensuremath{\leqslant}}
\renewcommand{\leq}{\ensuremath{\leqslant}}
\renewcommand{\ge}{\ensuremath{\geqslant}}
\renewcommand{\geq}{\ensuremath{\geqslant}}
\renewcommand{\emptyset}{\varnothing}

%%% Русская традиция начертания математических функций (на случай копирования из зарубежных источников)
\renewcommand{\tan}{\operatorname{tg}}
\renewcommand{\cot}{\operatorname{ctg}}
\renewcommand{\csc}{\operatorname{cosec}}

%%% Русская традиция начертания греческих букв (греческие буквы вертикальные, через пакет upgreek)
\renewcommand{\epsilon}{\ensuremath{\upvarepsilon}}   %  русская традиция записи
\renewcommand{\phi}{\ensuremath{\upvarphi}}
%\renewcommand{\kappa}{\ensuremath{\varkappa}}
\renewcommand{\alpha}{\upalpha}
\renewcommand{\beta}{\upbeta}
\renewcommand{\gamma}{\upgamma}
\renewcommand{\delta}{\updelta}
\renewcommand{\varepsilon}{\upvarepsilon}
\renewcommand{\zeta}{\upzeta}
\renewcommand{\eta}{\upeta}
\renewcommand{\theta}{\uptheta}
\renewcommand{\vartheta}{\upvartheta}
\renewcommand{\iota}{\upiota}
\renewcommand{\kappa}{\upkappa}
\renewcommand{\lambda}{\uplambda}
\renewcommand{\mu}{\upmu}
\renewcommand{\nu}{\upnu}
\renewcommand{\xi}{\upxi}
\renewcommand{\pi}{\uppi}
\renewcommand{\varpi}{\upvarpi}
\renewcommand{\rho}{\uprho}
%\renewcommand{\varrho}{\upvarrho}
\renewcommand{\sigma}{\upsigma}
%\renewcommand{\varsigma}{\upvarsigma}
\renewcommand{\tau}{\uptau}
\renewcommand{\upsilon}{\upupsilon}
\renewcommand{\varphi}{\upvarphi}
\renewcommand{\chi}{\upchi}
\renewcommand{\psi}{\uppsi}
\renewcommand{\omega}{\upomega}
 % Стили для специфических пользовательских задач

%%% Библиография. Выбор движка для реализации %%%
% Здесь только проверка установленного ключа. Сама настройка выбора движка
% размещена в common/setup.tex
\ifnumequal{\value{bibliosel}}{0}{%
    \input{biblio/predefined}   % Встроенная реализация с загрузкой файла через движок bibtex8
}{
    \input{biblio/biblatex}     % Реализация пакетом biblatex через движок biber
}

% Вывести информацию о выбранных опциях в лог сборки
\typeout{Selected options:}
\typeout{Draft mode: \arabic{draft}}
\typeout{Font: \arabic{fontfamily}}
\typeout{AltFont: \arabic{usealtfont}}
\typeout{Bibliography backend: \arabic{bibliosel}}
\typeout{Precompile images: \arabic{imgprecompile}}
% Вывести информацию о версиях используемых библиотек в лог сборки
\listfiles

%%% Управление компиляцией отдельных частей диссертации %%%
% Необходимо сначала иметь полностью скомпилированный документ, чтобы все
% промежуточные файлы были в наличии
% Затем, для вывода отдельных частей можно воспользоваться командой \includeonly
% Ниже примеры использования команды:
%
%\includeonly{Dissertation/part2}
%\includeonly{Dissertation/contents,Dissertation/appendix,Dissertation/conclusion}
%
% Если все команды закомментированы, то документ будет выведен в PDF файл полностью

\begin{document}
%%% Переопределение именований типовых разделов
% https://tex.stackexchange.com/a/156050
\gappto\captionsrussian{%%% Переопределение именований %%%
\renewcommand{\contentsname}{Contents}% (ГОСТ Р 7.0.11-2011, 4)
\renewcommand{\figurename}{Figure}% (ГОСТ Р 7.0.11-2011, 5.3.9)
\renewcommand{\tablename}{Table}% (ГОСТ Р 7.0.11-2011, 5.3.10)
\renewcommand{\listfigurename}{List of Figures}%
\renewcommand{\listtablename}{List of Tables}%
\renewcommand{\bibname}{\bibtitlefull}%
% Переопределения названий для nomencl. Так как опция russian не для utf8
\renewcommand{\nomname}{Список сокращений и условных обозначений}%
\renewcommand{\eqdeclaration}[1]{, см.~(#1)}%
\renewcommand{\pagedeclaration}[1]{, стр.~#1}%
\renewcommand{\nomAname}{Латинские буквы}%
\renewcommand{\nomGname}{Греческие буквы}%
\renewcommand{\nomXname}{Верхние индексы}%
\renewcommand{\nomZname}{Индексы}%\unskip} % for polyglossia and babel
%%% Переопределение именований %%%
\renewcommand{\contentsname}{Contents}% (ГОСТ Р 7.0.11-2011, 4)
\renewcommand{\figurename}{Figure}% (ГОСТ Р 7.0.11-2011, 5.3.9)
\renewcommand{\tablename}{Table}% (ГОСТ Р 7.0.11-2011, 5.3.10)
\renewcommand{\listfigurename}{List of Figures}%
\renewcommand{\listtablename}{List of Tables}%
\renewcommand{\bibname}{\bibtitlefull}%
% Переопределения названий для nomencl. Так как опция russian не для utf8
\renewcommand{\nomname}{Список сокращений и условных обозначений}%
\renewcommand{\eqdeclaration}[1]{, см.~(#1)}%
\renewcommand{\pagedeclaration}[1]{, стр.~#1}%
\renewcommand{\nomAname}{Латинские буквы}%
\renewcommand{\nomGname}{Греческие буквы}%
\renewcommand{\nomXname}{Верхние индексы}%
\renewcommand{\nomZname}{Индексы}%

%%% Структура диссертации (ГОСТ Р 7.0.11-2011, 4)
\thispagestyle{empty}

\noindent%
\begin{tabularx}{\textwidth}{@{}lXr@{}}%
    & & \large{На правах рукописи}\\
    \IfFileExists{images/logo.pdf}{\includegraphics[height=2.5cm]{logo}}{\rule[0pt]{0pt}{2.5cm}}  & &
    \ifnumequal{\value{showperssign}}{0}{%
        \rule[0pt]{0pt}{1.5cm}
    }{
        \includegraphics[height=1.5cm]{personal-signature.png}
    }\\
\end{tabularx}

\vspace{0pt plus1fill} %число перед fill = кратность относительно некоторого расстояния fill, кусками которого заполнены пустые места
\begin{center}
\textbf {\large \thesisAuthor}
\end{center}

\vspace{0pt plus3fill} %число перед fill = кратность относительно некоторого расстояния fill, кусками которого заполнены пустые места
\begin{center}
\textbf {\Large %\MakeUppercase
\thesisTitle}

\vspace{0pt plus3fill} %число перед fill = кратность относительно некоторого расстояния fill, кусками которого заполнены пустые места
{\large Специальность \thesisSpecialtyNumber\ "---\par <<\thesisSpecialtyTitle>>}

\ifdefined\thesisSpecialtyTwoNumber
{\large Специальность \thesisSpecialtyTwoNumber\ "---\par <<\thesisSpecialtyTwoTitle>>}
\fi

\vspace{0pt plus1.5fill} %число перед fill = кратность относительно некоторого расстояния fill, кусками которого заполнены пустые места
\Large{Автореферат}\par
\large{диссертации на соискание учёной степени\par \thesisDegree}
\end{center}

\vspace{0pt plus4fill} %число перед fill = кратность относительно некоторого расстояния fill, кусками которого заполнены пустые места
{\centering\thesisCity~--- \thesisYear\par}

\newpage
% оборотная сторона обложки
\thispagestyle{empty}
\noindent Работа выполнена на {\thesisInOrganization}.

\vspace{0.008\paperheight plus1fill}
\noindent%
\begin{tabularx}{\textwidth}{@{}lX@{}}
    \ifdefined\supervisorTwoFio
    Научные руководители:   & \supervisorRegalia\par
                              \ifdefined\supervisorDead
                              \framebox{\textbf{\supervisorFio}}
                              \else
                              \textbf{\supervisorFio}
                              \fi
                              \par
                              \vspace{0.013\paperheight}
                              \supervisorRegalia\par
                              \ifdefined\supervisorTwoDead
                              \framebox{\textbf{\supervisorTwoFio}}
                              \else
                              \textbf{\supervisorTwoFio}
                              \fi
                              \vspace{0.013\paperheight}\\
    \else
    Научный руководитель:   & \supervisorRegalia\par
                              \ifdefined\supervisorDead
                              \framebox{\textbf{\supervisorFio}}
                              \else
                              \textbf{\supervisorFio}
                              \fi
                              \vspace{0.013\paperheight}\\
    \fi
    Официальные оппоненты:  &
    \ifnumequal{\value{showopplead}}{0}{\vspace{13\onelineskip plus1fill}}{%
        \textbf{\opponentOneFio,}\par
        \opponentOneRegalia,\par
        \opponentOneJobPlace,\par
        \opponentOneJobPost\par
        \vspace{0.01\paperheight}
        \textbf{\opponentTwoFio,}\par
        \opponentTwoRegalia,\par
        \opponentTwoJobPlace,\par
        \opponentTwoJobPost
    \ifdefined\opponentThreeFio
        \par
        \vspace{0.01\paperheight}
        \textbf{\opponentThreeFio,}\par
        \opponentThreeRegalia,\par
        \opponentThreeJobPlace,\par
        \opponentThreeJobPost
    \fi
    }%
    \vspace{0.013\paperheight} \\
    \ifdefined\leadingOrganizationTitle
    Ведущая организация:    &
    \ifnumequal{\value{showopplead}}{0}{\vspace{6\onelineskip plus1fill}}{%
        \leadingOrganizationTitle
    }%
    \fi
\end{tabularx}
\vspace{0.008\paperheight plus1fill}

\noindent Защита состоится \defenseDate~на~заседании диссертационного совета \defenseCouncilNumber~при \defenseCouncilTitle~по адресу: \defenseCouncilAddress.

\vspace{0.008\paperheight plus1fill}
\noindent С диссертацией можно ознакомиться в библиотеке \synopsisLibrary.

\vspace{0.008\paperheight plus1fill}
\noindent Отзывы на автореферат в двух экземплярах, заверенные печатью учреждения, просьба направлять по адресу: \defenseCouncilAddress, ученому секретарю диссертационного совета~\defenseCouncilNumber.

\vspace{0.008\paperheight plus1fill}
\noindent{Автореферат разослан \synopsisDate.}

\noindent Телефон для справок: \defenseCouncilPhone.

\vspace{0.008\paperheight plus1fill}
\noindent%
\begin{tabularx}{\textwidth}{@{}%
>{\raggedright\arraybackslash}b{18em}@{}
>{\centering\arraybackslash}X
r
@{}}
    Ученый секретарь\par
    диссертационного совета\par
    \defenseCouncilNumber,\par
    \defenseSecretaryRegalia
    &
    \ifnumequal{\value{showsecrsign}}{0}{}{%
        \includegraphics[width=2cm]{secretary-signature.png}%
    }%
    &
    \defenseSecretaryFio
\end{tabularx}
           % Титульный лист
\include{Dissertation/contents}        % Оглавление
\ifnumequal{\value{contnumfig}}{1}{}{\counterwithout{figure}{chapter}}
\ifnumequal{\value{contnumtab}}{1}{}{\counterwithout{table}{chapter}}
\chapter*{Introduction}                         % Заголовок
\addcontentsline{toc}{chapter}{Introduction}    % Добавляем его в оглавление

\newcommand{\actuality}{}
\newcommand{\progress}{}
\newcommand{\aim}{{\textbf\aimTXT}}
\newcommand{\tasks}{\textbf{\tasksTXT}}
\newcommand{\novelty}{\textbf{\noveltyTXT}}
\newcommand{\influence}{\textbf{\influenceTXT}}
\newcommand{\methods}{\textbf{\methodsTXT}}
\newcommand{\defpositions}{\textbf{\defpositionsTXT}}
\newcommand{\reliability}{\textbf{\reliabilityTXT}}
\newcommand{\probation}{\textbf{\probationTXT}}
\newcommand{\contribution}{\textbf{\contributionTXT}}
\newcommand{\publications}{\textbf{\publicationsTXT}}


{\actuality} In game theory, modeling of conflicts with three or more parties independently pursuing own goals is, for many reasons, considered much more difficult comparing to modeling of classical bilateral confrontations. Among these reasons one plays very special part "--- the influence of information asymmetry exerted on course of a struggle. For games with two participants its effect generally comes down to the consequences of a priori incompleteness in their knowledge about the parameters of the conflict and is usually described with Bayesian models. Common knowledge in such cases leaves out players' payoff functions, so each of them acts according to their own, possibly different assumptions, expressed in the form of a probability distribution in the space of all payoff functions with a given set of strategies. %В теории игр конфликты с участием трёх и более преследующих собственные цели сторон по многим причинам считаются существенно более трудными для моделирования в сравнении с классическими парными противостояниями. Среди этих причин следует особо подчеркнуть влияние, оказываемое на их ход информационной асимметрией. В играх с двумя участниками её эффект, в целом, сводится к последствиям априорной неполноты их знаний о параметрах конфликта и обычно описывается при помощи байесовских моделей. Платёжная функция в таких случаях не является общим знанием игроков, каждый из них действует исходя из собственных, возможно, различающихся предположений, выраженных в форме распределения вероятностей на пространстве всевозможных платёжных функций с заданным множеством стратегий.

For multistage bilateral conflicts there can be another source of information asymmetry "--- the presence of a secret component in the actions of opponents. If this occurs, actions already performed by the player in the early stages may be completely or partially unknown to his opponent, who is forced to base the strategy of later stages on assumptions. Naturally, this is formalized through partitioning of the extensive-form game tree into information sets. Effectively, this two aspects exhaust the impact of information asymmetry on two-sided conflicts. However, addition of third participant induces a new phenomenon, noticed by Robert Aumann in his article \cite{Aumann74} introducing now-familiar new formalism "--- correlated extension of normal form games. %Ещё одним источником информационной асимметрии в парном конфликте может выступать присутствие в действиях оппонентов тайной составляющей в тех случаях, когда он проходит в несколько стадий. При этом действия, уже совершённые игроком на ранних стадиях, могут быть полностью или частично неизвестны его противнику, вынужденному основывать стратегию более поздних стадий на предположениях. Это естественным образом формализуется при помощи информационных разбиений дерева ходов в развёрнутой форме игры. По сути, двумя упомянутыми аспектами исчерпывается влияние информационной асимметрии на конфликты с двумя сторонами. Однако, увеличение количества участников ещё хотя бы на одного порождает новый феномен, замеченный ещё Робертом Ауманом в статье \ifsynopsis\smartcite{Aumann74}\else\cite{Aumann74}\fi, где впервые было сформулировано коррелированное расширение игр в нормальной форме.

This phenomenon reflects on comparison of the sets of Nash equilibria for the same games, calculated using mixed strategies on one level and involving external correlation mechanisms on the other. In case of two players, all vectors of expected utility in the correlated equilibria belong to the convex hull of the set of mixed equilibria utilities, i.e., correlation mechanisms can be thought of simply as a way to obtain linear combinations of classical solutions. With the advent of the third player, the picture changes "--- in some games, the presence of a non-public correlation mechanism allows reaching Nash equilibrium points with payoffs outside the convex hull of mixed solutions' utilities. Effectively, it means that asymmetry of knowledge can have a significant impact on the outcome of a multilateral conflict even in cases when its subject have no relation to the payout structure. Let's call the games prone to this effect \emph{sensitive to additional information asymmetry}. %Выражается этот феномен в том, как для одних и тех же игр соотносятся множества равновесий по Нэшу в смешанных стратегиях и с использованием внешних механизмов корреляции. Для случая двух игроков все вектора математических ожиданий выплат в точках коррелированных равновесий принадлежат выпуклой оболочке множества смешанных равновесий, т.\:е. механизмы корреляции можно рассматривать просто как способ получения линейных комбинаций классических решений. С появлением третьего игрока картина меняется "--- в некоторых играх присутствие непубличного корреляционного механизма позволяет достичь точек равновесия по Нэшу с выплатами за пределами выпуклой оболочки решений в смешанных стратегиях. Фактически, это означает, что асимметрия знаний может оказывать существенное влияние на исход многостороннего конфликта даже в тех случаях, когда она не касается существенных для структуры выплат факторов. Назовём игры, подверженные этому эффекту, \emph{чувствительными к дополнительной информационной асимметрии}.

While the described phenomenon has been known for a while, most of the models built by researchers bypassed it, as it is considered rather feature of curiosity in some games of many players. An important exception worth noting is, perhaps, the most significant area of game theory actively using the correlated extension formalism of games in normal form "--- <<mechanism design>> \cite{Nikolenko} by Leonid Gurvitch, Eric Maskin, and Roger Myerson. Their approach aims to create economic instruments that incentivize selfish rational agents to behave optimally in alignment with common goal functions that formalize various social goods. Being an extremely fruitful area of research, mechanism design has spawned many directions and branches, united nevertheless by a number of integral common features arising from the information structure of games for which its main provisions are proved. A typical interaction scheme looks like this: players"~agents, knowing their own preferences and capabilities, but being ignorant of these parameters for other participants, report their type to the central authority, which coins a set of correlated strategies based on this information. Next, the center actualizes it for an instance in the form of a pure strategies set and instructs each of the agents, who, in turn, make the final decision on a particular action. This implies that agents can lie at the first stage and disobey at the last. In paradigm of mechanism design the main goal is about the creation of such algorithms for the behavior of central authority that the strategies of truthfulness and obedience form a Nash equilibrium for agents. %Хотя описанный феномен известен уже давно, при построении моделей большинство исследователей обходили его стороной, рассматривая скорее как курьёзную особенность некоторых игр многих игроков. Важным исключением при этом выступает, пожалуй, наиболее значимая из активно использующих формализм коррелированного расширения игр в нормальной форме область теории игр "--- <<дизайн механизмов>> \ifsynopsis\smartcite{Nikolenko}\else\cite{Nikolenko}\fi Леонида Гурвича, Эрика Маскина и Роджера Майерсона. Их подход ставит своей целью создание экономических инструментов, стимулирующих эгоистичных рациональных агентов к поведению, оптимальному с точки зрения общих целевых функций, формализующих различные социальные блага. Будучи чрезвычайно плодотворной областью исследований, дизайн механизмов породил множество направлений и ответвлений, объединённых тем не менее рядом неотъемлемых общих черт, проистекающих из информационной структуры игр, для которых доказываются его основные положения. Типичная схема взаимодействий выглядит так: игроки"~агенты, знающие свои предпочтения и возможности, но находящиеся в неведении относительно этих параметров у других участников, информируют о них центр, формирующий на основе этой информации набор коррелированных стратегий. Далее центр реализует его для конкретного случая в виде набора чистых стратегий и инструктирует каждого из агентов, которые в свою очередь и принимают окончательное решение о том или ином действии. При этом подразумевается, что агенты могут лгать на первом этапе и не подчиняться на последнем. Главной задачей дизайна механизмов в этой парадигме становится создание таких алгоритмов поведения центра, что стратегии правдивости и послушания образуют для агентов равновесие Нэша.

Without diving into details, it can be said that mechanism design is based on a special case of information asymmetry "--- a kind of star connection structure, where a dedicated central agent controls the single correlation mechanism in his own interests, while subordinate agents are in complete isolation both from each other, and from the outside world. The name here accurately reflects the view of conflicts inherent for the model "--- through its lens, the players are seen as interchangeable parts of a single man-made mechanism, not connected by anything other than participation in it. Although modeling within the framework of this simplification can be useful in construction of formalized methods for conflict resolution, it doesn't cover cases where they are significantly affected by information asymmetries developing not as a result of calculated design, but naturally, as spontaneous interaction occurs between agents in an uncontrolled, external to the model environment. For example, it is difficult to overestimate the significance of corruption's influence on political and economic institutions, and yet it is made up of precisely such unplanned information links external to the institutions themselves. %Не вдаваясь в детали, можно сказать, что дизайн механизмов базируется на частном случае информационной асимметрии "--- своего рода звёздчатой структуре связей, где выделенный центральный агент может в своих интересах распоряжаться общим механизмом корреляции, а подчинённые ему агенты находятся в полной изоляции как друг от друга, так и от остального мира. Название здесь действительно неплохо отражает свойственный для модели взгляд на конфликты "--- через её призму стороны рассматриваются как взаимозаменимые детали единого рукотворного механизма, не связанные ничем кроме участия в нём. Хотя моделирование в рамках подобного упрощения вполне может быть полезно при конструировании формализованных способов решения конфликтов, оно никак не может помочь в тех случаях, когда на них оказывает существенное влияние информационная асимметрия, складывающаяся не в результате сознательного дизайна, а естественным образом, по мере спонтанного взаимодействия агентов в неконтролируемой, внешней с точки зрения модели среде. К примеру, сложно переоценить значимость влияния коррупции на политические и экономические институты, а ведь она складывается именно из таких незапланированных информационных связей, внешних по отношению к самим институтам.

For the reason described above, the study of the influence of additional information asymmetry on the solutions of multilateral games shouldn't be reduced to purely constructive models. Alas, outside of mechanism design disregard for this phenomenon became prevailing convention. For example, in the article \cite{Fudenberg} by Drew Fudenberg and Eric Maskin, you can find the following footnote: <<In essence, if $n \ge 3$, the rest of the players get the opportunity to omit player $j$ payouts even lower, using a correlated strategy against him, in which player $j$ cannot observe the signal of the correlation mechanism (...). Adhering to the tradition that has developed in publications devoted to repetitive games, however, we will not consider such correlated strategies.>> It would seem, however, that in the context of repeated games with discounting, the topic of using the secrecy of the correlation mechanism to strengthen punishment strategies is rather intriguing "--- examples of agent groups strengthening their collective long-term position through necessarily secret coordination of actions seemingly isn't hard to find in every field of research that uses the folk theorem, from anthropology to international politics. Disconcertingly, up to date game theory has very little to offer as an analysis tool for the described phenomenon to other sciences. %По описанной причине исследование влияния дополнительной информационной асимметрии на решения игр многих игроков никак нельзя сводить только к конструктивным моделям. Увы, но за пределами дизайна механизмов сложилась традиция игнорировать этот феномен. К примеру, в статье \cite{Fudenberg} Дрю Фуденберга и Эрика Маскина можно найти следующую сноску: <<Actually, if $n \ge 3$, the other players may be able to keep player $j$'s payoff even lower by using a correlated strategy against $j$, where the outcome of the correlating device is not observed by $j$ (...). In keeping with the rest of the literature on repeated games, however, we shall rule out such correlated strategies.>>\footnote{<<В сущности, если $n \ge 3$, остальные игроки получают возможность опустить выплаты игрока $j$ ещё ниже, используя против него коррелированную стратегию, в которой игрок $j$ не может наблюдать сигнала механизма корреляции (...). Придерживаясь сложившейся в посвящённых повторяющимся играм публикациях традиции, мы, впрочем, не будем рассматривать такие коррелированные стратегии.>>} А ведь казалось бы, в контексте повторяющихся игр с дисконтированием проблематика использования секретности механизма корреляции для усиления стратегий наказания довольно любопытна "--- наверное в каждой области исследований, использующей народную теорему, от антропологии до международной политики, несложно отыскать примеры того, как группы агентов усиливали свою коллективную долгосрочную позицию при помощи необходимо тайного согласования действий. Увы, но приходится констатировать, что теории игр до сих пор почти нечего предложить другим наукам в качестве инструмента анализа описанного феномена.

% {\progress}
% Этот раздел должен быть отдельным структурным элементом по
% ГОСТ, но он, как правило, включается в описание актуальности
% темы. Нужен он отдельным структурынм элемементом или нет ---
% смотрите другие диссертации вашего совета, скорее всего не нужен.

This study is \textbf{aim}ed toward development of new multilateral conflict model, that takes into account how its course is affected by additional informational asymmetry. %{\aim} данной работы является создание новой модели многосторонних конфликтов, учитывающей влияние на их ход дополнительной информационной асимметрии.

To approach stated goal following \textbf{tasks} had to be addressed: %Для~достижения поставленной цели необходимо было решить следующие {\tasks}:
\begin{enumerate}[beginpenalty=10000] % https://tex.stackexchange.com/a/476052/104425
  \item Analyze the correlated extension of normal form games model in scope of the research. %Исследовать формализм коррелированного обобщения игр в нормальной форме с точки зрения проблематики работы.
  \item Develop the notation aimed to describe the informational structures, that could be tying together participants of the arbitrary conflicts in a variety of fashions. %Разработать способ описания информационных структур, достаточно разнообразным образом связывающих участников произвольного конфликта.
  \item Analyze the impact of additional informational asymmetry on the solutions' conformity with the criteria of collective rationality. %Исследовать влияние дополнительной информационной асимметрии на соответствие равновесий критериям коллективной рациональности.
  \item Develop the reasonable solution concept for the games with additional informational asymmetry, taking into account the inter-agent relations. %Разработать приемлемую концепцию решения с учётом связей между агентами для игр с дополнительной информационной асимметрией.
\end{enumerate}


\textbf{Scientific novelty:}
\begin{enumerate}[beginpenalty=10000] % https://tex.stackexchange.com/a/476052/104425
  \item The games of many players, which are sensitive to additional information asymmetry, are singled out as an independent object of study. %Выделены в качестве самостоятельного объекта исследования игры многих игроков, проявляющие чувствительность к дополнительной информационной асимметрии.
  \item A formalism of the conspiracy space is proposed, purposefully narrowing the formalism of the correlation space in order to model additional information asymmetry. %Предложен формализм пространства заговоров, специальным образом сужающий в целях моделирования дополнительной информационной асимметрии формализм пространства корреляции.
  \item The concept of structurally consistent equilibrium is formulated, in many cases allowing to single out among the solutions of games in conspiracy spaces those that adhere to the principle of collective rationality. %Сформулирована концепция структурно согласованного равновесия, позволяющая во многих случаях выделять среди решений игр в пространствах заговоров отвечающие принципу коллективной рациональности.
  \item The possibility of extending the set of perfect subgame equilibria in repeated games sensitive to additional information asymmetry  with the help of modern cryptography tools is demonstrated. %Показана возможность расширения множества совершенных подыгровых равновесий в повторяющихся играх, чувствительных к дополнительной информационной асимметрии, при помощи инструментов современной криптографии.
\end{enumerate}

The \textbf{practical significance} of the work stems from the obvious need to take into account, when modeling multilateral conflicts, the fact that the composition of their participants is not, in most cases, a random sample of agents that are not related to each other in any way. The classical formalism of games in normal form relies on the implicit assumption that the only significant characteristic of each player is the order of preference regarding the outcome of the draw, expressed in the form of a payoff function. At the same time, it is quite obvious that real people entering into a confrontation are often connected by relationships, significant for its outcome, the structure of which cannot be expressed by a simple combination of payment functions. Such connection can be clearly illustrated by comparing two imaginary games of bridge or preference with equally strong players, differing in that one table is occupied by strangers, while at the other some have been playing together for many years. Any sufficiently experienced card-player can confirm that, given equal skill, the <<cohesion>> factor between partners reliably provides a decisive advantage. Naturally, this phenomenon can be generalized to more significant conflicts: politics, business, diplomacy "--- wherever the outcome of the confrontation depends significantly on the coordination and unpredictability of actions, mutual understanding that does not require communication can often turn defeat into victory. Thus, to improve accuracy of predictions for the outcomes of multilateral conflicts, there is a compelling need for models that take this factor into account. %{\influence} работы проистекает из явной необходимости учитывать при моделировании многосторонних конфликтов тот факт, что состав их участников не является в большинстве случаев случайной выборкой никак не связанных друг с другом агентов. Классический формализм игр в нормальной форме опирается на неявное допущение, состоящее в том, что единственной значимой характеристикой каждого игрока является порядок его предпочтений относительно исхода розыгрыша, выражающийся в форме платёжной функции. Совершенно очевидно при этом, что реальных людей, вступающих в противостояние, зачастую связывают значимые для его исхода отношения, структура которых не может быть выражена простым сочетанием платёжных функций. \ifsynopsis\else В качестве наглядной иллюстрации такой связи можно сравнить две воображаемые партии в бридж или преферанс с участием одинаково сильных игроков, различающиеся тем, что в одном случае за столом сидят незнакомцы, а в другом часть из них играют вместе уже много лет. Любой достаточно опытный картёжник скажет, что при равных навыках фактор <<сыгранности>> с партнёром надёжно обеспечивает решающее преимущество. Естественным образом этот феномен можно обобщить и на более значимые конфликты: политика, бизнес, дипломатия "--- везде, где исход противостояния существенно зависит от согласованности и непредсказуемости действий, взаимопонимание, не требующее коммуникации, зачастую может превратить поражение в победу. \fi Таким образом, для более точного предсказания исходов многосторонних конфликтов насущно необходимы модели, позволяющие учитывать этот фактор.

\textbf{Methodology and research methods.} The study utilizes frameworks of game theory, probability theory, topology and cryptography. %В работе используются методы теории игр, теории вероятности, топологии и криптографии.

\textbf{Defense positions:}
\begin{enumerate}[beginpenalty=10000] % https://tex.stackexchange.com/a/476052/104425
  \item Proof of the theorem on the isomorphism of correlation spaces, which defines equivalence classes for them, demarcating indistinguishable from the game-theoretic point of view spaces. %Доказательство теоремы об изоморфизме пространств корреляции, определяющей для них классы эквивалентности, в которые входят только неразличимые с теоретико-игровой точки зрения пространства.
  \item A proof of the theorem on conspiracy spaces of one structure, thanks to which the structure of a space describes it comprehensively. %Доказательство теоремы о пространствах заговоров одной структуры, благодаря которой структуру пространства можно считать его исчерпывающим описанием.
  \item Proof of sensitivity to additional information asymmetry of the symmetric task scheduling problem with nonmonotonic returns. %Доказательство чувствительности к дополнительной информационной асимметрии симметричной проблемы планирования заданий с немонотонной отдачей.
  \item A solution to a three-way symmetric scheduling problem with a nonmonotonic rush premium function in an asymmetric conspiracy space that satisfies the structural consistency criterion. %Решение трёхсторонней симметричной проблемы планирования с немонотонной функцией оплаты за срочность в ассиметричном пространстве заговоров, удовлетворяющее критерию структурной согласованности.
  \item A model of repeated games, that takes into account the cost of calculations required to choose a move at the each iteration. %Модель повторяющихся игр с учётом стоимости вычислений, необходимых для выбора хода на очередной итерации.
  \item Cryptographic punishment strategies for repeated three-player even-odd, widening the set of perfect subgame equilibria with points achievable only by taking into account the complexity of the algorithms. %Криптографические стратегии наказания для повторяющегося трёхстороннего чёт-нечета, пополняющие множество совершенных подыгровых равновесий точками, не достижимыми без принятия во внимание сложности алгоритмов.
\end{enumerate}

%{\reliability} полученных результатов обеспечивается при помощи формальных доказательств с применением топологических инструментов. Результаты находятся в соответствии с результатами, полученными другими авторами.

\textbf{Work approbation.}
The main results of the work were reported on: Lomonosov readings (2017, 2021) \cite{ownlmr2017, ownlmr2021}, IX Moscow International Conference on Operations Research \cite{ownorm2018} and Conference for Young Scientists in Mathematical Economics and Economic Theory (MEET-2021) \cite{meet2021}. %Основные результаты работы докладывались~на: Ломоносовских чтениях (2017, 2021 гг.) \cite{ownlmr2017, ownlmr2021}, IX Московской международной конференции по исследованию операций \cite{ownorm2018} и конференции молодых учёных по математической экономике и экономической теории (MEET-2021) \cite{meet2021}.

\textbf{Personal contribution.} The author independently obtained the results featured in the dissertation work in the form of theorems and other provisions submitted for defense. The obtained results were prepared for publication without co-authors. %Результаты, представленные в диссертационной работе теоремами и другими выносимыми на защиту положениями, получены автором самостоятельно. Подготовка к публикации полученных результатов проводилась без соавторов.

\textbf{Publications.} The main results on the topic of the dissertation are presented in~7~published papers, 3 of which \cite{ownconsp}\cite{owncards}\cite{ownsched} were published in a periodical scientific journal recommended by the HAC and indexed by Web of Science and Scopus. The central work has an English translation\cite{ownconsp_eng}. 3 papers were published in conference abstracts. %Основные результаты по теме диссертации изложены в~7~печатных работах, 3 из которых \cite{ownconsp}\cite{owncards}\cite{ownsched} изданы в периодическом научном журнале, рекомендованном ВАК и индексируемом Web of Science и Scopus. Центральная работа имеет перевод на английский язык\cite{ownconsp_eng}. 3 работы изданы в тезисах докладов.
 % Характеристика работы по структуре во введении и в автореферате не отличается (ГОСТ Р 7.0.11, пункты 5.3.1 и 9.2.1), потому её загружаем из одного и того же внешнего файла, предварительно задав форму выделения некоторым параметрам

\textbf{Volume and structure of work.} The dissertation consists of introduction, 3 chapters, conclusion and 4 appendices. The full volume of the dissertation is 86 pages, including 2 figures and 5 tables. The list of references contains 26 titles.
%\textbf{Объем и структура работы.} Диссертация состоит из~введения,
%\formbytotal{totalchapter}{глав}{ы}{}{},
%заключения и
%\formbytotal{totalappendix}{приложен}{ия}{ий}{}.
%%% на случай ошибок оставляю исходный кусок на месте, закомментированным
%%Полный объём диссертации составляет  \ref*{TotPages}~страницу
%%с~\totalfigures{}~рисунками и~\totaltables{}~таблицами. Список литературы
%%содержит \total{citenum}~наименований.
%%
%Полный объём диссертации составляет
%\formbytotal{TotPages}{страниц}{у}{ы}{}, включая
%\formbytotal{totalcount@figure}{рисун}{ок}{ка}{ков} и
%\formbytotal{totalcount@table}{таблиц}{у}{ы}{}.
%Список литературы содержит
%\formbytotal{citenum}{наименован}{ие}{ия}{ий}.
    % Введение
\ifnumequal{\value{contnumfig}}{1}{\counterwithout{figure}{chapter}
}{\counterwithin{figure}{chapter}}
\ifnumequal{\value{contnumtab}}{1}{\counterwithout{table}{chapter}
}{\counterwithin{table}{chapter}}
\chapter{Модель заговоров}\label{ch:ch1}

\section{Коррелированное расширение игры в нормальной форме}\label{sec:ch1/sec1}

Базовым формализмом для описания дополнительной информационной асимметрии в сложившейся научной традиции выступает коррелированное расширение нормальной формы игр, предложенное Робертом Ауманом в \cite{Aumann74}. Для удобства его центральные элементы будут изложены здесь в нотации, адаптированной к русскоязычной среде. Рассмотрим игру в нормальной форме $\Gamma = \langle A, S^a, u^a(s), a \in A \rangle$. Конечное множество игроков здесь и далее везде обозначается как $A = \{1, \ldots, m\}$, а конечное множество наборов чистых стратегий "--- $S = S^1 \times \ldots \times S^m$. Помимо множества стратегий $S^a$, каждый игрок определяется платёжной функцией $u^a : S \rightarrow \mathbb{R}$.

Также рассмотрим вероятностное пространство\cite{Kolmogorov} $\langle \Omega, \mathfrak{B}, \mathbb{P} \rangle$, в котором реализуется наблюдаемое игроками состояние природы. Здесь $\Omega$ "--- множество всевозможных таких состояний, $\mathfrak{B}$ "--- $\sigma$"~алгебра подмножеств $\Omega$, а $\mathbb{P} : \mathfrak{B} \rightarrow \mathbb{R}_{\ge 0}$ "--- вероятностная мера. Каждому игроку $a \in A$ поставим в соответствие \emph{собственное подпространство} $\langle \Omega, \mathfrak{I}^a, \mathbb{P} \rangle$ такое, что $\mathfrak{I}^a \subseteq \mathfrak{B}$. При этом набор $\sigma$"~алгебр $\mathfrak{I} = (\mathfrak{I}^a, a \in A)$ отражает информированность игроков о состоянии природы. В описываемой ситуации это состояние не влияет на функции выигрышей непосредственно, выступая исключительно как способ синхронизации действий игроков. Это значит, что $\sigma$"~алгебра $\mathfrak{B}$ сама по себе не является существенным параметром модели, и измеримость по ней для $\mathbb{P}$ можно заменить измеримостью по $\mathfrak{I}^a, \forall a \in A$.

Отдельно следует заметить, что в оригинале для собственных подпространств игроков Аумановский формализм предполагал индивидуальность не только $\sigma$"~алгебр, но и соответствующих им мер, учитывая тем самым возможную субъективность оценок вероятности наступления тех или иных событий, что немаловажно в случаях, когда в качестве механизма корреляции выступают процессы, слишком сложные для объективного анализа (например спортивные соревнования). Однако, для целей данного исследования этот аспект не имеет большого смысла, поскольку в модели заговоров подразумевается, что их участники могут произвольным образом выбирать механизм корреляции, а в такой ситуации разумно ожидать, что люди будут использовать простые источники случайности с известным распределением (рулетки, кости, жребий и т.д.). По этой причине здесь и далее модель коррелированных стратегий используется в её упрощённой форме, с общей для всех игроков объективной вероятностной мерой в пространстве состояний природы.

Таким образом, получается $\Phi = \langle A, \Omega, \mathfrak{I}^a, \mathbb{P}, a \in A \rangle$ "--- набор параметров, характеризующих некоторое \emph{пространство корреляции} для произвольной игры с множеством игроков $A$. Отметим, что в играх с одним множеством игроков, но различными множествами чистых стратегий и функциями выигрыша можно применять одно и то же пространство корреляции. Полностью же \emph{коррелированное расширение игры} определяет пара $\Gamma | \Phi$. Опишем полученную новую игру в терминах нормальной формы\footnote{Следуя нотации, введённой в \cite{Aumann74}, наборы стратегий и множества исходов коррелированного расширения игры обозначаются жирным шрифтом ($\mathbf{s}$ и $\mathbf{S}$ там, где в базовой игре $s$ и $S$).}:
\begin{equation*}
	\Gamma | \Phi = \langle A, \mathbf{S}^a, u^a(\mathbf{s}), a \in A \rangle.
\end{equation*}

Здесь множество $\mathbf{S}^a$ доступных игроку $a$ коррелированных стратегий состоит из всех $\mathfrak{I}^a$"~измеримых функций $\mathbf{s}^a : \Omega \rightarrow S^a$, отображающих множество возможных состояний природы на множество доступных ему чистых стратегий. Соответственно, функция выигрыша вычисляется по формуле математического ожидания случайной величины
\begin{equation*}
	u^a(\mathbf{s}) = \sum\limits_{s \in S} \mathbb{P}(\mathbf{s}^{-1}(s)) u^a(s),\;\mathbf{s}^{-1}(s) = \{\omega \in \Omega \mid \mathbf{s}^a(\omega) = s^a, \forall a \in A\},
\end{equation*}
где $\mathbb{P}(\mathbf{s}^{-1}(s))$ выступают в роли коэффициентов распределения на матрице игры.

В своей статья Ауман демонстрирует силу вводимого формализма, показывая на примерах, как с его помощью в играх можно получать новые точки равновесия по Нэшу. Подбирая параметры пространств корреляции, можно конструировать не только решения с любыми выплатами из выпуклой оболочки векторов выплат в точках классического смешанного равновесия Нэша, но для некоторых игр даже решения, лежащие за пределами такой выпуклой оболочки. Это позволяет сформулировать ключевое для данной работы понятие:
\begin{definition}
	Пусть $\Gamma$ --- игра в нормальной форме c $m$ участниками, а $U \subseteq \mathbb{R}^m$ --- множество всех векторов выплат, достижимых в её смешанных равновесиях по Нэшу. Игра $\Gamma$ называется \emph{чувствительной к дополнительной информационной асимметрии}, когда существует пространство корреляции $\Phi$ такое, что в игре $\Gamma | \Phi$ найдётся коррелированное равновесие по Нэшу с вектором выплат, не принадлежащим выпуклой оболочке множества $U$.
\end{definition}

Кроме того, в этой же статье доказывается, что наличие в пространстве корреляции публичной вещественной рулетки, т.е. подпространства событий с равномерно распределённым вещественным исходом в диапазоне $\left[0, 1\right)$, влечёт выпуклость как множества достижимых выплат, так и множества равновесий Нэша в любой игре. Касательно коррелированного расширения игр в нормальной форме, вышеизложенного вполне достаточно для понимания идей данной работы, более подробно же современное состояние знаний на эту тему можно проследить по публикациям, упомянутым в Приложении~\cref{app:A}.

\section{Изоморфизм пространств корреляции}\label{sec:ch1/sec2}

Следует отметить, что модель пространств корреляции в некотором смысле существенно избыточна, поскольку как таковые события из состояния природы значения не имеют и используются лишь в качестве сигналов для синхронизации стратегий игроков. Чтобы осмысленным образом рассуждать о влиянии, оказываемом дополнительной информационной асимметрией на исходы конфликтов, неизбежно требуется умение абстрагироваться от конкретных её источников, фокусируя внимание на структурных различиях в осведомлённости оппонентов. Если формально различные пространства корреляции оказываются полностью взаимозаменимы с теоретико"~игровой точки зрения, то они должны быть отнесены к общему классу эквивалентности, чьё описание и является в действительности существенным параметром модели. Причём сразу следует заметить, что это касается не только тривиальных замен множества состояний природы на другое множество той же мощности с соответствующей биекцией остальных параметров пространства, но и более сложных случаев. Например, если в контексте некоторой игры группа игроков наблюдает общий сигнал в виде колеса вещественной рулетки, будет ли иметь значение наблюдение ими ещё и броска монетки? Здравый смысл подсказывает, что любую общую стратегию с использованием рулетки и монетки можно легко превратить в эквивалентную для одной только рулетки, для чего достаточно поделить колесо пополам и отобразить отдельные варианты для орла и решки на полученные два сектора. Опишем этот феномен в виде изоморфизма:
\begin{definition}
	Разбиением пространства корреляции $\Phi = \langle A, \Omega, \mathfrak{I}^a, \mathbb{P}, a \in A \rangle$ в произвольное конечное множество исходов (кодомен) $X = X^1 \times ... \times X^m$ называется отображение $f : \Omega \rightarrow X$, состоящее из набора функций $(f^1, ..., f^m)$, где каждая $f^a : \Omega \rightarrow X^a$ измерима в $\mathfrak{I}^a$. Далее разбиение $f$ пространства корреляции $\Phi$ будем сокращённо обозначать $f \models \Phi$.
\end{definition}

В контексте коррелированного расширения множествам исходов $X^a$ соответствуют множества чистых стратегий $S^a$, а элементам разбиения $f^a$ "--- коррелированные стратегии $\mathbf{s}^a$. Далее также будут использоваться отображения $f^{-1} : X \rightarrow 2^\Omega$, обратные к разбиениям пространств корреляции:
\begin{equation*}
	f^{-1}(x) = \bigcap_{a \in A} (f^a)^{-1}(x^a).
\end{equation*}
\begin{definition}
	Пространство $\Phi_1$ с мерой $\mathbb{P}_1$ называется отобразимым на $\Phi_2$ с мерой $\mathbb{P}_2$ (далее $\Phi_1 \precsim \Phi_2$), если их множества игроков совпадают и для любого разбиения $f_1 \models \Phi_1$ существует разбиение $f_2 \models \Phi_2$ с тем же кодоменом такое, что $\mathbb{P}_1 \circ f_1^{-1} = \mathbb{P}_2 \circ f_2^{-1}$. Взаимно отобразимые друг на друга пространства корреляции называются изоморфными (далее $\Phi_1 \sim \Phi_2$).
\end{definition}

Это определение легко проиллюстрировать упомянутым выше примером "--- для любого разбиения $f_1 : \left[0, 1\right) \times \{0, 1\} \rightarrow X$ пространства, состоящего из вещественной рулетки и симметричной монетки, можно построить соответствующий образ $f_2 : \left[0, 1\right) \rightarrow X$ в пространстве из одной только рулетки:
\begin{equation*}
	f_2(\alpha) = \begin{cases}
		f_1(2 \alpha, 0), & 0 \le \alpha < \frac{1}{2} \\
		f_1(2 \alpha - 1, 1), & \frac{1}{2} \le \alpha < 1
	\end{cases}.
\end{equation*}

Рефлексивность, симметричность и транзитивность вводимого при помощи данного определения изоморфизма очевидны, а значит это действительно отношение эквивалентности на множестве пространств корреляции. При этом, хотя определение изоморфизма дано в отрыве от коррелированного расширения игр, можно сформулировать следующую теорему:
\begin{definition}
	Для игры $\Gamma = \langle A, S^a, u^a(s), a \in A \rangle$ множеством достижимых выплат по отклонениям группы игроков $A_*$ от профиля стратегий $s$ будем называть
	\begin{equation*}
		U_\Gamma^{A_*}(s) = \{\bar{u} \mid \exists s_* \in S : u(s_*) = \bar{u}, \forall a \in A \setminus A_*, s^a = s_*^a\}.
	\end{equation*}
\end{definition}
\begin{theorem}[Об изоморфных пространствах]
	Пусть $\Phi_1 \sim \Phi_2$. Тогда для любой игры в нормальной форме $\Gamma$ с конечными множествами стратегий игроков её коррелированные расширения $\Gamma | \Phi_1$ и $\Gamma | \Phi_2$ обладают следующим свойством. Пусть $\mathbf{s}_1$ "--- некоторый профиль стратегий игры $\Gamma | \Phi_1$. Тогда существует $\mathbf{s}_2$ "--- профиль стратегий игры $\Gamma | \Phi_2$ такой, что $U_{\Gamma | \Phi_1}^{A_*}(\mathbf{s}_1) = U_{\Gamma | \Phi_2}^{A_*}(\mathbf{s}_2)$ для любой группы игроков $A_*$.
\end{theorem}

Эта теорема позволяет считать изоморфные пространства корреляции неразличимыми в контексте поиска равновесий, устойчивых к как индивидуальным, так и групповым отклонениям. Доказательство, представляющее собой упражнение в топологии без тесной связи с центральными идеями работы, вынесено в Приложение~\cref{app:B}.

\section{Пространства заговоров}\label{sec:ch1/sec3}

Получив осмысленный изоморфизм для пространств корреляции, можно выделить из всевозможных классов эквивалентности те, что представляют интерес с точки зрения моделирования дополнительной информационной асимметрии. Как показал Ауманн, выход за пределы выпуклой оболочки множества решений в смешанных стратегиях возможен, если часть игроков (не менее двух) использует коррелированную стратегию, зависящую от события, о котором не проинформирован хотя бы один из остальных игроков. Подобную форму взаимовыгодного тайного согласования действий естественно называть \emph{заговором}, а используемый для синхронизации сигнал "--- \emph{тайной}. Пусть в пространстве корреляции $\Phi = \langle A, \Omega, \mathfrak{I}^a, \mathbb{P}, a \in A \rangle$ имеется тайна, т.е. вероятностное подпространство $\langle \Omega, \mathfrak{S}, \mathbb{P} \rangle$. Искомая асимметрия информированности предполагает, что некоторые игроки наблюдают события из $\mathfrak{S}$ (или другие, коррелирующие с ними), а некоторые "--- нет. Хотя теоретически можно представить заговор, степень вовлеченности в который варьируется от игрока к игроку (кто-то может наблюдать события из $\mathfrak{S}$ частично или опосредованно, через наблюдение других коррелирующих с ними событий), имеет смысл в первую очередь рассмотреть простейший случай "--- деление всех игроков на <<заговорщиков>> $A_{\mathfrak{S}} \subseteq A$, наблюдающих $\mathfrak{S}$ целиком, и <<аутсайдеров>> $A \setminus A_{\mathfrak{S}}$, в чьём поле зрения только события, не коррелирующие с элементами $\mathfrak{S}$.

Следующий логичный вопрос "--- о структуре самой тайны. Вполне можно представить, как заговорщики используют в её роли самые разные источники случайности: броски игральных костей, тасование колод карт, лотерейные розыгрыши и т.д., так что на первый взгляд неочевидно, можно ли ограничиться рассмотрением какого-то одного, естественного в контексте происходящего механизма. Положительный ответ можно получить, используя введённые выше отображения пространств корреляции. Если мы будем сравнивать всевозможные пространства, различающиеся только тайнами группы заговорщиков $A_* \subseteq A$, отношение $\precsim$ вводит на их множестве частичный порядок. Нижней гранью этого порядка будет вырожденное пространство корреляции, в котором тайна заговора состоит из единственного атомарного события с вероятностью $1$ "--- такое пространство отобразимо на любое другое и, очевидно, вообще не может быть использовано для корреляции стратегий. Верхняя грань интереснее "--- в её типичном представителе тайна заговора представляет собой произвольное безатомическое \cite[с.~81]{Bogachev} пространство. Проще всего представить такой механизм корреляции в виде вещественной рулетки, вращение которой генерирует равномерно распределённую случайную величину в единичном полуинтервале $\left[0, 1\right)$. Наблюдающие рулетку заговорщики могут, разделяя её на сектора необходимых размеров, согласовать любой набор коррелированных стратегий в играх с конечным множеством исходов. В первую очередь такой универсальный источник случайности и имеет смысл рассматривать как предоставляющий максимум свободы выбора.

Наконец, следует подумать о пространствах корреляции с множественными тайнами. По сути, ничего не мешает игрокам наблюдать сразу несколько рулеток, выбирая в зависимости от ситуации, с кем из оппонентов коррелировать свою стратегию. Более того, стратегия игрока может одновременно существенно зависеть от более чем одной тайны. Таким образом, естественным предметом рассмотрения можно считать пространства корреляции, состоящие из наборов независимых вещественных рулеток, каждая из которых характеризуется подмножеством игроков, имеющих возможность её наблюдать. Остаётся заметить, что при наличии в одном пространстве корреляции двух и более вещественных рулеток, принадлежащих одному и тому же кругу заговорщиков, все кроме одной можно без ущерба для модели выкинуть, так как в играх с конечными множествами исходов такое дублирование, очевидно, бесполезно. Перейдём теперь к более формальному определению предложенной концепции. Для этого рассмотрим произвольное пространство корреляции $\Phi = \langle A, \Omega, \mathfrak{I}^a, \mathbb{P}, a \in A \rangle$. В этом пространстве для каждой непустой группы игроков $A_* \subseteq A$ определим следующее семейство событий\footnote{Здесь и далее под $\sigma(\mathfrak{X})$ понимается пополнение семейства множеств $\mathfrak{X}$ до $\sigma$"~алгебры}:
\begin{equation*}
	\mathfrak{S}_\Phi^{A_*} = \{U \in \bigcap\limits_{a \in A_*} \mathfrak{I}^a \mid \mathbb{P}(U \cap V) = \mathbb{P}(U) \mathbb{P}(V), \forall V \in \sigma(\bigcup\limits_{a \in A \setminus A_*} \mathfrak{I}^a)\}.
\end{equation*}

Таким образом, $\mathfrak{S}_\Phi^{A_*}$ "--- множество всех таких событий, что о них осведомлены все члены $A_*$, и каждое событие попарно независимо со всеми событиями, известными не членам $A_*$ даже при объединении их знаний. Поскольку пересечение $\sigma$"~алгебр образует $\sigma$"~алгебру, и так как подмножество независимых с некоторым событием событий $\sigma$"~алгебры также образует $\sigma$"~алгебру, то $\mathfrak{S}_\Phi^{A_*}$ "--- $\sigma$"~алгебра. Это позволяет говорить о вероятностном подпространстве $\langle \Omega, \mathfrak{S}_\Phi^{A_*}, \mathbb{P} \rangle$, которое и называется тайной группы $A_*$. Выделим два случая: тайны с безатомическими мерами мы будем называть \emph{полными}, а тайны с тривиальными атомическими мерами с единственным атомом $\Omega$ "--- \emph{пустыми}. Это даёт возможность сформулировать следующее:
\begin{definition}
	Пространство корреляции $\langle A, \Omega, \mathfrak{I}^a, \mathbb{P}, a \in A \rangle$ назовём пространством заговоров структуры $\mathfrak{A} = \{A_1, ..., A_n\} \subseteq 2^A$, когда
	\begin{itemize}
		\item $\bigcup\limits_{i=0}^n A_i = A$;
		\item $\forall A_* \in \mathfrak{A}$ тайна $A_*$ полна;
		\item $\forall A_* \notin \mathfrak{A}$ тайна $A_*$ пуста;
		\item $\mathfrak{I}^a = \sigma(\bigcup\limits_{a \in A_* \in \mathfrak{A}}\mathfrak{S}_\Phi^{A_*})$, т.е. $\mathfrak{I}^a$ "--- наименьшая $\sigma$"~алгебра, включающая все $\sigma$"~алгебры тайн групп, в которые входит игрок $a$.
	\end{itemize}
\end{definition}

Проще говоря, пространствами заговоров называются такие пространства корреляции, в которых а) тайна любой группы игроков либо полна, либо пуста; б) каждый игрок входит хотя бы в одну группу с полной тайной и в) у игроков нет никаких знаний о состоянии природы, которые не порождались бы тайнами групп, которым они принадлежат. Построим для иллюстрации простейший пример такого пространства:
\begin{itemize}
	\item $A = \{1, 2, 3\}$,
	\item $\Omega = \left[0, 1\right)^2$,
	\item $\mathfrak{I}^1 = \sigma(\{\left[ 0, p_1 \right) \times \left[ 0, 1 \right) \mid 0 < p_1 \leq 1 \})$,
	\item $\mathfrak{I}^2 = \sigma(\{\left[ 0, 1 \right) \times \left[ 0, p_2 \right) \mid 0 < p_2 \leq 1 \})$,
	\item $\mathfrak{I}^3 = \sigma(\{\left[ 0, p_1 \right) \times \left[ 0, p_2 \right) \mid 0 < p_1 \leq 1, 0 < p_2 \leq 1 \})$,
	\item $\mathbb{P}$ "--- мера Лебега.
\end{itemize}

В этом примере пространство корреляции состоит из двух независимых вещественных рулеток, первый и второй игроки наблюдают по одной из них, а третий наблюдает обе. При этом выходит, что $\mathfrak{S}_\Phi^{\{1,3\}}$ совпадает с $\mathfrak{I}^1$, $\mathfrak{S}_\Phi^{\{2,3\}}$ совпадает с $\mathfrak{I}^2$, а для остальных групп $A_* \subseteq A$ соответствующая $\mathfrak{S}_\Phi^{A_*}$ тривиальна. Структурой пространства (или \emph{семейством заговоров}) называется множество всех групп игроков с полными тайнами. В вышеприведённом примере структура пространства $\mathfrak{A} = \{\{1,3\},\{2,3\}\}$. С точки зрения допустимых профилей стратегий это означает, что любая группа игроков, входящая в семейство заговоров, может использовать общую тайну для формирования коррелированной стратегии, причём игроки, не входящие в эту группу, не могут присоединиться к согласованному таким образом выбору стратегий. Напротив, группы игроков, не входящие в семейство заговоров, вышеописанной возможностью не располагают. Структуру пространства можно считать его исчерпывающим конечным описанием в силу истинности следующего утверждения:
\begin{theorem} \label{the:struct}
	Все пространства заговоров одной структуры изоморфны.
\end{theorem}

Доказательство этой теоремы снова представляет собой упражнение в топологии без тесной связи с основными идеями работы и вынесено в приложение \ref{app:C}. Теперь, когда установлено, что множество всех пространств заговоров разбивается на классы эквивалентности, нетрудно предложить способ конструирования стандартного представителя каждого класса по соответствующему семейству заговоров:
\begin{definition}
	Стандартным пространством структуры $\mathfrak{A} = \{A_1, A_2, ..., A_n\}$ называется пространство корреляции $\Phi_{\mathfrak{A}} = \langle A, \Omega, \mathfrak{I}^a, \mathbb{P}, a \in A \rangle$ со следующими параметрами:
	\begin{itemize}
		\item $A = \bigcup\limits_{i=1}^n A_i$,
		\item $\Omega = \left[0, 1\right)^n$,
		\item $\mathfrak{I}^a = \sigma(\{\prod\limits_{i=1}^n\left[ 0, p_i \right) \mid \operatorname{\mathbf{if}} \: a \in A_i \: \operatorname{\mathbf{then}} \: 0 < p_i \leq 1 \: \operatorname{\mathbf{else}} \: p_i = 1 \})$,
		\item $\mathbb{P}$ "--- мера Лебега.
	\end{itemize}
\end{definition}

Множество состояний природы представляет собой $n$"~мерный (по числу заговоров, входящих в семейство) единичный куб, а вероятностная мера соответствует непрерывному равномерному распределению. При этом $\sigma$"~алгебра каждого игрока борелева в проекциях на оси, соответствующие заговорам в которые он входит, и тривиальна в проекциях на остальные оси. Выделение стандартного представителя для любых семейств заговоров позволяет использовать нотацию $\Gamma | \mathfrak{A}$, под которой в дальнейшем будет пониматься $\Gamma | \Phi_{\mathfrak{A}}$. Эта нотация подчёркивает тот факт, что выбор конкретного пространства корреляции среди всех пространств заговоров необходимой структуры для нас значения не имеет, а стандартное пространство выступает в роли простейшего представителя, пригодного для практических вычислений.

\section{Трёхсторонний чёт-нечет}\label{sec:ch1/sec4}

В качестве элементарного примера конфликта, чувствительного к дополнительной информационной асимметрии, может выступать <<трёхсторонний чёт"~нечет>>. В этой игре каждый из трёх участников тайно выбирает <<орла>> или <<решку>> на своей монете и прижимает её к столу соответствующей стороной вверх, после чего все одновременно поднимают ладони и в зависимости от сложившейся комбинации делят фиксированный банк. Когда все три монеты лежат одной и той же стороной, раунд считается сыгранным вничью и игроки делят банк поровну. Если же совпали только две из них, то оказавшийся в меньшинстве игрок считается проигравшим и не получает доли при дележе банка. В матричной форме это можно описать так:
\begin{table} [htbp]
	\centering
	\begin{threeparttable}
		\caption{Трёхсторонний чёт"~нечет}
		\label{tab:coin3}
		\begin{tabular}{ |c|c|c|c|c| }
			\cline{1-2} \cline{4-5}
			\rule[-7pt]{0pt}{2em}$4, 4, 4$ &
			\rule[-7pt]{0pt}{2em}$6, 0, 6$ & \qquad\qquad\qquad &
			\rule[-7pt]{0pt}{2em}$6, 6, 0$ &
			\rule[-7pt]{0pt}{2em}$0, 6, 6$ \\
			\cline{1-2} \cline{4-5}
			\rule[-7pt]{0pt}{2em}$0, 6, 6$ &
			\rule[-7pt]{0pt}{2em}$6, 6, 0$ & \qquad\qquad\qquad &
			\rule[-7pt]{0pt}{2em}$6, 0, 6$ &
			\rule[-7pt]{0pt}{2em}$4, 4, 4$ \\
			\cline{1-2} \cline{4-5}
		\end{tabular}
	\end{threeparttable}
\end{table}

В таблице \ref{tab:coin3} первый игрок выбирает строку, второй "--- столбец, а третий "--- матрицу. Решением этой игры в чистых стратегиях являются два равновесия Нэша, соответствующие синхронным выборам одинаковых сторон всеми игроками. В смешанных стратегиях добавляется ещё одно вырожденное решение, когда каждый игрок делает случайный выбор между орлом и решкой с равными вероятностями. Все эти решения, очевидно, дают математическое ожидание платежей равное $(4,4,4)$. В рамках классической теории игр этим анализ конфликта и исчерпывается, однако добавление фактора информационной асимметрии делает ситуацию интереснее. Рассмотрим эту же игру в пространстве заговоров структуры $\{\{1,2\}\}$, т.е. в ситуации, когда игроки 1 и 2 имеют возможность использовать согласованные втайне от игрока 3 коррелированные стратегии. Пусть $\alpha \in \left[0, 1\right)$ "--- значение соответствующей секретной рулетки. Заговорщики могут использовать стратегии вида $\mathbf{s}^1, \mathbf{s}^2 : \left[0, 1\right) \rightarrow P_{\{\mathbf{head}, \mathbf{tail}\}}$, где под $P_{\{\mathbf{head}, \mathbf{tail}\}}$ понимаются всевозможные вероятностные меры на множестве $\{\mathbf{head}, \mathbf{tail}\}$, т.е. множество классических смешанных стратегий рассматриваемой игры. Аутсайдер же вынужден довольствоваться стратегиями $\mathbf{s}^3 \in P_{\{\mathbf{head}, \mathbf{tail}\}}$, поскольку не имеет доступа к рулетке заговора. Для того, чтобы оказаться в более выгодной по сравнению с равным дележом точке равновесия, игроки 1 и 2 могут выбрать любые стратегии, обеспечивающие им равновероятный синхронный выбор стороны монетки:
\begin{equation*}
	\mathbf{s}^1(\alpha) = \mathbf{s}^2(\alpha) = \begin{cases}
		(1, 0), & \alpha < \frac{1}{2},\\
		(0, 1), & \alpha \ge \frac{1}{2}.
	\end{cases}
\end{equation*}

При этом любая смешанная стратегия игрока 3 в силу независимости с рулеткой заговора обеспечивает ему совпадение с остальными ровно в половине случаев. Выплаты в сложившейся ситуации равны $(5,5,2)$, причём ни для одного из игроков нет выгодного индивидуального отклонения.

\section{Необходимая сложность модели заговоров}\label{sec:ch1/sec5}

Классический формализм матричной игры в нормальной форме отождествляет профиль чистых стратегий и исход розыгрыша "--- они являются буквально одним и тем же математическим объектом. Почти то же самое можно сказать и о его смешанном расширении "--- пространство профилей смешанных стратегий изоморфно пространству исходов, состоящему из вероятностных распределений на матрице игры, соблюдающих условие независимости выбора строк и столбцов. Увы, коррелированное расширение в формулировке Роберта Аумана ломает эту идиллическую картину. Пространство исходов розыгрышей в нём даже проще, чем в смешанном случае "--- распределение вероятностей на множестве элементов игровой матрицы может быть любым, без дополнительных условий. Но вот со стратегическими профилями всё резко становиться сложнее "--- стратегия каждого игрока представляет собой функцию, отображающую множество состояний природы в множество его чистых стратегий, причём с соблюдением измеримости по $\sigma$"~алгебре его информированности. Очевидно, что более ни о каком взаимно однозначном соответствии между профилями и исходами не может идти и речи "--- в зависимости от параметров пространства корреляции некоторые исходы могут оказаться вообще не достижимы, а равновероятные события можно без влияния на исход менять местами в домене стратегий. Вдобавок, работая с коррелированными стратегиями в формулировке Аумана, можно столкнуться с многомерными пространствами событий, борелевскими $\sigma$"~алгебрами и другими нетривиальными феноменами колмогоровской теории вероятностей, что, вероятно, тоже вносит свой вклад в то, сколь неохотно специалисты теории игр прибегают к этому инструменту в более прикладных исследованиях.

Предложенная здесь модель заговоров сужает коррелированное расширение, выделяя из континуума всевозможных пространств корреляции конечный (для любого конечного числа игроков) набор, по одному на каждую структуру заговоров. На первый взгляд такое радикальное упрощение пространства параметров даёт надежду на то, что и при практическом использовании модели удастся обойтись без <<эзотерических>> аспектов теории меры. В идеале хотелось бы, чтобы в играх с заговорами можно было тем или иным способом отождествить пространство стратегических профилей и множество исходов, так же, как это происходит в классических формализмах. Если бы мы могли, глядя только на распределение вероятностей реализации отдельных исходов в игровой матрице, определить, какой профиль стратегий (или любой представитель семейства неразличимых профилей) был сыгран игроками, то это означало бы, что мы можем без погружения в детали аумановской модели корреляции найти также и все распределения, достижимые в тех или иных отклонениях от сыгранного профиля стратегий, проверяя тем самым ситуацию на равновесность.

Увы, в общем случае это вряд ли возможно "--- наглядно показать потерю важной информации при переходе от наборов стратегий к результату корреляции можно при помощи того же самого трёхстороннего чёт-нечета, описанного в предыдущем разделе. Представим себе, что игра ведётся в пространстве заговоров $\{\{1, 2\}, \{1, 2, 3\}\}$, то есть, игроки 1 и 2 могут коррелировать свои действия как втайне от игрока 3, так и вместе с ним. Пространство корреляции при этом оказывается состоящим из двух рулеток $\alpha^{1, 2}$ и $\alpha^{1, 2, 3}$, наблюдаемых игроками, указанными в верхних индексах их обозначения. Рассмотрим два набора стратегий: первый "---
\begin{equation*}
	\mathbf{s}^1(\alpha^{1, 2}, \alpha^{1, 2, 3}) = \mathbf{s}^2(\alpha^{1, 2}, \alpha^{1, 2, 3}) = \begin{cases}
		(1, 0), & \alpha^{1, 2, 3} < \frac{1}{2},\\
		(0, 1), & \alpha^{1, 2, 3} \ge \frac{1}{2},
	\end{cases} \mathbf{s}^3(\alpha^{1, 2, 3}) = (\frac{1}{2}, \frac{1}{2}),
\end{equation*}
и второй "---
\begin{equation*}
	\mathbf{s}^1(\alpha^{1, 2}, \alpha^{1, 2, 3}) = \mathbf{s}^2(\alpha^{1, 2}, \alpha^{1, 2, 3}) = \begin{cases}
		(1, 0), & \alpha^{1, 2} < \frac{1}{2},\\
		(0, 1), & \alpha^{1, 2} \ge \frac{1}{2},
	\end{cases} \mathbf{s}^3(\alpha^{1, 2, 3}) = (\frac{1}{2}, \frac{1}{2}).
\end{equation*}

В обоих случаях все игроки делают равновероятный выбор между орлом и решкой, причём первые два игрока делают его синхронно, а третий "--- независимо от них. Если записать распределение вероятностей отдельных исходов в матрице $2 \times 2 \times 2$, соответствующей игровой из таблицы \ref{tab:coin3}, то для обоих наборов это будет выглядеть так:
\begin{tabular}{ |c|c|c|c|c| }
	\cline{1-2} \cline{4-5}
	$\frac{1}{4}$ &
	$0$ & \qquad &
	$\frac{1}{4}$ &
	$0$ \\
	\cline{1-2} \cline{4-5}
	$0$ &
	$\frac{1}{4}$ & \qquad &
	$0$ &
	$\frac{1}{4}$ \\
	\cline{1-2} \cline{4-5}
\end{tabular}.
Понятное дело, платежи здесь также равны и составляют $(5, 5, 2)$. Тем не менее, если присмотреться, можно обнаружить, что в первом случае игроки 1 и 2 используют для синхронизации общую с игроком 3 рулетку $\alpha^{1, 2, 3}$, что позволяет ему, изменив свою стратегию на
\begin{equation*}
	\mathbf{s}^3(\alpha^{1, 2, 3}) = \begin{cases}
		(1, 0), & \alpha^{1, 2, 3} < \frac{1}{2},\\
		(0, 1), & \alpha^{1, 2, 3} \ge \frac{1}{2},
	\end{cases}
\end{equation*}
присоединиться к ним и перейти в исход с распределением вероятностей
\begin{tabular}{ |c|c|c|c|c| }
	\cline{1-2} \cline{4-5}
	$\frac{1}{2}$ &
	$0$ & \qquad &
	$0$ &
	$0$ \\
	\cline{1-2} \cline{4-5}
	$0$ &
	$0$ & \qquad &
	$0$ &
	$\frac{1}{2}$ \\
	\cline{1-2} \cline{4-5}
\end{tabular},
где выплаты составляют $(4, 4, 4)$. Таким образом, первая точка не является равновесием Нэша. Во втором же случае игроки 1 и 2 используют для синхронизации приватную рулетку $\alpha^{1, 2}$, исход которой игрок 3 наблюдать не может. Какую бы стратегию он ни выбрал, совпадёт со стратегиями заговорщиков она ровно в половине случаев, а значит выплаты не изменятся. Так как у остальных игроков выгодных отклонений тоже нет, эта ситуация уже будет равновесием Нэша.

Таким образом, мы построили пример, в котором одному и тому же распределению вероятностей исходов в матрице игры соответствуют как минимум два настолько различных набора стратегий, что один из них является равновесием Нэша, а второй "--- нет. Очевидно, что никаким разумным способом объединить эти наборы в одном классе эквивалентности нельзя, а значит отождествление игровых ситуаций и исходов розыгрыша в общем случае невозможно, как бы мы ни старались. Следует, впрочем, заметить, что в более простых случаях, когда каждый игрок может входить не более чем в один заговор, построить такой наглядный и тривиальный контрпример уже не получится "--- если корреляция, проявляющаяся в распределении вероятностей на матрице игры, могла бы быть получена только одним способом, то при решении практических задач ничего не мешает оперировать именно распределениями вероятностей, а не функциями, отображающими сигналы в исходы. Однако, если рассматривать теорию заговоров как способ моделирования информационных асимметрий, складывающихся в ходе реальных процессов взаимодействия независимых агентов в неконтролируемой среде, то подобное самоограничение, увы, загнало бы нас в рамки, не имеющие естественного обоснования.

\FloatBarrier
           % Глава 1
\chapter{Коллективная рациональность в играх с заговорами}\label{ch:ch2}

\section{Проблема планирования заданий}\label{sec:ch2/sec1}

Концепция равновесия по Нэшу, будучи фундаментом классических моделей теории игр, сама по себе зачастую оказывается недостаточно сильным формализмом. При анализе многосторонних конфликтов нередко возникают ситуации, когда множество равновесий по Нэшу слишком велико, чтобы считать его полноценным решением игры. В этом случае на помощь исследователям приходят критерии коллективной рациональности "--- выбирая среди равновесных точек оптимальные по Парето или Слейтеру, иногда можно значительно сузить пространство решений за счёт вполне естественного исключения заведомо невыгодных для всех участников. Присутствие дополнительной информационной асимметрии при этом создаёт дополнительные затруднения, поскольку предлагаемая модель подразумевает неизбежный элемент антагонизма интересов "--- в новых точках равновесия увеличение выигрышей участвующих в корреляции происходит за счёт сокращения выплат тех, кто не может к ней присоединиться, тем самым делая непродуктивными обычные критерии коллективной рациональности. Для демонстрации этого эффекта рассмотрим тривиальное обобщение классической проблемы планирования заданий \cite{Koutsoupias}, занимающей важное место в концептуальном ландшафте теории игр. При помощи этого конфликта хорошо иллюстрируются понятия \emph{цены анархии} и \emph{цены стабильности}, давая вероятно самый выразительный пример того, насколько равновесия Нэша в одной и той же игре могут различаться в смысле их глобальной оптимальности. Нам, однако, необходимо взглянуть на эту игру под другим углом, при котором понятие цен анархии и стабильности теряет смысл, уступая место чувствительности к дополнительной информационной асимметрии.

Начнём с классической модели планирования заданий. В вычислительном центре работают $m$ сотрудников, каждому из которых поручено произвести некое вычисление. В их распоряжении находятся $n$ компьютеров, на каждом из которых может быть запущена одна или несколько программ, производящих вычисления сотрудников. Машины отличаются архитектурными особенностями, что задаётся матрицей констант $t_i^a \ge 0$, обозначающих время выполнения программы сотрудника $a=\overline{1,\ldots,m}$ на компьютере $i=\overline{1,\ldots,n}$. Каждое вычисление может производиться только одним устройством. Несколько программ на одном компьютере выполняются последовательно, но результаты их работы выводятся одновременно после остановки последней из них. Таким образом, выгода сотрудника заключается в выборе такого компьютера для своего вычисления, что для него оказывается минимальным суммарное время выполнения всех запущенных программ. Опишем происходящее в терминах игры нормальной формы
\begin{equation}\label{intro:game}
	\Gamma = \langle A, S^a, u^a(s), a \in A \rangle
\end{equation}
с параметрами:
\begin{itemize}
	\item $A = \{ 1, \ldots, m \}$;
	\item $S^1 = \ldots = S^m = \{ 1, \ldots, n \}$;
	\item $u^a(s) = -t_{s^a}(s)$, где $t_i(s) = \sum\limits_{a \in A, s^a = i} t_i^a$.
\end{itemize}

Здесь следует заострить внимание на определении функции выплат. Выбрав $u^a(s) = -t_{s^a}(s)$, мы моделируем ситуацию, когда все вычисления должны были быть готовы ещё вчера, и сотрудники напрямую штрафуются за каждую лишнюю секунду, пока их результаты не лягут на стол начальству. Однако, можно смоделировать и менее напряжённый рабочий момент, взяв, например, ступенчатую функцию выплат:
\begin{equation*}%\label{intro:stairs}
	u^a(s) = \begin{cases}
		u^a_{GOOD}, &t_{s^a}(s) < t^a_{DEADLINE};\\
		u^a_{LATE}, &t_{s^a}(s) \ge t^a_{DEADLINE}.
	\end{cases}
\end{equation*}

При этом получается, что для каждого игрока $a$ назначен срок успешного выполнения задания $t^a_{DEADLINE}$, уложившись в который, он получает фиксированную выплату $u^a_{GOOD}$ (с учётом премии), а не уложившись "--- $u^a_{LATE}$ (обычную ставку). Можно придумать и более сложные схемы поощрения сотрудников, так что сформулируем сразу в общем виде:
\begin{equation}\label{intro:rewards}
	u^a(s) = v^a(t_{s^a}(s)),
\end{equation}
где $v^a(t)$ "--- монотонно невозрастающая функция оплаты за срочность по заданию сотрудника $a$. Любая игра в нормальной форме, построенная по схеме \ref{intro:game} с платёжной функцией вида \ref{intro:rewards}, в сущности является проблемой планирования задач. При этом условие монотонного невозрастания $v^a(t)$ необходимо, поскольку на него в явном виде опирается доказательство считающегося важным свойства этой игры "--- равенство $1$ цены стабильности \cite{Agussurja}. Напомним, в оптимизационных задачах с эгоистичными агентами цена стабильности представляет собой соотношение $\frac{t_{NASH}}{t_{BEST}}$, где $t_{NASH}$ "--- значение наилучшего из равновесий Нэша, а $t_{BEST}$ "--- значение глобально оптимального решения. Это означает, что в проблеме планирования задач среди всех ситуаций, минимизирующих время до остановки последнего компьютера, обязательно найдутся равновесия по Нэшу. Однако, в данной работе предлагается на время забыть о минимизации общей продолжительности вычислений и вместо этого проанализировать, какие новые свойства модели могут проявиться в отсутствие такого ограничения на монотонность.

\section{Штраф за индивидуализм}\label{sec:ch2/sec2}

Представим вычислительный центр с компьютерами, требующими сложного техобслуживания после смены, если на них запускали хотя бы одну задачу. Если поощрять укладывающихся в дедлайны сотрудников невзирая на это, несложно представить ситуацию, когда они, в стремлении любой ценой гарантировать себе премию, будут раскидывать задачи по неразумно большому количеству компьютеров. Перед лицом такой перспективы у руководства может возникнуть соблазн стимулировать своих сотрудников избегать систематической недозагрузки машин при помощи штрафов. Это может быть смоделировано ступенчатой функцией оплаты за срочность следующего вида:
\begin{equation*}
	v^a(t) = \begin{cases}
		\begin{aligned}
			u^a_{HAST}, \qquad & & & \ t < t^a_{BREAKAWAY} \hspace{-10pt} & ;\\
			u^a_{GOOD}, \qquad & t^a_{BREAKAWAY} \hspace{-8pt} & \le & \ t < t^a_{DEADLINE} & ;\\
			u^a_{LATE}, \qquad & t^a_{DEADLINE} & \le & \ t & .
		\end{aligned}
	\end{cases}
\end{equation*}

Здесь для каждого сотрудника $a$ зафиксирован не только дедлайн $t^a_{DEADLINE}$, к которому требуется успеть, чтобы получить премию, но и минимальная загруженность используемого компьютера $t^a_{BREAKAWAY}$, которой нужно достигнуть, чтобы не нарваться на штраф за нерациональное расходование вычислительных ресурсов (причём $u^a_{HAST} < u^a_{LATE} < u^a_{GOOD}$). Размер минимальной загруженности может устанавливаться, например, в зависимости от важности соответствующей задачи "--- если срочное получение результата окупает использование дополнительных машин, то его можно сделать пониже или вовсе приравнять к нулю. Если же наоборот задача не так уж и важна, то большая минимальная загруженность заставит соответствующего сотрудника вспомнить об интересах фирмы и скооперироваться с коллегами.

Как видно, сотрудникам в данном случае приходится руководствоваться немонотонной функцией оплаты за срочность, что создаёт некоторые эффекты, несвойственные для классической формулировки проблемы планирования заданий. Во-первых, при такой постановке ожидаемо не во всякой игре цена стабильности обязана равняться $1$. Достаточно рассмотреть игру с $2$ сотрудниками и $2$ одинаковыми компьютерами:
\begin{itemize}
	\item $t_1^1 = t_2^1 = 8,\\ t_1^2 = t_2^2 = 2$;
	\item $t^1_{DEADLINE} = t^2_{DEADLINE} = 9,\\ t^1_{BREAKAWAY} = t^2_{BREAKAWAY} = 3$;
	\item $u^1_{HAST} < u^1_{LATE} < u^1_{GOOD},\\ u^2_{HAST} < u^2_{LATE} < u^2_{GOOD}$.
\end{itemize}

Поскольку первый игрок, заведомо не имеющий проблем с недозагрузкой, укладывается в дедлайн только если компьютер будет в его безраздельном пользовании, то очевидно, что сочетания стратегий, в которых оба игрока выбирают одну машину, равновесиями Нэша быть не могут. Похожим образом, поскольку второму игроку выгоднее опоздать с расчётами нежели быть наказанным за использование отдельного компьютера для недостаточно большой задачи ($u^a_{HAST} < u^a_{LATE}$), равновесиями Нэша не могут быть и те ситуации, где каждый из сотрудников использует свою машину. По структуре выплат игра оказывается неотличима от игры в чёт-нечет, вовсе не имеющей решений в чистых стратегиях и с единственным равновесием Нэша в точке независимого равновероятного выбора между альтернативами обоими игроками. При этом максимально загруженная машина с вероятностью $\frac{1}{2}$ проработает либо $10$ часов при совпадении их выбора, либо $8$ при несовпадении, что даёт математическое ожидание цены стабильности, равное $\frac{9}{8}$. По сути, при отказе от монотонности функции оплаты за срочность вряд ли имеет смысл вообще рассуждать о ценах стабильности и анархии, поскольку класс игр теперь начинает включать и такие, которые явно не имеют отношения к поискам минимума продолжительности вычислений.

\section{Смешанные равновесия игры $\Gamma^3_n$}\label{sec:ch2/sec3}

Планирование заданий с немонотонной отдачей представляет собой довольно обширный класс конфликтов, анализ которого в общем виде выходит за рамки данной работы. Для наглядной демонстрации необходимого эффекта вполне достаточно будет специально сконструированного примера.  Рассмотрим игру $\Gamma^3_n$ той же общей схемы, что в предыдущем разделе, но чуть боле сложную, с $3$ однотипными заданиями и $n \ge 2$ одинаковыми компьютерами:
\begin{itemize}
	\item $t_i^1 = t_i^2 = t_i^3 = 2, i = \overline{1,n}$;
	\item $t^1_{DEADLINE} = t^2_{DEADLINE} = t^3_{DEADLINE} = 5,\\ t^1_{BREAKAWAY} = t^2_{BREAKAWAY} = t^3_{BREAKAWAY} = 3$;
	\item $u^1_{HAST} = u^2_{HAST} = u^3_{HAST} = 0,\\ u^1_{GOOD} = u^2_{GOOD} = u^3_{GOOD} = 3,\\ u^1_{LATE} = u^2_{LATE} = u^3_{LATE} = 2$
	%	\item $u^1_{HAST} < u^1_{LATE} < u^1_{GOOD},\\ u^2_{HAST} < u^2_{LATE} < u^2_{GOOD},\\ u^3_{HAST} < u^3_{LATE} < u^3_{GOOD}$.
\end{itemize}

Проще говоря, в игре $\Gamma^3_n$ выгоднее всего использовать компьютер вдвоём "--- $4$ часа суммарной продолжительности работы оказываются как раз в оптимальном промежутке между границей недогруза и дедлайном. Следующий по выгоде вариант "--- использование одной машины втроём, что приводит к штрафу за опоздание. Наименее привлекателен выбор компьютера, на котором задача оказывается одна "--- за такое расходование общественного ресурса игрок не получает вообще ничего. На первый взгляд эта игра не выглядит слишком необычно. В ней без особого труда находятся равновесия Нэша в чистых стратегиях "--- все наборы $(i, i, i), i = \overline{1,n}$, причём очевидно и то, что других решений в чистых стратегиях быть не может. Со смешанными стратегиями дело становится чуть интереснее.

\begin{lemma}
	Пусть $S_* \subseteq \{1, \ldots, n\}$ "--- произвольное непустое подмножество компьютеров. Тогда в игре $\Gamma^3_n$ набор одинаковых смешанных стратегий $s^1 = s^2 = s^3 = \left(\frac{[1 \in S_*]}{\left| S_* \right|}, \ldots, \frac{[n \in S_*]}{\left| S_* \right|}\right)$, где каждый игрок независимо и равновероятно выбирает одну из машин множества $S_*$, является равновесием по Нэшу.\footnote{Здесь и далее для упрощения и сокращения записи используется нотация <<скобка Айверсона>>: $[\textsc{true}] = 1, [\textsc{false}] = 0$.}
\end{lemma}

\begin{proof}
	Воспользуемся тем, что достаточно проверить отклонения только в пользу чистых стратегий. Обозначим символом $\left. s \right|^a_i$ отклонение от набора $s$ игроком $a$ в пользу стратегии $i$ и заметим, что 
	\begin{align*}
		u^a(\left. s \right|^a_i) &= [i \in S_*] \left(\frac{(\left| S_* \right| - 1)^2}{\left| S_* \right|^2} u^a_{HAST} + 2 \frac{\left| S_* \right| - 1}{\left| S_* \right|^2} u^a_{GOOD} + \frac{1}{\left| S_* \right|^2} u^a_{LATE}\right)\\
		&= [i \in S_*] \frac{6 \left| S_* \right| - 4}{\left| S_* \right|^2}.
	\end{align*}
	
	Таким образом, для каждого игрока максимум ожидаемого выигрыша достигается при отклонении в пользу любого из компьютеров, входящих в $S_*$.
\end{proof}

Соответственно без отклонений математическое ожидание выигрышей предстаёт в виде $u^a(s) = \frac{6 \left| S_* \right| - 4}{\left| S_* \right|^2}, a = \overline{1,3}$. Для доказательства того, что других равновесий по Нэшу в смешанных стратегиях нет, нам понадобятся ещё пара лемм:

\begin{lemma}
	В игре $\Gamma^3_n$ набор смешанных стратегий $s = (s^1, s^2, s^3)$ может быть равновесием по Нэшу только в том случае, когда $s^1 = s^2 = s^3$.
\end{lemma}

\begin{proof}
	Пусть $s^a = (p_1^a, \ldots, p_n^a), a = \overline{1,3}$. Если стратегии игроков не совпадают, то найдётся компьютер $i$, для которого (без потери общности) вероятности выбора первым и вторым игроком $p_i^1 > p_i^2$. В силу элементарных свойств вероятностей, найдётся также и компьютер $j$, где $p_j^1 < p_j^2$. Выпишем математическое ожидание выигрышей тех же игроков при выборе ими $i$"~го компьютера (для $j$"~го всё совпадает с точностью до индекса, очевидно):
	\begin{align*}
		u^1(\left. s \right|^1_i) &= (1 - p_i^2) (1 - p_i^3) u^1_{HAST} + (p_i^2 + p_i^3 - 2 p_i^2 p_i^3) u^1_{GOOD} + p_i^2 p_i^3 u^1_{LATE}\\ &= 3 p_i^2 + 3 p_i^3 - 4 p_i^2 p_i^3;\\
		u^2(\left. s \right|^2_i) &= 3 p_i^1 + 3 p_i^3 - 4 p_i^1 p_i^3;\\
		u^3(\left. s \right|^3_i) &= 3 p_i^1 + 3 p_i^2 - 4 p_i^1 p_i^2.
	\end{align*}
	
	Про функцию $f(x,y) = 3x + 3y - 4xy$ на области определения $0 \le x \le 1, 0 \le y \le 1$ можно заметить следующее "--- если $f(x_0, y_0) < 2$, то для любых $x_1 > x_0, y_1 \ge y_0$ или $x_1 \ge x_0, y_1 > y_0$ выполняется $f(x_0, y_0) < f(x_1, y_1)$. Предположим, что $p_i^3 \le p_j^3$ (если $p_i^3 \ge p_j^3$, рассуждения аналогичны с точностью до индексов). Рассмотрим следующие случаи:
	\begin{itemize}
		\item $p_i^2 < p_j^2$. В силу элементарных свойств вероятностей мы можем быть уверены, что $p_i^2 < \frac{1}{2}$ и $p_i^3 \le \frac{1}{2}$, откуда $u^1(\left. s \right|^1_i) < 2$, а значит $u^1(\left. s \right|^1_i) < u^1(\left. s \right|^1_j)$. Поскольку $p_i^1 > p_i^2 \ge 0$, первый игрок выбирает неоптимальную стратегию с ненулевой вероятностью. $\bot$
		\item $p_i^2 \ge p_j^2$, из чего следует $p_i^1 > p_j^1$, а значит аналогично предыдущему пункту $u^3(\left. s \right|^3_j) < u^3(\left. s \right|^3_i)$. Если $p_j^3 > 0$, то третий игрок выбирает неоптимальную стратегию с ненулевой вероятностью. $\bot$
		\item $p_i^2 \ge p_j^2$, как и в предыдущем случае, но теперь $p_i^3 = p_j^3 = 0$. Из $p_j^1 < p_i^1$ аналогичным образом следует $u^2(\left. s \right|^2_j) < u^2(\left. s \right|^2_i)$, и поэтому $p_j^2 > p_j^1 \ge 0$ влечёт выбор неоптимальной стратегии с ненулевой вероятностью уже вторым игроком. $\bot$
	\end{itemize}
	
	Таким образом, предположение о существовании компьютера, для которого вероятность выбора его одним игроком отличается от вероятности выбора другим, в любом случае противоречит необходимому условию равновесия Нэша.
\end{proof}

\begin{lemma}
	В игре $\Gamma^3_n$ набор одинаковых смешанных стратегий может быть равновесием по Нэшу только в том случае, когда все компьютеры, выбираемые с ненулевой вероятностью, выбираются с равными вероятностями.
\end{lemma}

\begin{proof}
	Возьмём любой набор, состоящий из одинаковых стратегий $(p_1, \ldots, p_n)$, где $0 < p_i < p_j$. Пользуясь формулой выплат из предыдущей леммы, $u^a(\left. s \right|^a_i) = 2 p_i (3 - 2 p_i)$. Опять же, из $p_i < \frac{1}{2}$ следует $u^a(\left. s \right|^a_i) < u^a(\left. s \right|^a_j)$, а значит все игроки выбрали неоптимальную стратегию с ненулевой вероятностью, что противоречит необходимому условию равновесия по Нэшу.
\end{proof}

Доказав, что предложенные точки равновесия исчерпывают пространство решений в смешанных стратегиях, мы можем сконструировать выпуклую оболочку множества достижимых векторов выплат. Поскольку множество целиком лежит на прямой $(u, u, u)$, достаточно найти минимум и максимум ожидаемых выигрышей:

\begin{align*}
	\min_{\emptyset \subset S_* \subseteq \{1,\ldots,n\}} \frac{6 \left| S_* \right| - 4}{\left| S_* \right|^2} &= \frac{6 n - 4}{n^2};\\
	\max_{\emptyset \subset S_* \subseteq \{1,\ldots,n\}} \frac{6 \left| S_* \right| - 4}{\left| S_* \right|^2} &= 2.
\end{align*}

Таким образом искомая выпуклая оболочка представляет собой отрезок, соединяющий точки $(2, 2, 2)$ и $(\frac{6 n - 4}{n^2}, \frac{6 n - 4}{n^2}, \frac{6 n - 4}{n^2})$. Если бы рассматриваемая игра не была чувствительна к дополнительной информационной асимметрии, на этом её анализ можно было бы закончить "--- все игроки находятся в равном положении и, действуя оптимально, могут ожидать равных выигрышей из указанного промежутка. Однако, взгляд на этот конфликт через призму модели заговоров делает картину происходящего куда интереснее.

\section{Коррелированные равновесия игры $\Gamma^3_n$ в пространстве заговоров}\label{sec:ch2/sec4}

Проанализируем этот же конфликт с позиций теории заговоров, перейдя к игре $\Gamma^3_n | \{\{1,2\}\}$. При этом игра $\Gamma^3_n$ дополняется одной вещественной рулеткой, результат вращения которой перед выбором стратегии узнают игроки $1$ и $2$, но не $3$. Для получения равновесного набора коррелированных стратегий, выводящего платежи за выпуклую оболочку множества решений в смешанных стратегиях, заговорщикам достаточно взять любую из классических точек равновесия по Нэшу с $\left| S_* \right| \ge 2$, но вместо того, чтобы выбирать между элементами $S_* \subseteq \{1, \ldots, n\}$ независимо, они должны разделить тайную рулетку на $\left| S_* \right|$ равных секторов и делать свой выбор синхронно, в зависимости от выпавшего сектора. Опишем это более формально, используя пространство корреляции
\begin{equation*}
	\Phi = \langle A, \Omega, \mathfrak{I}^a, \mathbb{P}, a \in A \rangle
\end{equation*}

В данном случае множество состояний природы $\Omega = \left[0, 1\right)$, $\sigma$"~алгебры информированности игроков $\mathfrak{I}^1 = \mathfrak{I}^2$ "--- борелевские, $\mathfrak{I}^3 = \{\emptyset, \Omega\}$ и мера $\mathbb{P}(X) = \left| X \right|$. Вышеописанные стратегии в игре $\Gamma^3_n | \Phi$ можно представить в виде функций, отображающих множество состояний природы в пространство смешанных стратегий:
\begin{align*}
	\mathbf{s}^1(\omega) = \mathbf{s}^2(\omega) &= ([\zeta(\omega) = 1], \ldots, [\zeta(\omega) = n]); \\ \mathbf{s}^3(\omega) &= \left(\frac{[1 \in S_*]}{\left| S_* \right|}, \ldots, \frac{[n \in S_*]}{\left| S_* \right|}\right),
\end{align*}
где общая для игроков 1 и 2 функция $\zeta : \Omega \rightarrow S_*$ определяет разбиение рулетки на $\left| S_* \right|$ равных секторов.

Выплаты в этом наборе уже не симметричны:
\begin{align*}
	u^1(\mathbf{s}) = u^2(\mathbf{s}) &= \frac{\left| S_* \right| - 1}{\left| S_* \right|} u^a_{GOOD} + \frac{1}{\left| S_* \right|} u^a_{LATE} = \frac{3 \left| S_* \right| - 1}{\left| S_* \right|};\\
	u^3(\mathbf{s}) &= \frac{\left| S_* \right| - 1}{\left| S_* \right|} u^a_{HAST} + \frac{1}{\left| S_* \right|} u^a_{LATE} = \frac{2}{\left| S_* \right|}.
\end{align*}

При этом ситуация действительно является равновесием по Нэшу, поскольку первый и второй игроки могут в качестве стратегий использовать любые функции, отображающие $\Omega$ в пространство вероятностных мер на $\{1, \ldots, n\}$, а вот третий игрок вынужден довольствоваться только константными в виду того, что механизм корреляции не информирует его о состоянии природы. Заметив, что $\frac{3 \left| S_* \right| - 1}{\left| S_* \right|} > 2$ при $\left| S_* \right| \ge 2$, мы подтверждаем чувствительность игры $\Gamma^3_n$ к дополнительной информационной асимметрии "--- в новых решениях игроки, наблюдающие не связанный напрямую с выплатами случайный эксперимент, увеличивают свой выигрыш по сравнению с наилучшим результатом, достижимым в классическом смешанном случае.

\section{Коллективная рациональность решений}\label{sec:ch2/sec5}

При рассмотрении множества решений игры $\Gamma^3_n$ в обычных смешанных стратегиях следует обратить внимание на то, что при $\left| S_* \right| > 2$ получающиеся точки равновесия по Нэшу лишены оптимальности не только в смысле Парето, но и по Слейтеру. В самом деле, выплаты всем игрокам в точках, где $\left| S_* \right| = 1$ или $\left| S_* \right| = 2$, равны $2$, а вот при больших размерах множества $\frac{6 \left| S_* \right| - 4}{\left| S_* \right|^2} < 2$. Такие обстоятельствах наделяют решения с использованием одного или двух компьютеров особым статусом "--- можно ожидать, что агенты, знакомые с принципом коллективной рациональности, всё же сумеют договориться о том, чтобы не оказаться в ситуации, которую другое решение доминирует по всем платежам. Естественно сразу задаться вопросом, а нельзя ли и решения с учётом дополнительной информационной асимметрии отфильтровать похожим образом, выделив из них удовлетворяющие принципам коллективной рациональности хоть в каком-то смысле?

При добавлении в игру пространства заговоров, состоящего из группы $\{1, 2\}$, сразу бросается в глаза то, что классические принципы коллективной рациональности становится бесполезны. В новых точках равновесия выигрыши первых двух игроков теперь $u^1(\mathbf{s}) = u^2(\mathbf{s}) > 2$ и растут с ростом $k$, а третьего "--- $u^3(\mathbf{s}) < 2$ и падают, что говорит о прямом антагонизме интересов. Это делает невозможным коллективно рациональный выбор в обычном смысле между смешанным равновесием с $\left| S_* \right|$ равным $1$ или $2$ и коррелированными решениями с различными $k$. Тем не менее, можно попытаться применить более тонкий критерий оптимальности, опирающийся на слегка расширенную интерпретацию происходящего в игре "--- назовём этот формализм \emph{структурно согласованным равновесием по Нэшу}.

Идея структурной согласованности равновесий в играх с заговорами довольно проста "--- если допустить, что входящие в семейство заговоров группы игроков объединяет не только общий механизм корреляции, но и в целом большие возможности для согласования действий, то среди обычных равновесий по Нэшу в коррелированных стратегиях можно особо выделить обладающие устойчивостью не только к индивидуальным отклонениям, но и к групповым, имея в виду исключительно входящие в семейство заговоров группы. Для игры $\Gamma^3_n | \{\{1,2\}\}$, например, это могло бы означать, что структурно несогласованными окажутся те равновесия, для которых найдётся отклонение, в котором участвуют первый и второй игроки, обоюдно увеличивая при этом свои выигрыши. В наиболее общих терминах это можно выразить так:
\begin{definition}
	В игре с заговорами $\Gamma | \mathfrak{A}$ равновесие по Нэшу $\mathbf{s}$ называется структурно согласованным, если для всех заговоров $A_* \in \mathfrak{A}$ отсутствуют приемлемые отклонения от ситуации $\mathbf{s}$.
\end{definition}

Остаётся сформулировать, что в рамках данной модели можно считать отклонением, приемлемым для того или иного заговора. Если взглянуть на этот вопрос как на проблему многокритериальной оптимизации, где критериями являются выигрыши отдельных участников, то напрашиваются два варианта:
\begin{enumerate}
	\item Отклонение приемлемо, если оно увеличивает выигрыш всех участников заговора; (по Слейтеру)
	\item Отклонение приемлемо, если оно увеличивает выигрыш хотя бы одного участника заговора, при этом не уменьшая выигрыша остальных. (по Парето)
\end{enumerate}

Увы, оба варианта нельзя назвать подходящими для наших целей. Разумно было бы ожидать, что добавление в любой заговор <<болвана>>, т.е. игрока с единственной чистой стратегией и константным выигрышем, не должно ничего менять в решении игры. Однако, заговор с подобным <<болваном>> в составе вообще не может иметь приемлемых отклонений в смысле Слейтера, а значит его участники теряют возможность использовать коллективную рациональность. Это делает первый из предложенных вариантов, очевидно, слишком слабым. С другой же стороны, рассмотрение игры в трёхсторонний чёт"~нечет (см. \ref{sec:ch1/sec4}) с двумя заговорами порождает неприятную проблему и для второго варианта. Если 1"~й игрок находится в заговоре со 2"~м, а 2"~й с 3"~м, то можно ожидать исходов с платежами, образованными любым смешением $(5,5,2)$ и $(2,5,5)$, причём выбор пропорции происходит по воле 2"~го игрока. Загвоздка в том, что приемлемость в смысле Парето побуждает 2"~го игрока в ситуации $(5,5,2)$ отклоняться в рамках заговора с 3"~м для улучшения чужого выигрыша. Аналогично, в ситуации $(2,5,5)$ 2"~му игроку приходится спасать 1"~го. В промежуточных же ситуациях 2"~й игрок может помочь обоим, а значит подобный альтруизм вообще исключает существование структурно согласованного равновесия в задаче, делая второй вариант определения приемлемости слишком сильным. Для обхода обоих проблем следует предложить промежуточное определение отклонения, которое будет сильнее чем по Слейтеру, но слабее чем по Парето:
\begin{definition}
	В игре с заговорами $\Gamma | \mathfrak{A}$ ситуацию $\mathbf{s}_* \neq \mathbf{s}$ назовём отклонением от $\mathbf{s}$, приемлемым для заговора $A_* \in \mathfrak{A}$, если
	\begin{itemize}
		\item $\forall a \notin A_* \quad \mathbf{s}^a = \mathbf{s}_*^a$;
		\item $\forall a \in A_* \quad u^a(\mathbf{s}_*) \geq u^a(\mathbf{s})$;
		\item $\forall a : \mathbf{s}^a \neq \mathbf{s}_*^a \quad u^a(\mathbf{s}_*) > u^a(\mathbf{s})$.
	\end{itemize}
\end{definition}

Если бы мы говорили только о равновесиях Нэша в чистых и смешанных стратегиях, такой критерий оптимальности оказался бы слишком сильным "--- действительно, одновременный выбор первым и вторым игроками компьютера $i$, выбираемого третьим с вероятностью $p_i^3 < 1$, даёт обоим выигрыш $(1 - p_i^3) u^a_{GOOD} + p_i^3 u^a_{LATE} = 3 - p_i^3 > 2$, что отсеивает вообще все решения в игре $\Gamma^3_n$. Для игры с заговорами же можно сформулировать следующее:

\begin{theorem}
	В игре с заговорами $\Gamma^3_n | \mathfrak{A}$ структурно согласованное равновесие Нэша существует при любых $n$ и $\mathfrak{A}$. При невырожденном $\mathfrak{A}$, каждому заговору из двух участников соответствует единственное такое равновесие.
\end{theorem}

\begin{proof}
	Возможны три случая в зависимости от $\mathfrak{A}$:
	\begin{enumerate}
		\item $\mathfrak{A} = \emptyset$. При пустом семействе заговоров групповые отклонения невозможны, так что любое равновесие по Нэшу в чистых или смешанных стратегиях будет структурно согласованным.
		\item $\mathfrak{A} = \{\{1, 2, 3\}\}$. При вырожденном семействе заговоров с одним публичным механизмом корреляции множество решений представляет собой выпуклую оболочку векторов платежей смешанных равновесий, что, как было показано в разделе \ref{sec:ch2/sec3}, даёт отрезок, соединяющий $(2, 2, 2)$ и $(\frac{6 n - 4}{n^2}, \frac{6 n - 4}{n^2}, \frac{6 n - 4}{n^2})$. При этом критерий структурной согласованности, очевидно, отсеивает всё, кроме точки $(2, 2, 2)$, соответствующей чистым равновесиям с выбором одной любой общей стратегии на всех и смешанным равновесиям с независимым равновероятным выбором из двух любых одних и тех же стратегий каждым игроком.
		\item $\mathfrak{A}$ содержит $\{1, 2\}$, $\{1, 3\}$ или $\{2, 3\}$. Как было показано в разделе \ref{sec:ch2/sec4}, использование тайных механизмов корреляции даёт новые точки равновесия с выплатами $\frac{3k - 1}{k}$ участникам заговора и $\frac{2}{k}$ аутсайдеру, где $k$ "--- количество задействованных в наборе стратегий машин. Поскольку при $k \ge 2$ выигрыш каждого из заговорщиков превосходит наилучший классический результат, никакие точки смешанного равновесия структурно согласованными уже не будут. Поскольку с ростом $k$ растут и выигрыши заговорщиков, отсеиваются также и все коррелированные равновесия, кроме максимизирующих доход любой из пар заговорщиков при $k = n$. Таким образом, каждой группе, входящей в структуру заговоров, соответствует одна (с точностью до пермутаций рулеток) точка структурно согласованного равновесия, в которой игроки выбирают равновероятно из всего доступного парка машин, причём выбор членов группы всегда совпадает между собой и независим с выбором оставшегося игрока.
	\end{enumerate}
\end{proof}

Для новых структурно согласованных равновесий несложно подобрать вполне естественную интерпретацию. Если представить, что два сотрудника могут координировать свои действия втайне от третьего, то нет ничего неожиданного в их стремлении выбрать один компьютер на двоих, чтобы избежать штрафа за недозагруз, минимизируя при этом шанс для третьего игрока наткнуться на них по воле случая, лишив их премии за срочность. Достигается это логичным образом тогда, когда наибольшая вероятность выбора каждой из машин минимальна, т.е. при равновероятном выборе из всех. Аналогичным образом, для третьего игрока целью становится максимизация наименьшей вероятности выбора каждого из компьютеров, поскольку он понимает, что заговорщики действуют заодно и пытаются избежать встречи с ним, и это тоже достигается в ситуации равновероятного выбора из всего парка машин.

\section{Сохранение тайн заговоров в процессе выработки консенсуса}\label{sec:ch2/sec6}

Как известно, в матричных играх нередко наличествуют несколько точек равновесия по Нэшу (например, координационные игры), причём наборы, состоящие из стратегий, принадлежащих к разным точкам, сами равновесиями не являются. При интерпретации таких точек равновесия в качестве решений игры приходится оговаривать, что выбор игроками равновесных стратегий в сущности не является независимым "--- предпочтение одной точки равновесия другой должно носить характер консенсуса среди игроков. В классических играх с полной информацией это не создаёт больших проблем, поскольку процедура, приводящая к консенсусу, вполне может быть гласной. Однако, когда речь заходит о моделировании конфликтов с заговорами, данный вопрос начинает требовать гораздо более осторожного подхода. При появлении в игре информационной асимметрии, существенно влияющей на её исход, игрокам может становиться выгодно изменять картину этой асимметрии (оповещая, к примеру, о значении тайного сигнала игроков, которые не должны его знать в соответствии со структурой пространства корреляции), для чего при неправильном дизайне могут использоваться те самые механизмы выработки консенсуса. Для того чтобы понять, что может пойти не так, имеет смысл начать с классического случая. Если представить произвольную матричную игру в виде реального процесса с живыми игроками под управлением беспристрастного ведущего, следящего за соблюдением протокола игры, то для получения равновесия по Нэшу можно применить что-то вроде следующей процедуры:

\begin{enumerate}
	\item Ведущий объявляет матрицу выплат;
	\item Игроки гласно обсуждают выбор стратегий;
	\item Игроки втайне друг от друга извещают ведущего о своих ходах;
	\item Ведущий оглашает собранный набор стратегий;
	\item Ведущий может с ненулевой вероятностью выбрать любого из игроков и предложить ему переходить;
	\item Ведущий вычисляет и объявляет выигрыши.
\end{enumerate}

Этой процедуры вполне достаточно, если мы говорим о классических равновесиях Нэша в чистых и смешанных стратегиях. Во втором случае следует только уточнить, что реализация конкретного исхода, определяемого набором смешанных стратегий, либо не происходит вообще (ведущий объявляет математические ожидания выигрышей), либо происходит только на последнем этапе. Однако, при добавлении к модели пространств корреляции ситуация несколько усложняется. Поскольку коррелированными стратегиями могут быть любые функции, отображающие получаемые игроками сигналы в смешанные стратегии, кажется естественным представлять себе, как игроки сами вычисляют избранные ими же функции, получив все релевантные сигналы:

\begin{enumerate}
	\item Ведущий объявляет матрицу выплат;
	\item Игроки гласно обсуждают выбор коррелированных стратегий;
	\item Ведущий генерирует состояние природы и извещает игроков о соответствующих событиях из их $\sigma$"~алгебр информированности;
	\item Игроки втайне друг от друга вычисляют смешанные стратегии и извещают ведущего о своих ходах;
	\item Ведущий оглашает собранный набор смешанных стратегий;
	\item Ведущий может с ненулевой вероятностью выбрать любого из игроков и предложить ему переходить;
	\item Ведущий вычисляет и объявляет выигрыши.
\end{enumerate}

К сожалению, такой наивный подход обладает существенным недостатком "--- он работает ожидаемым образом только в симметричных пространствах корреляции, где $\sigma$"~алгебры информированности всех игроков совпадают. Если же мы говорим о пространствах заговоров, то возникают сразу две проблемы. Во-первых, между этапами 3 и 4 у кого-то из игроков может возникнуть искушение разгласить значение приватного сигнала, если это может сподвигнуть неосведомлённых о нём игроков на выбор более выгодной для разглашающего стратегии. Этот вопрос ещё можно было бы закрыть, добавив в алгоритм запрет на коммуникацию между игроками, начинающийся после этапа 2, но, увы, это только одна из проблем.

Во-вторых, что поправить несколько сложнее, при наличии информационной асимметрии возможность одного из игроков изменить выбранную стратегию на этапе 6 перестаёт соответствовать концепции равновесия по Нэшу. В симметричном случае между первоначальным выбором смешанной стратегии на этапе 4 и возможным отклонением от неё на этапе 6 игрок не получает никакой дополнительной информации, поскольку публичность сигнала и так позволяет вычислить избранные другими игроками смешанные стратегии "--- оглашая их, ведущий, фактически, только фиксирует результат публичной договорённости, достигнутой на этапе 2. Асимметричные пространства корреляции же содержат события, о которых на этапе 3 оповещается только часть игроков. При этом каждый игрок может достоверно вычислить свою смешанную стратегию, но не стратегии оппонентов, завязанные на скрытые от него сигналы. В этой ситуации оглашение ведущим собранных смешанных стратегий на этапе 5 увеличивает знание игроков перед принятием кем-то из них решения об отклонении, что противоречит идее равновесия по Нэшу. Для приведения вышеописанной процедуры в соответствие с моделируемым формализмом её необходимо изменить несколько противоречащим интуиции образом:

\begin{enumerate}
	\item Ведущий объявляет матрицу выплат;
	\item Игроки гласно обсуждают выбор коррелированных стратегий;
	\item Игроки втайне друг от друга извещают ведущего о выбранных коррелированных стратегиях;
	\item Ведущий оглашает собранный набор коррелированных стратегий;
	\item Ведущий генерирует состояние природы и вычисляет смешанные стратегии игроков;
	\item Ведущий может с ненулевой вероятностью выбрать любого из игроков, известить его обо всех реализовавшихся событиях из его $\sigma$"~алгебры информированности и предложить ему изменить вычисленную ведущим смешанную стратегию;
	\item Ведущий вычисляет и объявляет выигрыши.
\end{enumerate}

Таким образом, в отличие от симметричного случая модель заговоров не позволяет смотреть на коррелированные стратегии как на <<чёрные ящики>> в головах игроков, просто подсказывающие им синхронную реакцию на раздражители. Здесь наборы коррелированных стратегий приходится интерпретировать как проговариваемые и формально фиксируемые соглашения, поскольку концепция равновесия по Нэшу подразумевает возможность отклонения именно в тот момент, когда общим знанием являются только намерения игроков реагировать тем или иным образом на тайные сигналы, но ещё не конкретные их реакции.

Интересно то, что при попытке обобщить эту процедуру для получения структурно согласованных равновесий в пространствах заговоров мы снова сталкиваемся с похожей проблемой. На первый взгляд достаточно было бы уточнить только этап 6, чтобы ведущий мог предлагать отклониться от выбранных стратегий как отдельным игрокам, так и целым группам заговорщиков. Однако, это срабатывает ожидаемым образом только для наиболее простых семейств с непересекающимися заговорами. В том же случае, когда у двух заговоров могут быть общие участники, уже групповые отклонения перестают соответствовать формализму "--- ведь если перед их обсуждением ведущий оповестит каждого заговорщика о рулетках всех заговоров, в которые тот входит, то кто-то из них тогда мог бы разгласить тайну другого заговора, тем самым неправомерно увеличив знания не входящих в него игроков до принятия решения об отклонении. Проще всего поправить это с помощью выноса групповых отклонений в отдельный этап, предшествующий генерации состояния природы:

\begin{enumerate}
	\item Ведущий объявляет матрицу выплат;
	\item Игроки гласно обсуждают выбор коррелированных стратегий;
	\item Игроки втайне друг от друга извещают ведущего о выбранных коррелированных стратегиях;
	\item Ведущий оглашает собранный набор коррелированных стратегий;
	\item Ведущий может с ненулевой вероятностью выбрать любую из групп заговорщиков и предложить им изменить выбранные коррелированные стратегии;
	\item Ведущий генерирует состояние природы и вычисляет смешанные стратегии игроков;
	\item Ведущий может с ненулевой вероятностью выбрать любого из игроков, известить его обо всех реализовавшихся событиях из его $\sigma$"~алгебры информированности и предложить ему изменить вычисленную ведущим смешанную стратегию;
	\item Ведущий вычисляет и объявляет выигрыши.
\end{enumerate}

Заметим, что этап 5 (обсуждение и утверждение заговорщиками группового отклонения) "--- сам по себе многостадийный процесс, в котором те, кто хочет увеличить свой выигрыш сменой стратегии, предлагают проект отклонения, а остальные участники заговора имеют индивидуальное право вето, если этот проект приносит им убыток. Таким образом, структурная согласованность равновесия подразумевает, что группы заговорщиков могут отклоняться на стадии планирования, перед получением игроками информации о состоянии природы, а индивидуальные отклонения возможны уже после срабатывания механизмов корреляции, но до оглашения конкретный смешанных стратегий, сыгранных противниками.

\section{Немонотонная отдача в других конфликтах планирования}\label{sec:ch2/sec7}

При помощи игры $\Gamma^3_n$ мы продемонстрировали, во-первых, что планирование заданий с немонотонной функцией оплаты за срочность может быть чувствительно к дополнительной информационной асимметрии, и, во-вторых, что в играх с заговорами, несмотря на присущий им частичный антагонизм интересов, возможны решения, отвечающие принципу коллективной рациональности. Подчеркнём полезность этих результатов, заметив, что проблема планирования заданий гораздо шире рассмотренного нами примера, причём как в смысле возможных значений параметров (матрицы коэффициентов $(t_i^a) \in \mathbb{R}_{\ge 0}^{m \times n}$ и функций отдачи $v^a : \mathbb{R}_{\ge 0} \rightarrow \mathbb{R}, a = \overline{1,m}$), так и в смысле разнообразия практических приложений модели. В контексте этой работы нет смысла заниматься слишком подробным разбором более сложных случаев, но для того, чтобы показать возможную связь модели с реальным миром за пределами вычислительных центров со странными схемами поощрений сотрудников, попробуем построить пару примеров с более солидной предметной областью.

Начнём с экономики, представив себе, как $m$ компаний готовятся выйти на рынок с предложениями высокотехнологичного товара, и перед ними встаёт выбор между $n$ различными открытыми стандартами на один и тот же его важный аспект. К примеру, это могут быть разнообразные промышленные роботы и стандарты их интеграции в <<умный>> цех. Когда компания $a \in \{1, \ldots, m\}$ выходит на рынок стандарта $i \in \{1, \ldots, n\}$, она тем самым осуществляет вклад в его развитие, характеризующийся векторной константой $t_i^a \in \mathbb{R}_{\ge 0} \times \ldots \times \mathbb{R}_{\ge 0}$, компоненты которой соответствуют отдельным независимым аспектам (например, функциям, для исполнения которых приобретаются роботы). Если в ситуации $s$ одним стандартом $i$ пользуются несколько компаний, то простым суммированием их вкладов можно посчитать общий индекс развития $t_i(s) = [s^1 = i] t_i^1 + \ldots + [s^m = i] t_i^m$. На ожидаемый доход от инвестиций в каждый из стандартов существенным образом влияют два дисконтирующих фактора: сетевой эффект и насыщение рынка.

Под сетевым эффектом мы понимаем зависимость покупательского энтузиазма от общего индекса развития стандарта "--- функция $0 \le \alpha^a(t) \le 1$ характеризует долю покупателей, готовых приобретать роботов компании $a$, выполненных по стандарту с общим индексом развития $t$. Более развитый стандарт всегда привлекает больше потребителей, так что функции $\alpha^a(t)$ монотонно неубывающие, т.е. $\alpha^a(t) \le \alpha^a(t + \Delta), \forall t, \Delta \succeq (0, \ldots, 0)$. Насыщение рынка, с другой стороны, подразумевает ограниченность спроса "--- при избытке инвестиций в любой из стандартов, платёжеспособности покупателей перестаёт хватать на всех, цены приходится снижать, а с ними падают и доходы. Соответственно, ещё одна функция $0 \le \beta^a(t) \le 1$ характеризует, какой долей прибыли придётся ограничиться компании $a$ для сохранения конкурентоспособности своих роботов на рынке стандарта с общим индексом развития $t$. Эта функция, по понятным причинам, монотонно невозрастающая, т.е. $\beta^a(t) \ge \beta^a(t + \Delta), \forall t, \Delta \succeq (0, \ldots, 0)$. Целью компании $a$ при выборе стратегии $s^a$ является максимизация комбинации дисконтирующих факторов $u^a(s) = \alpha^a(t_{s^a}(s)) \beta^a(t_{s^a}(s))$.

Можно представить и политическую интерпретацию этой же игры. Пусть в некий коллегиальный выборный орган пытаются избираться $n$ кандидатов (самостоятельно, без партийных списков), а $m$ эффективных менеджеров выбирают, за кого из них развернуть агитацию в подведомственных учреждениях. Когда олигарх $a$ принимает решение о поддержке кандидата $i$, тем самым он вносит вклад в его популярность, характеризующийся векторной константой $t_i^a \in \mathbb{R}_{\ge 0} \times \ldots \times \mathbb{R}_{\ge 0}$, компоненты которой соответствуют электорально значимым демографическим группам. Если в ситуации $s$ кандидата $i$ поддерживают несколько олигархов, то простым суммированием их вкладов можно получить общий индекс популярности кандидата $t_i(s) = [s^1 = i] t_i^1 + \ldots + [s^m = i] t_i^m$. На ожидаемую выгоду от поддержки того или иного кандидата влияют два дисконтирующих фактора: политическое влияние и готовность к сотрудничеству.

Политическое влияние кандидата в вопросах, интересующих поддержавшего его олигарха, очевидно, растёт вместе с общим индексом его популярности, что выражается функцией $0 \le \alpha^a(t) \le \alpha^a(t + \Delta) \le 1, \forall t, \Delta \succeq (0, \ldots, 0)$. Готовность же кандидата к сотрудничеству с каждым из своих сторонников, наоборот, падает с ростом его суммарной популярности, что выражается функцией $1 \ge \beta^a(t) \ge \beta^a(t + \Delta) \ge 0, \forall t, \Delta \succeq (0, \ldots, 0)$. Целью олигарха $a$ при выборе стратегии $s^a$ является максимизация комбинации дисконтирующих факторов $u^a(s) = \alpha^a(t_{s^a}(s)) \beta^a(t_{s^a}(s))$.

Сведём описание обоих конфликтов к матричной игре в нормальной форме:
\begin{equation*}
	\Gamma = \langle A, S^a, u^a(s), a \in A \rangle;
\end{equation*}
\begin{equation*}
	A = \{1, \ldots, m\}, S^1 = \ldots = S^m = \{1, \ldots, n\};
\end{equation*}
\begin{equation*}
	u^a(s) = \alpha^a(t_{s^a}(s)) \beta^a(t_{s^a}(s)), a = \overline{1,m};
\end{equation*}
\begin{equation*}
	\alpha^a, \beta^a : \mathbb{R}_{\ge 0} \times \ldots \times \mathbb{R}_{\ge 0} \rightarrow [0, 1];
\end{equation*}
\begin{equation*}
	0 \le \alpha^a(t) \le \alpha^a(t + \Delta) \le 1, \forall t, \Delta \succeq (0, \ldots, 0);
\end{equation*}
\begin{equation*}
	1 \ge \beta^a(t) \ge \beta^a(t + \Delta) \ge 0, \forall t, \Delta \succeq (0, \ldots, 0);
\end{equation*}
\begin{equation*}
	t_i(s) = [s^1 = i] t_i^1 + \ldots + [s^m = i] t_i^m, i = \overline{1,n}.
\end{equation*}

Несложно заметить сходство этой игры с проблемой планирования заданий, что мы здесь постарались подчеркнуть использованием тех же символов для переменных и констант. Фактически, единственным значимым отличием оказывается то, что в планировании заданий время $t$ было скаляром, а не вектором. Функциям оплаты за срочность в новых формулировках соответствуют $v^a(t) = \alpha^a(t) \beta^a(t)$, форма которых и определяет наши ожидания от исхода конфликта. Ранее мы уже говорили, что анализ проблемы планирования заданий в общем виде, для произвольных $(t_i^a)$ и $(v^a)$ представляет собой чрезвычайно сложную задачу, и, конечно же, переход от скаляров к векторам в домене функций отдачи дело нисколько не упрощает. Лучшее, что тут можно сделать "--- дать качественный прогноз для некоторых неформально описанных подклассов конфликта с опорой на здравый смысл, интуицию и аналогию с разобранным выше частным случаем $\Gamma^3_n$.

Во избежание переусложнения ограничимся случаем квазивогнутых функций отдачи $v^a(t) = \alpha^a(t) \beta^a(t)$, естественным образом обобщив это понятие на многомерные области определения. Для начала обозначим символом $T^a_{\nearrow}$ множество всех таких $t$, что $t_* \prec t \Rightarrow v^a(t_*) \le v^a(t) \wedge t_* \in T^a_{\nearrow}$. Аналогично, символом $T^a_{\searrow}$ обозначим множество всех таких $t$, что $t_* \succ t \Rightarrow v^a(t_*) \le v^a(t) \wedge t_* \in T^a_{\searrow}$. Это будут области непрерывного неубывания и невозрастания $v^a(t)$ соответственно. Если эти два множества покрывают всю область определения, т.е. $T^a_{\nearrow} \cup T^a_{\searrow} = \mathbb{R}_{\ge 0} \times \ldots \times \mathbb{R}_{\ge 0}$, то функция $v^a(t)$ квазивогнута. Для подобных функций также можно обозначить <<гребень>> $T^a_{\sim} = T^a_{\nearrow} \cap T^a_{\searrow}$, в одномерном случае соответствующий максимуму.

Попробуем представить логику конфликта для простейшего случая с двухэлементным семейством заговоров $\mathfrak{A} = \{A_p, A_q\}, A_p \cap A_q = \emptyset, A_p \cup A_q = A$. Сократим запись с помощью следующих обозначений:
\begin{equation*}
	t_i^{A_*} = \sum_{a \in A_*} t_i^a, \forall A_* \subseteq A;
\end{equation*}
\begin{equation*}
	\check{u}^a = \min_{1 \le i \le n} v^a(t_i^A), \hat{u}^a = \max_{1 \le i \le n} v^a(t_i^A), \forall a \in A.
\end{equation*}

Здесь $t_i^{A_*}$ соответствует суммарному индексу развития стандарта (популярности кандидата) $i$ при его выборе группой игроков $A_* \subseteq A$. Так же, для каждого игрока $a$ базовым платёжным интервалом называется промежуток от $\check{u}^a$ до $\hat{u}^a$, т.е. от наименьшего до наибольшего возможных выигрышей при единогласном выборе общей стратегии всеми участниками конфликта. Рассмотрим игры, ограниченные следующими условиями:

\begin{itemize}
	\item $\forall a \in A, i = \overline{1,n}, v^a(t_i^{A_p \cup \{a\}}) < \check{u}^a$, т.е. ни один игрок не может достигнуть нижней границы своего базового платёжного интервала, если ту же стратегию выбирает только группа $A_p$;
	\item $\forall a \in A, i = \overline{1,n}, v^a(t_i^{A_q \cup \{a\}}) > \check{u}^a$, т.е. каждый игрок преодолевает нижнюю границу своего базового платёжного интервала при выборе любой стратегии совместно с группой $A_q$;
	\item $\forall a \in A_q, i = \overline{1,n}, t_i^{A_q} \in T^a_{\searrow}$, т.е. для всех членов группы $A_q$ выбор общей стратегии выводит суммарный индекс в область невозрастания функции отдачи.
\end{itemize}

Таким образом мы очертили круг ситуаций, в которых участники конфликта поделены на две непересекающиеся группы заговорщиков. Если рассматривать происходящее с точки зрения каждого заговора в предположении, что аутсайдеры вообще не участвует в игре, несложно заметить, что любой стратегический набор, в котором заговорщики выбирают одну стратегию на всех, будет хорошим кандидатом в равновесия по Нэшу. Это не означает, что других равновесий не может быть, но мы в целях наглядности сознательно ограничимся анализом наборов, характеризующихся двумя независимыми распределениями вероятностей $p = (p_1, \ldots, p_n)$ и $q = (q_1, \ldots, q_n)$, где участники заговоров $A_p$ и $A_q$ синхронно внутри групп но независимо между группами выбирают стратегии $i = \overline{1,n}$ c вероятностями $p_i$ и $q_i$ соответственно, используя соответствующий приватный механизм корреляции. Выплаты при этом считаются по формуле
\begin{equation*}
	u^a(p, q) = \sum_{i=1}^n p_i ((1 - q_i) v^a(t_i^{A_p}) + q_i v^a(t_i^A)), \forall a \in A_p,
\end{equation*}
\begin{equation*}
	u^a(p, q) = \sum_{i=1}^n q_i ((1 - p_i) v^a(t_i^{A_q}) + p_i v^a(t_i^A)), \forall a \in A_q,
\end{equation*}

В соответствии с условиями, заговор $A_p$ недостаточно велик для максимизации прибылей своих участников, так что каждый из них предпочёл бы присоединиться к стратегии, избранной заговором $A_q$. Однако, поскольку тайна чужого заговора игрокам не доступна, от стратегии предписанной механизмом корреляции они могут отклониться только в пользу другой чистой стратегии. Таким образом, для получения прибыли в результате индивидуального отклонения от пары распределений $(p, q)$, заговорщику $a \in A_p$ необходимо и достаточно найти такие индексы $i \neq j \in \{1, \ldots, n\}$, что при $p_i > 0$ выполняется неравенство
\begin{equation*}
	(1 - q_i) v^a(t_i^{A_p}) + q_i v^a(t_i^A) < (1 - q_j) v^a(t_j^a) + q_j v^a(t_j^{A_q \cup \{a\}}).
\end{equation*}

Несложно заметить, что с ростом любого $q_j$ постепенно увеличивается и кол-во индексов $i$, для которых выполняется это неравенство, а при приближении $q_j$ к $1$ оно рано или поздно начинает выполняться для всех $i \neq j$. Зафиксировав произвольное распределение $q$, можно для каждого заговорщика $a \in A_p$ вычислить множество $S^a_q \subseteq \{1, \ldots, n\}$, состоящее из стратегий, допускающих подобные продуктивные отклонения. При этом, поскольку участники заговора $A_q$, отклонившись от предписанной стратегии, неизбежно терпят убытки, для них проверять ничего не нужно. В результате, произвольная пара распределений $(p, q)$ описывает равновесие по Нэшу тогда и только тогда, когда
\begin{equation*}
	\forall i \in \bigcup_{a \in A_p} S^a_q, p_i = 0.
\end{equation*}

Таким образом, если мы говорим о равновесиях только в классическом Нэшевском смысле без учёта коллективной рациональности, то в описанном противостоянии участникам большого заговора вообще не приходится думать о возможном предательстве со стороны соратников, тогда как малый заговор должен внимательно выбирать общую стратегию так, чтобы у его участников не было искушения попытаться угадать стратегию, избранную большим. Стремление же удостовериться в структурной согласованности обозначенных решений дают чуть более интересную картину.

Сразу оговоримся, что в установленных ограничениях сложно точно убедиться в структурной согласованности даже для узкого множества рассматриваемых $(p, q)$"~наборов, поскольку вполне возможны, например, коллективные отклонения, разделяющие заговор $A_q$ на две группы, выбирающие разные стратегии. Одна при этом получает прибыль в результате избавления от лишних участников (см. ограничение $t_i^{A_q} \in T^a_{\searrow}$, т.е. принадлежность точек единогласного выбора к области невозрастания функции отдачи). Вторая же потенциально увеличивает доход, присоединившись к стратегии, избранной заговором $A_p$, если существует достаточно большое $p_i$. Можно, конечно, попытаться наложить на параметры конфликта дополнительные ограничения, предупреждающие подобные и даже более экзотические отклонения, однако это сильно усложнит постановку, не слишком добавляя иллюстративности.

Вместо этого мы ограничимся поиском только тех $(p, q)$"~наборов, от которых не существует успешных коллективных отклонений в пользу других $(p, q)$"~наборов. Найденные точки равновесия всё ещё можно будет подозревать в отсутствии структурной согласованности, однако мы хотя бы исключим большой класс заведомо несогласованных. Итак, для прибыльности коллективного отклонения малого заговора достаточно найти такие индексы $i \neq j \in \{1, \ldots, n\}$, что при $p_i > 0$ для каждого $a \in A_p$ выполняется неравенство
\begin{equation*}
	(1 - q_i) v^a(t_i^{A_p}) + q_i v^a(t_i^A) < (1 - q_j) v^a(t_j^{A_p}) + q_j v^a(t_j^A).
\end{equation*}

Аналогично, для большого заговора отклонение успешно, если есть такие индексы $i \neq j \in \{1, \ldots, n\}$, что при $q_i > 0$ для каждого $a \in A_q$ выполняется неравенство
\begin{equation*}
	(1 - p_i) v^a(t_i^{A_q}) + p_i v^a(t_i^A) < (1 - p_j) v^a(t_j^{A_q}) + p_j v^a(t_j^A).
\end{equation*}

Рассматривая эти неравенства в свете ограничений, наложенных нами на параметры конфликта, несложно заметить, что коллективная рациональность стимулирует обе группы игроков к минимизации наибольших вероятностей выбора отдельных стратегий, однако по противоположным причинам. Участникам малого заговора выгодно сообща присоединиться к стратегии, выбранной большим заговором, выплаты членам которого такое совпадение стратегий заведомо уменьшает. Переводя на язык условленных ранее интерпретаций, слабый картель с удовольствием воспользовался бы развитостью стандарта (или влиятельностью кандидата) избираемого крупным картелем, однако крупному картелю, напротив, совершенно не улыбается делить ограниченный спрос на и без того насыщенном рынке (или соревноваться за внимание и без того уверенного в избрании политика) с лишними конкурентами.

На уровне коллективно рациональных решений игра, таким образом, превращается в разновидность двухсторонних пряток, где одна группа ищет встречи с другой, пытающейся этого столкновения избежать, а заговоры служат для получения преимущества от объединения усилий при минимизации вероятности присоединения к дележу прибылей нежелательных попутчиков. По сути, формализм структурно согласованного равновесия в играх с заговорами является не какой-то сложной экономической концепцией, а всего лишь воплощением интуитивного принципа, применявшегося, вероятно, ещё в дописьменную эпоху. Вполне возможно, что накой-нибудь охотник, заметив раненного мамонта при обходе племенных угодий, прикинул: <<В одиночку я его, пожалуй, не завалю, но и племя всё звать смысла нет. Шепну-ка я лучше на ушко паре друзей "--- это ж сколько почёта и славы будет, втроём столько мяса добыть.>> С похожего рассуждения и могла начаться мировая история заговоров.

\FloatBarrier
           % Глава 2
\chapter{Вычислительная сложность стратегий в повторяющихся играх с дисконтированием}\label{ch:ch3}

\section{Повторяющийся трёхсторонний чёт-нечет}\label{sec:ch3/sect1}

Чувствительность к дополнительной информационной асимметрии интересна не только тем, как на исходы проявляющих её игр влияет собственно информационная асимметрия. Далее будет показано, что в тех случаях, когда этим свойством обладают отдельные розыгрыши повторяющихся игр, необычные эффекты могут возникать даже при отсутствии приватных механизмов корреляции, в рамках классического смешанного расширения и совершенных подыгровых равновесий. Разберём этот феномен на примере игры в <<трёхсторонний чёт-нечет>>, описанной в разделе \ref{sec:ch1/sec4}. Сперва проанализируем повторяющуюся версию данного конфликта с классической точки зрения, применяя <<народную>> теорему.

\begin{table} [htbp]
	\centering
	\begin{threeparttable}
		\caption{Трёхсторонний чёт"~нечет}
		\label{tab:coin3b}
		\begin{tabular}{ |c|c|c|c|c| }
			\cline{1-2} \cline{4-5}
			\rule[-7pt]{0pt}{2em}$4, 4, 4$ &
			\rule[-7pt]{0pt}{2em}$6, 0, 6$ & \qquad\qquad\qquad &
			\rule[-7pt]{0pt}{2em}$6, 6, 0$ &
			\rule[-7pt]{0pt}{2em}$0, 6, 6$ \\
			\cline{1-2} \cline{4-5}
			\rule[-7pt]{0pt}{2em}$0, 6, 6$ &
			\rule[-7pt]{0pt}{2em}$6, 6, 0$ & \qquad\qquad\qquad &
			\rule[-7pt]{0pt}{2em}$6, 0, 6$ &
			\rule[-7pt]{0pt}{2em}$4, 4, 4$ \\
			\cline{1-2} \cline{4-5}
		\end{tabular}
	\end{threeparttable}
\end{table}

В соответствии с формулировкой теоремы для получения множества равновесий по Нэшу в повторяющейся игре необходимо построить минимаксные стратегии наказания в отдельной её итерации. Однако трёхсторонний чёт-нечет любопытен тем, что в смешанных стратегиях никакая попытка двух игроков минимизировать выигрыш третьего не может опустить его результат ниже того, что он получает в точках равновесия по Нэшу (т.е. 4). Действительно, пусть первый игрок избрал стратегию $(p, 1 - p)$, а второй "--- $(q, 1 - q)$. Тогда выигрыши третьего при выборе чистых стратегий составляют:
\begin{align*}
	u^3(p, q, 1) &= 6 (p + q) - 8 p q; \\
	u^3(p, q, 0) &= 2 (p + q) - 8 p q + 4.
\end{align*}

Несложно заметить, что
\begin{align*}
	p \ge &q \ge \frac{1}{2} \Rightarrow u^3(p, q, 1) \ge 4; \\
	p \ge 1 - &q \ge \frac{1}{2} \Rightarrow u^3(p, q, 1) \ge 4; \\
	q \ge &p \ge \frac{1}{2} \Rightarrow u^3(p, q, 1) \ge 4; \\
	q \ge 1 - &p \ge \frac{1}{2} \Rightarrow u^3(p, q, 1) \ge 4; \\
	p \le &q \le \frac{1}{2} \Rightarrow u^3(p, q, 0) \ge 4; \\
	p \le 1 - &q \le \frac{1}{2} \Rightarrow u^3(p, q, 0) \ge 4; \\
	q \le &p \le \frac{1}{2} \Rightarrow u^3(p, q, 0) \ge 4; \\
	q \le 1 - &p \le \frac{1}{2} \Rightarrow u^3(p, q, 0) \ge 4. \\
\end{align*}

Эти варианты исчерпывают всё пространство возможных ситуаций, так что никакое сочетание смешанных стратегий двух игроков не может быть эффективным наказанием для третьего. В классическом случае это означало бы, что повторяющаяся игра имеет в качестве решения то же самое множество равновесий по Нэшу, что и отдельный розыгрыш. Однако, вывод этот можно интересным образом уточнить, перейдя к модели, учитывающей стоимость вычислений, необходимых для выбора стратегии в очередной итерации.

\section{Модель повторяющихся игр с учётом стоимости вычислений}\label{sec:ch3/sect2}

Построение искомой модели подразумевает конкретизацию способа, при помощи которого рациональные агенты вычисляют стратегию поведения. Для этой цели подойдёт любой формализм, позволяющий алгоритмически полные вычисления с вероятностным ветвлением. Кроме того, необходимы возможность сохранения произвольного внутреннего состояния в памяти для использования на последующих итерациях и, очевидно, метод численной оценки сложности произведённого на каждой итерации вычисления. Общая схема взаимодействия агентов и среды может выглядеть так:
\begin{figure}[ht]
	\centerfloat{
		\ifdefmacro{\tikzsetnextfilename}{\tikzsetnextfilename{tikz_repeat_compiled}}{}
		\begin{tikzpicture}[scale=2]
			\node[circle,draw] (game1) at (1, 0) {$\Gamma$};
			\node[circle,draw] (game2) at (3, 0) {$\Gamma$};
			\node[circle,draw] (game3) at (5, 0) {$\Gamma$};
			\node[rectangle,draw] (calc11) at (0,  1) {$M^1$};
			\node[rectangle,draw] (calc12) at (2,  1) {$M^1$};
			\node[rectangle,draw] (calc13) at (4,  1) {$M^1$};
			\node[] (calc1i) at (6, 1) {$\cdots$};
			\node[rectangle,draw] (calc21) at (0, -1) {$M^2$};
			\node[rectangle,draw] (calc22) at (2, -1) {$M^2$};
			\node[rectangle,draw] (calc23) at (4, -1) {$M^2$};
			\node[] (calc2i) at (6, -1) {$\cdots$};
			\draw [->,thick,decorate,decoration={snake,post length=1mm}] (calc11) -- (game1) node[midway,sloped,below] {$s^1_1$};
			\draw [->,thick,decorate,decoration={snake,post length=1mm}] (calc21) -- (game1) node[midway,sloped,above] {$s^2_1$};
			\draw [->,thick,decorate,decoration={snake,post length=1mm}] (calc12) -- (game2) node[midway,sloped,below] {$s^1_2$};
			\draw [->,thick,decorate,decoration={snake,post length=1mm}] (calc22) -- (game2) node[midway,sloped,above] {$s^2_2$};
			\draw [->,thick,decorate,decoration={snake,post length=1mm}] (calc13) -- (game3) node[midway,sloped,below] {$s^1_3$};
			\draw [->,thick,decorate,decoration={snake,post length=1mm}] (calc23) -- (game3) node[midway,sloped,above] {$s^2_3$};
			\draw [->,thick,decorate,decoration={snake,post length=1mm}] (game1) -- (calc12) node[midway,sloped,below] {$s_1$};
			\draw [->,thick,decorate,decoration={snake,post length=1mm}] (game1) -- (calc22) node[midway,sloped,above] {$s_1$};
			\draw [->,thick,decorate,decoration={snake,post length=1mm}] (game2) -- (calc13) node[midway,sloped,below] {$s_2$};
			\draw [->,thick,decorate,decoration={snake,post length=1mm}] (game2) -- (calc23) node[midway,sloped,above] {$s_2$};
			\draw [->,thick,decorate,decoration={snake,post length=1mm}] (game3) -- (calc1i) node[midway,sloped,below] {$s_3$};
			\draw [->,thick,decorate,decoration={snake,post length=1mm}] (game3) -- (calc2i) node[midway,sloped,above] {$s_3$};
			\node[diamond,draw] (sum1) at (-1, 2) {$\Sigma^1$};
			\node[] (sum10) at (0,  2) [label=$-\delta \mathfrak{w}^1_1$] {};
			\node[] (sum11) at (1,  2) [label=$+\delta u^1(s_1)$] {};
			\node[] (sum12) at (2,  2) [label=$-\delta^2 \mathfrak{w}^1_2$] {};
			\node[] (sum13) at (3,  2) [label=$+\delta^2 u^1(s_2)$] {};
			\node[] (sum14) at (4,  2) [label=$-\delta^3 \mathfrak{w}^1_3$] {};
			\node[] (sum15) at (5,  2) [label=$+\delta^3 u^1(s_3)$] {};
			\node[] (sum1i) at (6,  2) {$\cdots$};
			\node[diamond,draw] (sum2) at (-1, -2) {$\Sigma^2$};
			\node[] (sum20) at (0, -2) [label=below:$-\delta \mathfrak{w}^2_1$] {};
			\node[] (sum21) at (1, -2) [label=below:$+\delta u^2(s_1)$] {};
			\node[] (sum22) at (2, -2) [label=below:$-\delta^2 \mathfrak{w}^2_2$] {};
			\node[] (sum23) at (3, -2) [label=below:$+\delta^2 u^2(s_2)$] {};
			\node[] (sum24) at (4, -2) [label=below:$-\delta^3 \mathfrak{w}^2_3$] {};
			\node[] (sum25) at (5, -2) [label=below:$+\delta^3 u^2(s_3)$] {};
			\node[] (sum2i) at (6, -2) {$\cdots$};
			\draw [->,thick] (sum1i) -- (sum1);
			\draw [->,thick] (calc11) -- (sum10);
			\draw [->,thick] (game1) -- (sum11);
			\draw [->,thick] (calc12) -- (sum12);
			\draw [->,thick] (game2) -- (sum13);
			\draw [->,thick] (calc13) -- (sum14);
			\draw [->,thick] (game3) -- (sum15);
			\draw [->,thick] (sum2i) -- (sum2);
			\draw [->,thick] (calc21) -- (sum20);
			\draw [->,thick] (game1) -- (sum21);
			\draw [->,thick] (calc22) -- (sum22);
			\draw [->,thick] (game2) -- (sum23);
			\draw [->,thick] (calc23) -- (sum24);
			\draw [->,thick] (game3) -- (sum25);
			\fill [white] (1,  1) circle (2pt);
			\fill [white] (3,  1) circle (2pt);
			\fill [white] (5,  1) circle (2pt);
			\fill [white] (1, -1) circle (2pt);
			\fill [white] (3, -1) circle (2pt);
			\fill [white] (5, -1) circle (2pt);
			\draw [->,double,thick] (-1,  1) -- (calc11) node[midway,above] {$\psi^1_0$};
			\draw [->,double,thick] (calc11) -- (calc12) node[near end,above] {$\psi^1_1$};
			\draw [->,double,thick] (calc12) -- (calc13) node[near end,above] {$\psi^1_2$};
			\draw [->,double,thick] (calc13) -- (calc1i) node[near end,above] {$\psi^1_3$};
			\draw [->,double,thick] (-1, -1) -- (calc21) node[midway,below] {$\psi^2_0$};
			\draw [->,double,thick] (calc21) -- (calc22) node[near end,below] {$\psi^2_1$};
			\draw [->,double,thick] (calc22) -- (calc23) node[near end,below] {$\psi^2_2$};
			\draw [->,double,thick] (calc23) -- (calc2i) node[near end,below] {$\psi^2_3$};
		\end{tikzpicture}
	}
	\legend{}
	\caption[Повторяющаяся игра с учётом стоимости вычислений]{Повторяющаяся игра с учётом стоимости вычислений}\label{fig:repeat}
\end{figure}

Диаграмма на рисунке \ref{fig:repeat} схематично изображает моделируемый процесс для двух игроков (естественным образом обобщающийся на любое конечное их число). В узлах, помеченных буквой $\Gamma$, происходят последовательные розыгрыши произвольной игры $\Gamma = \langle A, S^a, u^a(s), a \in A \rangle$. В $i$-м розыгрыше игрок $a$ выбирает свою стратегию $s^a_i$, применяя вероятностный алгоритм $M^a$ к результату предыдущей итерации $(\psi^a_{i-1}, s_{i-1})$, включающему сохранённое состояние памяти самого алгоритма и набор сыгранных на итерации $i-1$ стратегий. Узлы $\Sigma^a$ изображают последовательное суммирование разностей выигрыша $u^a(s_i)$ и затрат на произведённое вычисление $\mathfrak{w}^a_i$, с учётом экспоненциально уменьшающегося коэффициента дисконтирования $\delta^i$. При такой схеме взаимодействий имеет смысл рассматривать только универсальные алгоритмы $M^a$, что позволяет закодировать любой набор вычислимых стратегий повторяющейся игры в начальных состояниях памяти $\psi^a_0$.

\section{Криптографические стратегии наказания}

Хотя трёхсторонний чёт-нечет не даёт возможности двум игрокам наказывать третьего, используя только лишь смешанные стратегии, ещё в первой главе мы установили, что, например, наличие механизма корреляции, которым игроки 1 и 2 могут пользоваться втайне от игрока 3, добавляет решения с выплатами до $(5, 5, 2)$. Таким образом, если бы мы рассматривали повторяющуюся игру в пространстве заговоров структуры $\{\{1, 2\}, \{2, 3\}, \{3, 1\}\}$, то в соответствии с <<народной>> теоремой её множество совершенных подыгровых равновесий включало бы все стратегические наборы с выплатами, обеспечивающие каждому игроку выигрыш не менее $2$. Однако, даже в том случае, когда игроки не могут использовать тайные механизмы корреляции, учёт стоимости вычислений позволяет в повторяющихся играх применять стратегии наказания, опирающиеся на достижения современной криптографии.

%Тогда стратегию игрока $a$ можно представить в следующем виде:
%\begin{equation*}
%	\overbrace{s}^a = \langle \psi^a \in \Psi^a, f^a : \Psi^a \rightarrow \mathcal{M}(X^a), g^a :  / \rangle
%\end{equation*}

\clearpage
           % Глава 3
\chapter*{Заключение}                       % Заголовок
\addcontentsline{toc}{chapter}{Заключение}  % Добавляем его в оглавление

%% Согласно ГОСТ Р 7.0.11-2011:
%% 5.3.3 В заключении диссертации излагают итоги выполненного исследования, рекомендации, перспективы дальнейшей разработки темы.
%% 9.2.3 В заключении автореферата диссертации излагают итоги данного исследования, рекомендации и перспективы дальнейшей разработки темы.
%% Поэтому имеет смысл сделать эту часть общей и загрузить из одного файла в автореферат и в диссертацию:

Основные результаты работы заключаются в следующем.
%% Согласно ГОСТ Р 7.0.11-2011:
%% 5.3.3 В заключении диссертации излагают итоги выполненного исследования, рекомендации, перспективы дальнейшей разработки темы.
%% 9.2.3 В заключении автореферата диссертации излагают итоги данного исследования, рекомендации и перспективы дальнейшей разработки темы.
\begin{enumerate}
  \item На основе анализа понятия коррелированного равновесия в контексте многосторонних конфликтов было сформулировано свойство чувствительности игр к дополнительной информационной асимметрии.
  \item Исследование изоморфизма пространств корреляции позволило ввести сужающий их формализм пространства заговоров, с применением которого удобно рассуждать о влиянии дополнительной информационной асимметрии на решения игр.
  \item Моделирование проблемы планирования заданий при допущении немонотонности функций отдачи показало, что в пространствах заговоров концепция структурной согласованности равновесий может использоваться как функциональный аналог классических критериев коллективной рациональности по отношению к равновесиям Нэша в смешанных стратегиях. 
  \item Для демонстрации значимости феномена чувствительности игр к дополнительной информационной асимметрии была построена модель повторяющихся конфликтов с учётом стоимости вычисления очередного шага стратегии.
  \item В рамках построенной модели было показано, как в повторяющихся играх можно даже без дополнительной информационной асимметрии как таковой использовать современные криптографические примитивы для конструирования эффективных стратегий наказания, использующих чувствительность к ней.
\end{enumerate}

Хочется надеяться, что эта работа привлечёт внимание специалистов к проблематике влияния дополнительной информационной асимметрии на исходы многосторонних конфликтов.

В заключение автор выражает благодарность и большую признательность научному руководителю
Васину~А.\,А. за поддержку, помощь, обсуждение результатов и научное руководство. Также автор благодарит Морозова~В.\,В. за активное участие в работе над доказательствами теорем и авторов шаблона *Russian-Phd-LaTeX-Dissertation-Template* за помощь в оформлении диссертации.
      % Заключение
\include{Dissertation/acronyms}        % Список сокращений и условных обозначений
\chapter*{Словарь терминов}             % Заголовок
\addcontentsline{toc}{chapter}{Словарь терминов}  % Добавляем его в оглавление

\textbf{пространство корреляции} : Дополнительный параметр коррелированного расширения игр в нормальной форме, характеризующий априорные знания игроков о событиях, напрямую не влияющих на исход конфликта

\textbf{пространство заговоров} : Пространство корреляции специальной структуры с простым конечным описанием в виде семейства групп игроков, называемых заговорами

\textbf{заговор} : Подмножество игроков, объединённое возможностью наблюдать за общим тайным механизмом корреляции, сигналы которого непредсказуемы для аутсайдеров
      % Словарь терминов
\include{Dissertation/references}      % Список литературы
\clearpage
\ifdefmacro{\microtypesetup}{\microtypesetup{protrusion=false}}{} % не рекомендуется применять пакет микротипографики к автоматически генерируемым спискам
\listoffigures  % Список изображений

%%% Список таблиц %%%
% (ГОСТ Р 7.0.11-2011, 5.3.10)
%\clearpage
%\listoftables   % Список таблиц
%\ifdefmacro{\microtypesetup}{\microtypesetup{protrusion=true}}{}
%\newpage           % Списки таблиц и изображений (иллюстративный материал)

\setcounter{totalchapter}{\value{chapter}} % Подсчёт количества глав

%%% Настройки для приложений
\appendix
% Оформление заголовков приложений ближе к ГОСТ:
\setlength{\midchapskip}{20pt}
\renewcommand*{\afterchapternum}{\par\nobreak\vskip \midchapskip}
\renewcommand\thechapter{\Asbuk{chapter}} % Чтобы приложения русскими буквами нумеровались

\chapter{Краткий обзор литературы, посвящённой коррелированному расширению игр в нормальной форме}\label{app:A}

Концепция коррелированного расширения стала важным общеупотребимым инструментом во многих областях теории игр, и в частности на неё опирается знаменитый дизайн механизмов. В связи с этим нельзя не упомянуть статью Роджера Майерсона <<Optimal coordination mechanisms in generalized principal"~agent problems>> \cite{Myerson82}. В ней формулируется обобщённая задача принципал"~агентов, в которой агенты обладают как тайной информацией, так и возможностью принимать решения, неподконтрольные принципалу. Показывается, что принципал может ограничиваться целенаправленно побуждающими прямыми механизмами координации, в которых агенты докладывают свою информацию принципалу, рекомендующему в ответ стратегии, образующие коррелированное равновесие. В конечном случае оптимальные механизмы координации могут быть найдены при помощи линейного программирования. Кроме того обсуждается проблематика систем с многими принципалами, в которых может не существовать некооперативного равновесия, так что вводится определение и показывается существование квази"~равновесия.

Коррелированное расширение нашло своё место и применительно к исследованию игр в развёрнутой форме. Ещё одна статья Роджера Майерсона <<Multistage Games with Communication>> \cite{Myerson86} рассматривает многостадийные игры с коммуникационным механизмом, функционирующим по принципу централизованного посредника. В коммуникационном равновесии ни один игрок не должен иметь возможность в одиночку увеличить свой выигрыш, манипулируя своими отчётами или действиями. Последовательное коммуникационное равновесие "--- это коммуникационное равновесие с системой условных вероятностей при которой ни один игрок не может получить выгоду от подобных манипуляций, даже если случаются события нулевой вероятности. Кодоминируемые действия определяются таким образом, что любое коммуникационное равновесие последовательно в том и только том случае, когда никто не использует кодоминируемых действий. Преобладающее коммуникационное равновесие определяется как результат последовательного исключения кодоминируемых действий, и показывается его существование.

Ещё одной важной вехой стала статья <<Correlated equilibrium as an expression of Bayesian rationality>> \cite{Aumann87}, в которой Ауман показал, что формализм коррелированного равновесия снимает противоречие между <<байесовским>> и <<теоретико"~игровым>> взглядом на мир. С байесовской позиции вероятности могут быть сопоставлены с чем угодно, даже с возможностью для игрока выбрать какую"~либо стратегию в некоторой игре. С точки зрения самой теории игр, напротив, традиционно считается, что нельзя говорить о вероятностях событий, происходящих по воле рациональных агентов, нужно вместо этого использовать понятие равновесности (или другие теоретико-игровые конструкты). Предложенный формализм же объединяет две эти точки зрения "--- на коррелированное равновесие можно смотреть как на следствие байесовской рациональности, поскольку условие равновесности представляет собой простую максимизацию выгоды каждым из игроков с учётом известной им информации. При таком подходе не требуется явной рандомизации в действиях игроков. Даже если игрок выбирает конкретную чистую стратегию без элемента случайности, вероятностная природа стратегий отражает неуверенность остальных игроков в его выборе, что и показывается на примерах.

Вопросы совместимости коррелированных равновесий с более строгими принципами оптимальности поднимались в статье Амриты Дхиллона и Жана Франсуа Мертенса <<Perfect Correlated Equilibria>> \cite{Dhillon}. В ней вводится понятие $(\epsilon)$"~совершенных кореллированных равновесий (PCE), обусловленных $(\epsilon)$"~совершенным равновесием некоторого корреляционного устройства. Показывается, что <<принцип выявления>> для этой концепции теряет силу "--- прямой механизм может и не обеспечивать совершенного равновесия. Так называемые приблизительно совершенные коррелированные равновесия (APCE) оказываются пределами $\epsilon$"~PCE, и авторы достигают для них полной характеризации. В ходе рассуждений о <<приемлемости>> APCE в некотором смысле, приводятся однако иллюстрированные доводы в пользу того, что среди них именно PCE представляются <<хорошими>>.

Динамический аспект формализма коррелированного расширения тоже не оставался без внимания исследователей. Предлагалось немало процедур разыгрывания для итеративных игр, обеспечивающих сходимость к коррелированным равновесиям. Среди множества работ на эту тему выделяется статья Серджиу Харта и Андреу Мас"~Колелла <<A simple adaptive procedure leading to correlated equilibrium>> \cite{Hart}, в которой авторы предложили так называемую процедуру сопоставления потерь. Применяя её, игроки каждый раз отклоняются от своих текущих стратегий пропорционально мере понесённого ими на предыдущих ходах ущерба от не использования иных стратегий. Показывается, что такая адаптивная процедура гарантирует в любой игре сходимость с вероятностью $1$ эмпирического распределения розыгрышей к множеству коррелированных равновесий.

\chapter{Доказательство теоремы об изоморфизме пространств корреляции}\label{app:B}

Прежде чем перейти к доказательству основного утверждения потребуется ввести дополнительный инструментарий:

\begin{definition}
	Измельчением множества исходов $X = X^1 \times \ldots \times X^m$ до конечного множества исходов $Y = Y^1 \times \ldots \times Y^m$ называется любое отображение $\rho = (\rho^1, \ldots, \rho^m)$, где каждая компонента $\rho^a$ отображает $Y^a$ в $X^a$.
\end{definition}

При помощи измельчений можно задавать связи между разбиениями с различными кодоменами. Если разбиения $f : \Omega \rightarrow X$ и $g : \Omega \rightarrow Y$ таковы, что $f = \rho \circ g$, то $f^{-1}(x) = \bigcup_{y \in \rho^{-1}(x)} g^{-1}(y), \forall x \in X$. При этом $f$ можно называть измельчимым до $g$.

Разбиения одного и того же пространства можно комбинировать. Например, из разбиений $g_i : \Omega \rightarrow Y_i = Y_i^1 \times \ldots \times Y_i^m, i=\overline{1,n}$ можно построить их комбинацию $g_1 \diamond \ldots \diamond g_n : \Omega \rightarrow Y_{(n)}$, где $Y_{(n)}^a = Y_1^a \times \ldots \times Y_n^a$ и $(g_1 \diamond \ldots \diamond g_n)^a(\omega) = (g_1^a(\omega), \ldots, g_n^a(\omega)), \forall \omega \in \Omega, a =\overline{1,m}$. Эта комбинация разбиений связана со своими компонентами измельчениями"~проекциями: $g_i  = \pi_i \circ (g_1 \diamond \ldots \diamond g_n), \pi_i^a(x_1^a, \ldots, x_n^a) = x_i^a$.

Аналогично комбинируются и измельчения с общим кодоменом. Например, из измельчений $\rho_i : Y_i \rightarrow X, i=\overline{1,n}$ можно построить комбинацию $\rho_1 \wr \ldots \wr \rho_n : Y_{[n]} \rightarrow X$, где $Y_{[n]}^a = \{(y_1^a, \ldots, y_n^a) \in Y_{(n)}^a \mid \rho_1^a(y_1^a) = \ldots = \rho_n^a(y_n^a)\}, a =\overline{1,m}$, причём на своей области определения $\rho_1 \wr \ldots \wr \rho_n$ совпадает со всеми $\rho_i \circ \pi_i$. Заметим, что
\begin{equation*}
	f = \rho_i \circ g_i, i=\overline{1,n} \;\Leftrightarrow\; f = (\rho_1 \wr \ldots \wr \rho_n) \circ (g_1 \diamond \ldots \diamond g_n).
\end{equation*}

\begin{definition}
	В пространстве корреляции $\Phi = \langle A, \Omega, \mathfrak{I}^a, \mathbb{P}, a \in A \rangle$ структурой разбиения $f : \Omega \rightarrow X$, порождённой измельчением $\rho : Y \rightarrow X$, называется множество $H_{\Phi,\rho}(f) = \{\mathbb{P} \circ g^{-1} \mid g \models \Phi, f = \rho \circ g\}$, состоящее из мер $\mu : Y \rightarrow \mathbb{R}_{\ge 0}$. Также обозначим $H_{\Phi,\rho}^{-1}(\mu) = \{f \models \Phi \mid \mu \in H_{\Phi,\rho}(f)\}$.
\end{definition}

\begin{lemma} \label{lemma:tyhon}
	Для всех $\rho : Y \rightarrow X$ и $\mu : Y \rightarrow \mathbb{R}_{\ge 0}$ множество $H_{\Phi,\rho}^{-1}(\mu) \subseteq X^{\Omega}$ компактно в полуметрике
	\begin{equation*}
		\operatorname{dis}(f_1, f_2) = \frac{1}{2} \sum_{x \in X} \left| \mathbb{P}(f_1^{-1}(x)) - \mathbb{P}(f_2^{-1}(x)) \right|.
	\end{equation*}
\end{lemma}

\begin{proof}
	Переформулируем $H_{\Phi,\rho}^{-1}(\mu) = \rho \circ H_{\Phi}^{-1}(\mu)$, определив для этого отображение $H_{\Phi}^{-1}(\mu) = \{g \models \Phi \mid \mathbb{P} \circ g^{-1} = \mu\}$. Докажем сперва компактность $H_{\Phi}^{-1}(\mu)$, вводя $\operatorname{dis}(g_1, g_2)$ аналогично $\operatorname{dis}(f_1, f_2)$. Полуметрика $\operatorname{dis}$ вполне ограниченна, так как $\operatorname{dis}(g_1, g_2) = d(\mu^Y_1, \mu^Y_2)$, где $\mu^Y_k = \mathbb{P} \circ g_k^{-1}, k=1,2$, а пространство вероятностных мер на любом конечном множестве вполне ограниченно. Замкнутость $H_{\Phi}^{-1}(\mu)$ очевидно следует из $\operatorname{dis}(g_1, g_2) = 0 \Leftrightarrow \mathbb{P} \circ g_1^{-1} = \mathbb{P} \circ g_2^{-1}$. Таким образом $H_{\Phi}^{-1}(\mu)$ компактно в полуметрике $\operatorname{dis}$. Докажем непрерывность отображения $\rho\ \circ : Y^\Omega \rightarrow X^\Omega$, выразив ту же полуметрику по-другому:
	\begin{equation*}
		\operatorname{dis}(f_1, f_2) = 1 - \sum_{x \in X} \min \left[\mathbb{P}(f_1^{-1}(x)), \mathbb{P}(f_2^{-1}(x))\right].
	\end{equation*}
	
	Пусть теперь $f_1 = \rho \circ g_1$ и $f_2 = \rho \circ g_2$:
	\begin{align*}
		\operatorname{dis}(\rho \circ g_1, \rho \circ g_2) &= 1 - \sum_{x \in X} \min \left[\mathbb{P}(g_1^{-1}(\rho^{-1}(x))), \mathbb{P}(g_2^{-1}(\rho^{-1}(x)))\right] \\
		&= 1 - \sum_{x \in X} \min \left[\sum_{y \in \rho^{-1}(x)} \mathbb{P}(g_1^{-1}(y)), \sum_{y \in \rho^{-1}(x)} \mathbb{P}(g_2^{-1}(y))\right] \\
		&\le 1 - \sum_{x \in X} \sum_{y \in \rho^{-1}(x)} \min \left[\mathbb{P}(g_1^{-1}(y)), \mathbb{P}(g_2^{-1}(y))\right] \\
		&= 1 - \sum_{y \in Y} \min \left[\mathbb{P}(g_1^{-1}(y)), \mathbb{P}(g_2^{-1}(y))\right] = \operatorname{dis}(g_1, g_2).
	\end{align*}
	
	Отображение $\rho\ \circ$ непрерывно, поскольку $\operatorname{dis}(\rho \circ g_1, \rho \circ g_2) \le \operatorname{dis}(g_1, g_2)$. Так как непрерывные отображения сохраняют компактность\cite[с.~199]{Engelking}, $H_{\Phi,\rho}^{-1}(\mu) = \rho \circ H_{\Phi}^{-1}(\mu)$ компактно в полуметрике $\operatorname{dis}$.
\end{proof}

\begin{definition}
	Разбиение $f_2 \models \Phi_2$ называется точным образом разбиения $f_1 \models \Phi_1$ (далее $f_1 \precsim f_2$), если их кодомены совпадают ($X_1 = X_2 = X$) и $H_{\Phi_1,\rho}(f_1) \subseteq H_{\Phi_2,\rho}(f_2)$ для всех измельчений $\rho$ с тем же кодоменом. Множество всех точных образов далее будем обозначать $\widehat{\Phi}_2(f_1) = \{f_2 \models \Phi_2 \mid f_1 \precsim f_2\}$.
\end{definition}

Отношение $f_1 \precsim f_2$ можно понять так "--- на какие бы измеримые части мы не делили компоненты разбиения $f_1$, в разбиении $f_2$ соответствующие компоненты всегда можно разделить на равные им по мере части.

\begin{note}
	Очевидно, что $f_1 \precsim f_2 \wedge f_2 \precsim f_3 \Rightarrow f_1 \precsim f_3$.
\end{note}

\begin{lemma} \label{lemma:up}
	Пусть в пространствах корреляции $\Phi_1$ и $\Phi_2$ разбиения $g_1 : \Omega_1 \rightarrow Y$ и $g_2 : \Omega_2 \rightarrow Y$ таковы, что $g_1 \precsim g_2$. Тогда $\rho \circ g_1 \precsim \rho \circ g_2$ для всех измельчений $\rho : Y \rightarrow X$.
\end{lemma}

\begin{proof}
	Возьмём любые $\rho_* : Y_* \rightarrow X$ и $\mu \in H_{\Phi_1,\rho_*}(\rho \circ g_1)$. По определению структуры разбиения, $\exists g_{1*}: \rho_* \circ g_{1*} = \rho \circ g_1, \mathbb{P}_1 \circ g_{1*}^{-1} = \mu$, а доказать требуется, по определению точного образа, что $\exists g_{2*}: \rho_* \circ g_{2*} = \rho \circ g_2, \mathbb{P}_2 \circ g_{2*}^{-1} = \mu$. Рассмотрим комбинацию $g_{1+} = g_1 \diamond g_{1*}$, где $g_1 = \pi \circ g_{1+}$ и $g_{1*} = \pi_* \circ g_{1+}$. Здесь $g_{1+} : \Omega_1 \rightarrow Y_+, Y_+^a = Y^a \times Y_*^a, a=\overline{1,m}$. По определению структуры разбиения, $\mathbb{P}_1 \circ g_{1+}^{-1} \in H_{\Phi_1,\pi}(g_1)$, а значит, поскольку $g_1 \precsim g_2$, существует $g_{2+} : \Omega_2 \rightarrow Y_+$ такое, что $\mathbb{P}_1 \circ g_{1+}^{-1} = \mathbb{P}_2 \circ g_{2+}^{-1} \in H_{\Phi_2,\pi}(g_2)$, т.е. $\pi \circ g_{2+} = g_2$. Из этого с очевидностью следует, что и $\mathbb{P}_2 \circ (\pi_* \circ g_{2+})^{-1} = \mathbb{P}_1 \circ (\pi_* \circ g_{1+})^{-1}$, а значит $g_{2*} = \pi_* \circ g_{2+}$ искомое.
\end{proof}

\begin{lemma} \label{lemma:down}
	Пусть в пространствах корреляции $\Phi_1$ и $\Phi_2$ разбиения $f_1 : \Omega_1 \rightarrow X$ и $f_2 : \Omega_2 \rightarrow X$ таковы, что $f_1 \precsim f_2$. Тогда для каждого измельчения $\rho : Y \rightarrow X$ и каждого разбиения $g_1 : \Omega_1 \rightarrow Y$ такого, что $f_1 = \rho \circ g_1$ существует разбиение $g_2 : \Omega_2 \rightarrow Y$ такое, что $f_2 = \rho \circ g_2$ и $g_1 \precsim g_2$.
\end{lemma}

\begin{proof}
	Сформулируем требуемое как $\exists g_2 \in \widehat{\Phi}_2(g_1) : f_2 = \rho \circ g_2$ и выразим $\widehat{\Phi}_2$ через структуры разбиений:
	\begin{equation*}% \label{eq:goal1}
		\widehat{\Phi}_2(g_1) = \bigcap_{\forall Z, \xi : Z \rightarrow Y, \mu \in H_{\Phi_1,\xi}(g_1)} H_{\Phi_2,\xi}^{-1}(\mu).
	\end{equation*}
	
	По лемме \ref{lemma:tyhon} множество $\widehat{\Phi}_2(g_1)$ является пересечением семейства компактов. Следовательно, для доказательства содержания в нём элемента $g_2 : f_2 = \rho \circ g_2$, достаточно доказать, что такой элемент содержится в пересечении каждого конечного подсемейства тех же компактов:
	\begin{equation*}% \label{eq:goal1}
		\exists g_{2*} \in \bigcap_{i=1}^n H_{\Phi_2,\xi_i}^{-1}(\mu_i) : f_2 = \rho \circ g_{2*},
	\end{equation*}
	где $\xi_i : Z_i \rightarrow Y$ "--- произвольные измельчения с произвольными доменами $Z_i$ и $\mu_i \in H_{\Phi_1,\xi_i}(g_1)$ также выбраны произвольно.
	
	По определению структуры разбиения $\exists h_{1,i} \models \Phi_1 : g_1 = \xi_i \circ h_{1,i}, \mathbb{P} \circ h_{1,i}^{-1} = \mu_i$. Построим их комбинацию $h_1 = h_{1,1} \diamond ... \diamond h_{1,n}$, где $h_{1,i} = \pi_i \circ h_1$, и обозначим $\xi = \xi_1 \wr ... \wr \xi_n$. По определению точного отображения $\exists h_2 \models \Phi_2 : f_2 = \rho \circ \xi \circ h_2, \mathbb{P}_1 \circ h_1^{-1} = \mathbb{P}_2 \circ h_2^{-1}$, а значит можно взять $g_{2*} = \xi \circ h_2$. По построению $f_2 = \rho \circ g_{2*}$ и $\mathbb{P}_1 \circ h_{1,i}^{-1} = \mathbb{P}_1 \circ (\pi_i \circ h_1)^{-1} = \mathbb{P}_2 \circ (\pi_i \circ h_2)^{-1} = \mathbb{P}_2 \circ h_{2,i}^{-1}$, следовательно $g_{2*}$ "--- искомое.
\end{proof}

\begin{corollary} \label{cor:down}
	Если пространства корреляции $\Phi_1 \precsim \Phi_2$, то для каждого разбиения $f_1 \models \Phi_1$ существует $f_2 \models \Phi_2$ такое, что $f_1 \precsim f_2$.
\end{corollary}

\begin{corollary}
	Леммы \ref{lemma:up}, \ref{lemma:down} и следствие \ref{cor:down} также верны для строгого отношения $f_1 \prec f_2 \equiv f_1 \precsim f_2 \cap \neg (f_1 \succsim f_2)$.
\end{corollary}

\begin{lemma}
	$f_1 \precsim f_2 \Leftrightarrow f_1 \succsim f_2$ для любых разбиений одного и того же пространства корреляции.
\end{lemma}

\begin{proof}
	Предположим обратное "--- существование $f_1 \prec f_2$ с кодоменом $X$. Тривиальное измельчение $\theta(x) = (0, \ldots, 0), \forall x \in X$ очевидно даёт $\theta \circ f_1 = \theta \circ f_2$. Это противоречит $\theta \circ f_1 \prec \theta \circ f_2$, следующему из леммы \ref{lemma:up}.
\end{proof}

\begin{corollary} \label{cor:sym}
	$f_1 \precsim f_2 \Leftrightarrow f_1 \succsim f_2$ для любых разбиений изоморфных пространств корреляции.
\end{corollary}

\begin{proof}[Доказательство теоремы об изоморфных пространствах]
	Рассмотрим в игре $\Gamma | \Phi_1$ произвольный профиль стратегий $\mathbf{s}_1$. Этот профиль, очевидно, является разбиением пространства корреляции $\Phi_1$. По следствию \ref{cor:down} существует разбиение $\mathbf{s}_2$ пространства корреляции $\Phi_2$ такое, что $\mathbf{s}_1 \precsim \mathbf{s}_2$, причём, аналогично, $\mathbf{s}_2$ является ещё и профилем стратегий в игре $\Gamma | \Phi_2$. Докажем вложения в обоих направлениях: 1. $U_{\Gamma | \Phi_1}^{A_*}(\mathbf{s}_1) \subseteq U_{\Gamma | \Phi_2}^{A_*}(\mathbf{s}_2)$ и 2. $U_{\Gamma | \Phi_1}^{A_*}(\mathbf{s}_1) \supseteq U_{\Gamma | \Phi_2}^{A_*}(\mathbf{s}_2)$ для любой группы игроков $A_*$:
	\begin{enumerate}
		\item Рассмотрим произвольный профиль $\mathbf{s}_{1*} \models \Phi_1$, отличающийся от $\mathbf{s}_1$ стратегиями группы $A_*$. Обозначим $\mathbf{s}_{1+} = \mathbf{s}_1 \diamond \mathbf{s}_{1*}$, где $\mathbf{s}_1 = \pi \circ \mathbf{s}_{1+}$ и $\mathbf{s}_{1*} = \pi_* \circ \mathbf{s}_{1+}$. По определению точного образа $H_{\Phi_1,\pi}(\mathbf{s}_1) \subseteq H_{\Phi_2,\pi}(\mathbf{s}_2)$, т.е. $\exists \mathbf{s}_{2+} \models \Phi_2 : \mathbb{P}_1 \circ \mathbf{s}_{1+} = \mathbb{P}_2 \circ \mathbf{s}_{2+}, \mathbf{s}_2 = \pi \circ \mathbf{s}_{2+}$. По построению $\mathbf{s}_{2*} = \pi_* \circ \mathbf{s}_{2+}$ отличается от $\mathbf{s}_2$ ходами тех же игроков, что отличают $\mathbf{s}_{1*}$ от $\mathbf{s}_1$, и $\mathbb{P}_1 \circ \mathbf{s}_{1*}^{-1} = \mathbb{P}_2 \circ \mathbf{s}_{2*}^{-1}$, а значит аналогичным образом $u^a(\mathbf{s}_{1*}) = u^a(\mathbf{s}_{2*})$. В силу произвольности выбора $\mathbf{s}_{1*}$ это влечёт $U_{\Gamma | \Phi_1}^{A_*}(\mathbf{s}_1) \subseteq U_{\Gamma | \Phi_2}^{A_*}(\mathbf{s}_2)$.
		\item Так как $\mathbf{s}_1 \succsim \mathbf{s}_2$ по следствию \ref{cor:sym}, рассуждения предыдущего пункта применимы и в обратном направлении.
	\end{enumerate}
\end{proof}

\chapter{Доказательство теоремы о пространствах заговоров одной структуры}\label{app:C}

Для доказательства теоремы об изоморфизме пространств заговоров так же понадобится несколько лемм.

\begin{lemma}\label{lemma:card}
	Для любого счётного семейства множеств $\mathfrak{F}$ найдётся цепь множеств $\mathfrak{T}$ такая, что $\sigma(\mathfrak{F}) = \sigma(\mathfrak{T})$.
\end{lemma}

\begin{proof}
	Пусть $\mathfrak{F} = \{F_1, F_2, \ldots\}$. Построим индуктивно последовательность цепей $(\mathfrak{T}_i)$, где каждая следующая цепь включает в себя предыдущую и $\sigma(\mathfrak{T}_i) = \sigma(\{F_1, \ldots, F_i\})$. В качестве базы возьмём $\mathfrak{T}_1 = \{F_1\}$. Шаг индукции: пусть $\mathfrak{T}_{i-1} = \{T_1, \ldots, T_n\}, T_1 \subset \ldots \subset T_n$ и $\sigma(\mathfrak{T}_{i-1}) = \sigma(\{F_1, \ldots, F_{i-1}\})$. Разложим следующий элемент $\mathfrak{F}$ на непересекающиеся дизъюнкты: $F_i = (F_i \cap T_1) \cup (F_i \cap T_2 \setminus T_1) \cup \ldots \cup (F_i \cap T_n \setminus T_{n-1}) \cup (F_i \setminus T_n)$. В этой записи $j$"~й дизъюнкт вложен в соответствующую разность $T_j \setminus T_{j-1}$ соседних элементов цепи. Следовательно, для его порождения достаточно пополнить $\mathfrak{T}_{i-1}$ множеством $T_{j-} = F_i \cap T_j \cup T_{j-1}$, сохраняющим структуру цепи, поскольку $T_{j-1} \subseteq T_{j-} \subseteq T_j$. Таким образом, чтобы получить $F_i$ целиком,
	\begin{equation*}
		\mathfrak{T}_i = \mathfrak{T}_{i-1} \cup \left\{
		\begin{aligned}
			&F_i \cap T_1,\\
			&F_i \cap T_2 \cup T_1,\\
			&\cdots\\
			&F_i \cap T_n \cup T_{n-1},\\
			&F_i \cup T_n
		\end{aligned}
		\right\}
	\end{equation*}
	
	Покажем, что предел последовательности $(\mathfrak{T}_i)$ "--- искомая цепь. В самом деле, любой элемент из $\sigma(\mathfrak{F})$ "--- это счётное объединение конечных пересечений множеств $F_i$. Поэтому
	\begin{equation*}
		\sigma(\mathfrak{F}) = \bigcup_{i=1}^{\infty} \sigma(\{F_1, \ldots, F_i\}) = \bigcup_{i=1}^{\infty} \sigma(\mathfrak{T}_i) = \sigma(\mathfrak{T}).
	\end{equation*}
\end{proof}

\begin{lemma} \label{lemma:max}
	Максимальная цепь измеримых множеств в безатомическом пространстве порождает безатомическую $\sigma$"~алгебру.
\end{lemma}

\begin{proof}
	Пусть максимальная цепь $\mathfrak{T}$ измеримых множеств безатомического пространства $\langle \Omega, \mathfrak{B}, \mathbb{P} \rangle$ порождает алгебру $\sigma(\mathfrak{T})$. Докажем, что для любого $B \in \sigma(\mathfrak{T})$ меры $\mathbb{P}(B) > 0$ найдётся $B' \in \sigma(\mathfrak{T})$ такое, что $B' \subset B$ и $\mathbb{P}(B) > \mathbb{P}(B') > 0$. Для этого, очевидно, достаточно доказать, что в цепи $\mathfrak{T}$ найдётся множество $T$ такое, что $0 < \mathbb{P}(T \cap B) < \mathbb{P}(B)$. Рассмотрим множества
	\begin{equation*}
		\underline{T} = \bigcup_{T_- \in \mathfrak{T} : \mathbb{P}(T_- \cap B) = 0} T_- \quad\text{и}\quad \overline{T} = \bigcap_{T_+ \in \mathfrak{T} : \mathbb{P}(T_+ \cap B) = \mathbb{P}(B)} T_+,
	\end{equation*}
	по построению вложенные $\underline{T} \subset \overline{T}$ так, что $\mathbb{P}(\overline{T}) - \mathbb{P}(\underline{T}) \ge \mathbb{P}(B)$. Поскольку цепь $\mathfrak{T}$ максимальна в безатомическом пространстве, существует $T_0 \in \mathfrak{T}$ такое, что $\underline{T} \subset T_0 \subset \overline{T}$. Так как $\underline{T} \subset T_0 \Rightarrow \mathbb{P}(T_0 \cap B) > 0$ и $T_0 \subset \overline{T} \Rightarrow \mathbb{P}(T_0 \cap B) < \mathbb{P}(B)$, значит $T_0$ искомое.
\end{proof}

\begin{definition}
	Для любых семейств измеримых множеств $\mathfrak{T} \subseteq 2^\Omega$ и мер $\mathbb{P} : \mathfrak{T} \rightarrow \mathbb{R}_{\ge 0}$ определим отображение $\operatorname{mim}\langle\mathfrak{T}, \mathbb{P}\rangle : \Omega \rightarrow \mathbb{R}_{\ge 0}$, называемое наименьшей мерой включения и вычисляемое по формуле $\operatorname{mim}\langle\mathfrak{T}, \mathbb{P}\rangle(\omega) = \inf\{\mathbb{P}(T) \mid \omega \in T \in \mathfrak{T}\}$.
\end{definition}

\begin{lemma}
	Если $\mathfrak{T} \subset 2^\Omega$ "--- цепь множеств, порождающая безатомическую $\sigma$"~алгебру, то $\mathbb{P} \circ \operatorname{mim}\langle\mathfrak{T}, \mathbb{P}\rangle^{-1}$ совпадает с мерой Лебега на отрезке $[0, \mathbb{P}(\Omega)]$.
\end{lemma}

\begin{proof}
	Поскольку $\mathfrak{T}$ "--- цепь, $\omega \in T \Leftrightarrow \operatorname{mim}\langle\mathfrak{T}, \mathbb{P}\rangle(\omega) \le \mathbb{P}(T), \forall \omega \in \Omega, T \in \mathfrak{T}$. Так как $\mathfrak{T}$ вдобавок порождает безатомическую $\sigma$"~алгебру, то для каждого $0 < t < \mathbb{P}(\Omega)$ найдётся $T \in \mathfrak{T}$ такое, что $\mathbb{P}(T) = t$. Следовательно, функция $\operatorname{mim}\langle\mathfrak{T}, \mathbb{P}\rangle$ отображает множества $T \in \mathfrak{T}$ на отрезки $[0, \mathbb{P}(T)]$, что очевидно влечёт цель доказательства.
\end{proof}

\begin{lemma} \label{lemma:consp}
	Пусть $\langle \Omega, \mathfrak{B}, \mathbb{P} \rangle$ "--- любое безатомическое вероятностное пространство с $\sigma$"~алгеброй, разложимой на $n$ безатомических компонент $\mathfrak{B} = \sigma(\mathfrak{B}_1 \cup \ldots \cup \mathfrak{B}_n)$ таких, что все события из разных компонент совместно независимы, т.е. $\mathbb{P}(B_1 \cap \ldots \cap B_n) = \mathbb{P}(B_1) \ldots \mathbb{P}(B_n)$ для любых $B_i \in \mathfrak{B}_i, i=\overline{1,n}$. Тогда любая измеримая функция с конечным кодоменом $f : \Omega \rightarrow X$ может быть представлена в виде $f = \varphi \circ \mathfrak{r}$, где $\mathfrak{r} : \Omega \rightarrow [0,1]^n$ такова, что $\mathbb{P} \circ \mathfrak{r}^{-1}$ совпадает с мерой Лебега, а $\varphi : [0,1]^n \rightarrow X$ "--- борелевская.
\end{lemma}

\begin{proof}
	Рассмотрим обратную функцию $f^{-1} : X \rightarrow \mathfrak{B}$. В силу разложимости $\mathfrak{B}$ её можно представить как предел последовательности конъюнкций:
	\begin{equation*}
		f^{-1}(x) = \bigcup_{j=1}^\infty F_1^j(x) \cap \ldots \cap F_n^j(x), \; F_i^j : X \rightarrow \mathfrak{B}_i.
	\end{equation*}
	
	Обозначим семейства множеств $\mathfrak{F}_i = \{F_i^j(x) \mid j \in \mathbb{N}, x \in X\}$ и заметим, что $f$ измерима по $\sigma(\mathfrak{F}_1 \cup \ldots \cup \mathfrak{F}_n)$. По лемме \ref{lemma:card}, существуют цепи множеств $\mathfrak{T}_i \subset \mathfrak{B}_i$ такие, что $\sigma(\mathfrak{F}_i) = \sigma(\mathfrak{T}_i)$. Согласно принципу максимума Хаусдорфа каждая такая цепь вложена в максимальную цепь $\overline{\mathfrak{T}}_i \subset \mathfrak{B}_i$, порождающую безатомическую $\sigma$"~алгебру по лемме \ref{lemma:max}. Построим искомые функции: $\mathfrak{r} = (\operatorname{mim}\langle\overline{\mathfrak{T}}_1, \mathbb{P}\rangle, \ldots, \operatorname{mim}\langle\overline{\mathfrak{T}}_n, \mathbb{P}\rangle)$ и $\varphi = f \circ \mathfrak{r}^{-1}$. Необходимые свойства соблюдаются по построению.
\end{proof}

\begin{proof}[Доказательство теоремы \ref{the:struct}]
	Применим лемму \ref{lemma:consp} к произвольному пространству заговоров $\Phi_1$ структуры $\mathfrak{A} = \{A_1, \ldots, A_n\}$, используя в качестве компонент разложения $\mathfrak{B}_1, \ldots, \mathfrak{B}_n$ $\sigma$"~алгебры тайны соответствующих групп заговорщиков. Это даёт для любого разбиения $f_1 \models \Phi_1$ разложение $f_1 = \varphi \circ \mathfrak{r}$. В любом другом пространстве заговоров $\Phi_2$ той же структуры $\mathfrak{A}$ соответствующее разбиение $f_2 \models \Phi_2$ построим похожим образом: $f_2 = \varphi \circ \mathfrak{u}$. Здесь $\varphi$ то же самое, а $\mathfrak{u} = (\operatorname{mim}\langle\mathfrak{W}_1, \mathbb{P}_2\rangle, \ldots, \operatorname{mim}\langle\mathfrak{W}_n, \mathbb{P}_2\rangle)$, где $\mathfrak{W}_i$ - произвольные максимальные цепи, вложенные в $\sigma$"~алгебры соответствующих тайн пространства заговоров $\Phi_2$. Поскольку и $\mathbb{P}_1 \circ \mathfrak{r}^{-1}$, и $\mathbb{P}_2 \circ \mathfrak{u}^{-1}$ обе совпадают с мерой Лебега, то и $\mathbb{P}_1 \circ f_1^{-1} = \mathbb{P}_2 \circ f_2^{-1}$, а значит теорема доказана.
\end{proof}

\chapter{Примеры вставки листингов программного кода}

Для крупных листингов есть два способа. Первый красивый, но в нём могут быть
проблемы с поддержкой кириллицы (у вас может встречаться в~комментариях
и~печатаемых сообщениях), он представлен на листинге~\cref{lst:hwbeauty}.
\begin{ListingEnv}[!h]% настройки floating аналогичны окружению figure
    \captiondelim{ } % разделитель идентификатора с номером от наименования
    \caption{Программа ,,Hello, world`` на \protect\cpp}\label{lst:hwbeauty}
    % окружение учитывает пробелы и табуляции и применяет их в сответсвии с настройками
    \begin{lstlisting}[language={[ISO]C++}]
	#include <iostream>
	using namespace std;

	int main() //кириллица в комментариях при xelatex и lualatex имеет проблемы с пробелами
	{
		cout << "Hello, world" << endl; //latin letters in commentaries
		system("pause");
		return 0;
	}
    \end{lstlisting}
\end{ListingEnv}%
Второй не~такой красивый, но без ограничений (см.~листинг~\cref{lst:hwplain}).
\begin{ListingEnv}[!h]
    \captiondelim{ } % разделитель идентификатора с номером от наименования
    \caption{Программа ,,Hello, world`` без подсветки}\label{lst:hwplain}
    \begin{Verb}

        #include <iostream>
        using namespace std;

        int main() //кириллица в комментариях
        {
            cout << "Привет, мир" << endl;
        }
    \end{Verb}
\end{ListingEnv}

Можно использовать первый для вставки небольших фрагментов
внутри текста, а второй для вставки полного
кода в приложении, если таковое имеется.

Если нужно вставить совсем короткий пример кода (одна или две строки),
то~выделение  линейками и нумерация может смотреться чересчур громоздко.
В таких случаях можно использовать окружения \texttt{lstlisting} или
\texttt{Verb} без \texttt{ListingEnv}. Приведём такой пример
с указанием языка программирования, отличного от~заданного по умолчанию:
\begin{lstlisting}[language=Haskell]
fibs = 0 : 1 : zipWith (+) fibs (tail fibs)
\end{lstlisting}
Такое решение "--- со вставкой нумерованных листингов покрупнее
и~вставок без выделения для маленьких фрагментов "--- выбрано,
например, в~книге Эндрю Таненбаума и Тодда Остина по архитектуре
компьютера.

Наконец, для оформления идентификаторов внутри строк
(функция \lstinline{main} и~тому подобное) используется
\texttt{lstinline} или, самое простое, моноширинный текст
(\texttt{\textbackslash texttt}).

Пример~\cref{lst:internal3}, иллюстрирующий подключение переопределённого
языка. Может быть полезным, если подсветка кода работает криво. Без
дополнительного окружения, с подписью и ссылкой, реализованной встроенным
средством.
\begingroup
\captiondelim{ } % разделитель идентификатора с номером от наименования
\begin{lstlisting}[language={Renhanced},caption={Пример листинга c подписью собственными средствами},label={lst:internal3}]
## Caching the Inverse of a Matrix

## Matrix inversion is usually a costly computation and there may be some
## benefit to caching the inverse of a matrix rather than compute it repeatedly
## This is a pair of functions that cache the inverse of a matrix.

## makeCacheMatrix creates a special "matrix" object that can cache its inverse

makeCacheMatrix <- function(x = matrix()) {#кириллица в комментариях при xelatex и lualatex имеет проблемы с пробелами
    i <- NULL
    set <- function(y) {
        x <<- y
        i <<- NULL
    }
    get <- function() x
    setSolved <- function(solve) i <<- solve
    getSolved <- function() i
    list(set = set, get = get,
    setSolved = setSolved,
    getSolved = getSolved)

}


## cacheSolve computes the inverse of the special "matrix" returned by
## makeCacheMatrix above. If the inverse has already been calculated (and the
## matrix has not changed), then the cachesolve should retrieve the inverse from
## the cache.

cacheSolve <- function(x, ...) {
    ## Return a matrix that is the inverse of 'x'
    i <- x$getSolved()
    if(!is.null(i)) {
        message("getting cached data")
        return(i)
    }
    data <- x$get()
    i <- solve(data, ...)
    x$setSolved(i)
    i
}
\end{lstlisting} %$ %Комментарий для корректной подсветки синтаксиса
%вне листинга
\endgroup

Листинг~\cref{lst:external1} подгружается из внешнего файла. Приходится
загружать без окружения дополнительного. Иначе по страницам не переносится.
\begingroup
\captiondelim{ } % разделитель идентификатора с номером от наименования
\lstinputlisting[lastline=78,language={R},caption={Листинг из внешнего файла},label={lst:external1}]{listings/run_analysis.R}
\endgroup

\chapter{Очень длинное название второго приложения, в~котором продемонстрирована работа с~длинными таблицами}

\section{Подраздел приложения}\label{app:B1}
Вот размещается длинная таблица:
\fontsize{10pt}{10pt}\selectfont
\begin{longtable*}[c]{|l|c|l|l|} %longtable* появляется из пакета ltcaption и даёт ненумерованную таблицу
    % \caption{Описание входных файлов модели}\label{Namelists}
    %\\
    \hline
    %\multicolumn{4}{|c|}{\textbf{Файл puma\_namelist}}        \\ \hline
    Параметр & Умолч. & Тип & Описание               \\ \hline
    \endfirsthead   \hline
    \multicolumn{4}{|c|}{\small\slshape (продолжение)}        \\ \hline
    Параметр & Умолч. & Тип & Описание               \\ \hline
    \endhead        \hline
    % \multicolumn{4}{|c|}{\small\slshape (окончание)}        \\ \hline
    % Параметр & Умолч. & Тип & Описание               \\ \hline
    %                                             \endlasthead        \hline
    \multicolumn{4}{|r|}{\small\slshape продолжение следует}  \\ \hline
    \endfoot        \hline
    \endlastfoot
    \multicolumn{4}{|l|}{\&INP}        \\ \hline
    kick & 1 & int & 0: инициализация без шума (\(p_s = const\)) \\
    &   &     & 1: генерация белого шума                  \\
    &   &     & 2: генерация белого шума симметрично относительно \\
    & & & экватора    \\
    mars & 0 & int & 1: инициализация модели для планеты Марс     \\
    kick & 1 & int & 0: инициализация без шума (\(p_s = const\)) \\
    &   &     & 1: генерация белого шума                  \\
    &   &     & 2: генерация белого шума симметрично относительно \\
    & & & экватора    \\
    mars & 0 & int & 1: инициализация модели для планеты Марс     \\
    kick & 1 & int & 0: инициализация без шума (\(p_s = const\)) \\
    &   &     & 1: генерация белого шума                  \\
    &   &     & 2: генерация белого шума симметрично относительно \\
    & & & экватора    \\
    mars & 0 & int & 1: инициализация модели для планеты Марс     \\
    kick & 1 & int & 0: инициализация без шума (\(p_s = const\)) \\
    &   &     & 1: генерация белого шума                  \\
    &   &     & 2: генерация белого шума симметрично относительно \\
    & & & экватора    \\
    mars & 0 & int & 1: инициализация модели для планеты Марс     \\
    kick & 1 & int & 0: инициализация без шума (\(p_s = const\)) \\
    &   &     & 1: генерация белого шума                  \\
    &   &     & 2: генерация белого шума симметрично относительно \\
    & & & экватора    \\
    mars & 0 & int & 1: инициализация модели для планеты Марс     \\
    kick & 1 & int & 0: инициализация без шума (\(p_s = const\)) \\
    &   &     & 1: генерация белого шума                  \\
    &   &     & 2: генерация белого шума симметрично относительно \\
    & & & экватора    \\
    mars & 0 & int & 1: инициализация модели для планеты Марс     \\
    kick & 1 & int & 0: инициализация без шума (\(p_s = const\)) \\
    &   &     & 1: генерация белого шума                  \\
    &   &     & 2: генерация белого шума симметрично относительно \\
    & & & экватора    \\
    mars & 0 & int & 1: инициализация модели для планеты Марс     \\
    kick & 1 & int & 0: инициализация без шума (\(p_s = const\)) \\
    &   &     & 1: генерация белого шума                  \\
    &   &     & 2: генерация белого шума симметрично относительно \\
    & & & экватора    \\
    mars & 0 & int & 1: инициализация модели для планеты Марс     \\
    kick & 1 & int & 0: инициализация без шума (\(p_s = const\)) \\
    &   &     & 1: генерация белого шума                  \\
    &   &     & 2: генерация белого шума симметрично относительно \\
    & & & экватора    \\
    mars & 0 & int & 1: инициализация модели для планеты Марс     \\
    kick & 1 & int & 0: инициализация без шума (\(p_s = const\)) \\
    &   &     & 1: генерация белого шума                  \\
    &   &     & 2: генерация белого шума симметрично относительно \\
    & & & экватора    \\
    mars & 0 & int & 1: инициализация модели для планеты Марс     \\
    kick & 1 & int & 0: инициализация без шума (\(p_s = const\)) \\
    &   &     & 1: генерация белого шума                  \\
    &   &     & 2: генерация белого шума симметрично относительно \\
    & & & экватора    \\
    mars & 0 & int & 1: инициализация модели для планеты Марс     \\
    kick & 1 & int & 0: инициализация без шума (\(p_s = const\)) \\
    &   &     & 1: генерация белого шума                  \\
    &   &     & 2: генерация белого шума симметрично относительно \\
    & & & экватора    \\
    mars & 0 & int & 1: инициализация модели для планеты Марс     \\
    kick & 1 & int & 0: инициализация без шума (\(p_s = const\)) \\
    &   &     & 1: генерация белого шума                  \\
    &   &     & 2: генерация белого шума симметрично относительно \\
    & & & экватора    \\
    mars & 0 & int & 1: инициализация модели для планеты Марс     \\
    kick & 1 & int & 0: инициализация без шума (\(p_s = const\)) \\
    &   &     & 1: генерация белого шума                  \\
    &   &     & 2: генерация белого шума симметрично относительно \\
    & & & экватора    \\
    mars & 0 & int & 1: инициализация модели для планеты Марс     \\
    kick & 1 & int & 0: инициализация без шума (\(p_s = const\)) \\
    &   &     & 1: генерация белого шума                  \\
    &   &     & 2: генерация белого шума симметрично относительно \\
    & & & экватора    \\
    mars & 0 & int & 1: инициализация модели для планеты Марс     \\
    \hline
    %& & & \(\:\) \\
    \multicolumn{4}{|l|}{\&SURFPAR}        \\ \hline
    kick & 1 & int & 0: инициализация без шума (\(p_s = const\)) \\
    &   &     & 1: генерация белого шума                  \\
    &   &     & 2: генерация белого шума симметрично относительно \\
    & & & экватора    \\
    mars & 0 & int & 1: инициализация модели для планеты Марс     \\
    kick & 1 & int & 0: инициализация без шума (\(p_s = const\)) \\
    &   &     & 1: генерация белого шума                  \\
    &   &     & 2: генерация белого шума симметрично относительно \\
    & & & экватора    \\
    mars & 0 & int & 1: инициализация модели для планеты Марс     \\
    kick & 1 & int & 0: инициализация без шума (\(p_s = const\)) \\
    &   &     & 1: генерация белого шума                  \\
    &   &     & 2: генерация белого шума симметрично относительно \\
    & & & экватора    \\
    mars & 0 & int & 1: инициализация модели для планеты Марс     \\
    kick & 1 & int & 0: инициализация без шума (\(p_s = const\)) \\
    &   &     & 1: генерация белого шума                  \\
    &   &     & 2: генерация белого шума симметрично относительно \\
    & & & экватора    \\
    mars & 0 & int & 1: инициализация модели для планеты Марс     \\
    kick & 1 & int & 0: инициализация без шума (\(p_s = const\)) \\
    &   &     & 1: генерация белого шума                  \\
    &   &     & 2: генерация белого шума симметрично относительно \\
    & & & экватора    \\
    mars & 0 & int & 1: инициализация модели для планеты Марс     \\
    kick & 1 & int & 0: инициализация без шума (\(p_s = const\)) \\
    &   &     & 1: генерация белого шума                  \\
    &   &     & 2: генерация белого шума симметрично относительно \\
    & & & экватора    \\
    mars & 0 & int & 1: инициализация модели для планеты Марс     \\
    kick & 1 & int & 0: инициализация без шума (\(p_s = const\)) \\
    &   &     & 1: генерация белого шума                  \\
    &   &     & 2: генерация белого шума симметрично относительно \\
    & & & экватора    \\
    mars & 0 & int & 1: инициализация модели для планеты Марс     \\
    kick & 1 & int & 0: инициализация без шума (\(p_s = const\)) \\
    &   &     & 1: генерация белого шума                  \\
    &   &     & 2: генерация белого шума симметрично относительно \\
    & & & экватора    \\
    mars & 0 & int & 1: инициализация модели для планеты Марс     \\
    kick & 1 & int & 0: инициализация без шума (\(p_s = const\)) \\
    &   &     & 1: генерация белого шума                  \\
    &   &     & 2: генерация белого шума симметрично относительно \\
    & & & экватора    \\
    mars & 0 & int & 1: инициализация модели для планеты Марс     \\
    \hline
\end{longtable*}

\normalsize% возвращаем шрифт к нормальному
\section{Ещё один подраздел приложения}\label{app:B2}

Нужно больше подразделов приложения!
Конвынёры витюпырата но нам, тебиквюэ мэнтётюм позтюлант ед про. Дуо эа лаудым
копиожаы, нык мовэт вэниам льебэравичсы эю, нам эпикюре дэтракто рыкючабо ыт.

Пример длинной таблицы с записью продолжения по ГОСТ 2.105:

\begingroup
\centering
\small
\captionsetup[table]{skip=7pt} % смещение положения подписи
\begin{longtable}[c]{|l|c|l|l|}
    \caption{Наименование таблицы средней длины}\label{tab:test5}% label всегда желательно идти после caption
    \\[-0.45\onelineskip]
    \hline
    Параметр & Умолч. & Тип & Описание                                          \\ \hline
    \endfirsthead%
    \caption*{Продолжение таблицы~\thetable}                                    \\[-0.45\onelineskip]
    \hline
    Параметр & Умолч. & Тип & Описание                                          \\ \hline
    \endhead
    \hline
    \endfoot
    \hline
    \endlastfoot
    \multicolumn{4}{|l|}{\&INP}                                                 \\ \hline
    kick     & 1      & int & 0: инициализация без шума (\(p_s = const\))       \\
             &        &     & 1: генерация белого шума                          \\
             &        &     & 2: генерация белого шума симметрично относительно \\
             &        &     & экватора                                          \\
    mars     & 0      & int & 1: инициализация модели для планеты Марс          \\
    kick     & 1      & int & 0: инициализация без шума (\(p_s = const\))       \\
             &        &     & 1: генерация белого шума                          \\
             &        &     & 2: генерация белого шума симметрично относительно \\
             &        &     & экватора                                          \\
    mars     & 0      & int & 1: инициализация модели для планеты Марс          \\
    kick     & 1      & int & 0: инициализация без шума (\(p_s = const\))       \\
             &        &     & 1: генерация белого шума                          \\
             &        &     & 2: генерация белого шума симметрично относительно \\
             &        &     & экватора                                          \\
    mars     & 0      & int & 1: инициализация модели для планеты Марс          \\
    kick     & 1      & int & 0: инициализация без шума (\(p_s = const\))       \\
             &        &     & 1: генерация белого шума                          \\
             &        &     & 2: генерация белого шума симметрично относительно \\
             &        &     & экватора                                          \\
    mars     & 0      & int & 1: инициализация модели для планеты Марс          \\
    kick     & 1      & int & 0: инициализация без шума (\(p_s = const\))       \\
             &        &     & 1: генерация белого шума                          \\
             &        &     & 2: генерация белого шума симметрично относительно \\
             &        &     & экватора                                          \\
    mars     & 0      & int & 1: инициализация модели для планеты Марс          \\
    kick     & 1      & int & 0: инициализация без шума (\(p_s = const\))       \\
             &        &     & 1: генерация белого шума                          \\
             &        &     & 2: генерация белого шума симметрично относительно \\
             &        &     & экватора                                          \\
    mars     & 0      & int & 1: инициализация модели для планеты Марс          \\
    kick     & 1      & int & 0: инициализация без шума (\(p_s = const\))       \\
             &        &     & 1: генерация белого шума                          \\
             &        &     & 2: генерация белого шума симметрично относительно \\
             &        &     & экватора                                          \\
    mars     & 0      & int & 1: инициализация модели для планеты Марс          \\
    kick     & 1      & int & 0: инициализация без шума (\(p_s = const\))       \\
             &        &     & 1: генерация белого шума                          \\
             &        &     & 2: генерация белого шума симметрично относительно \\
             &        &     & экватора                                          \\
    mars     & 0      & int & 1: инициализация модели для планеты Марс          \\
    kick     & 1      & int & 0: инициализация без шума (\(p_s = const\))       \\
             &        &     & 1: генерация белого шума                          \\
             &        &     & 2: генерация белого шума симметрично относительно \\
             &        &     & экватора                                          \\
    mars     & 0      & int & 1: инициализация модели для планеты Марс          \\
    kick     & 1      & int & 0: инициализация без шума (\(p_s = const\))       \\
             &        &     & 1: генерация белого шума                          \\
             &        &     & 2: генерация белого шума симметрично относительно \\
             &        &     & экватора                                          \\
    mars     & 0      & int & 1: инициализация модели для планеты Марс          \\
    kick     & 1      & int & 0: инициализация без шума (\(p_s = const\))       \\
             &        &     & 1: генерация белого шума                          \\
             &        &     & 2: генерация белого шума симметрично относительно \\
             &        &     & экватора                                          \\
    mars     & 0      & int & 1: инициализация модели для планеты Марс          \\
    kick     & 1      & int & 0: инициализация без шума (\(p_s = const\))       \\
             &        &     & 1: генерация белого шума                          \\
             &        &     & 2: генерация белого шума симметрично относительно \\
             &        &     & экватора                                          \\
    mars     & 0      & int & 1: инициализация модели для планеты Марс          \\
    kick     & 1      & int & 0: инициализация без шума (\(p_s = const\))       \\
             &        &     & 1: генерация белого шума                          \\
             &        &     & 2: генерация белого шума симметрично относительно \\
             &        &     & экватора                                          \\
    mars     & 0      & int & 1: инициализация модели для планеты Марс          \\
    kick     & 1      & int & 0: инициализация без шума (\(p_s = const\))       \\
             &        &     & 1: генерация белого шума                          \\
             &        &     & 2: генерация белого шума симметрично относительно \\
             &        &     & экватора                                          \\
    mars     & 0      & int & 1: инициализация модели для планеты Марс          \\
    kick     & 1      & int & 0: инициализация без шума (\(p_s = const\))       \\
             &        &     & 1: генерация белого шума                          \\
             &        &     & 2: генерация белого шума симметрично относительно \\
             &        &     & экватора                                          \\
    mars     & 0      & int & 1: инициализация модели для планеты Марс          \\
    \hline
    %& & & $\:$ \\
    \multicolumn{4}{|l|}{\&SURFPAR}                                             \\ \hline
    kick     & 1      & int & 0: инициализация без шума (\(p_s = const\))       \\
             &        &     & 1: генерация белого шума                          \\
             &        &     & 2: генерация белого шума симметрично относительно \\
             &        &     & экватора                                          \\
    mars     & 0      & int & 1: инициализация модели для планеты Марс          \\
    kick     & 1      & int & 0: инициализация без шума (\(p_s = const\))       \\
             &        &     & 1: генерация белого шума                          \\
             &        &     & 2: генерация белого шума симметрично относительно \\
             &        &     & экватора                                          \\
    mars     & 0      & int & 1: инициализация модели для планеты Марс          \\
    kick     & 1      & int & 0: инициализация без шума (\(p_s = const\))       \\
             &        &     & 1: генерация белого шума                          \\
             &        &     & 2: генерация белого шума симметрично относительно \\
             &        &     & экватора                                          \\
    mars     & 0      & int & 1: инициализация модели для планеты Марс          \\
    kick     & 1      & int & 0: инициализация без шума (\(p_s = const\))       \\
             &        &     & 1: генерация белого шума                          \\
             &        &     & 2: генерация белого шума симметрично относительно \\
             &        &     & экватора                                          \\
    mars     & 0      & int & 1: инициализация модели для планеты Марс          \\
    kick     & 1      & int & 0: инициализация без шума (\(p_s = const\))       \\
             &        &     & 1: генерация белого шума                          \\
             &        &     & 2: генерация белого шума симметрично относительно \\
             &        &     & экватора                                          \\
    mars     & 0      & int & 1: инициализация модели для планеты Марс          \\
    kick     & 1      & int & 0: инициализация без шума (\(p_s = const\))       \\
             &        &     & 1: генерация белого шума                          \\
             &        &     & 2: генерация белого шума симметрично относительно \\
             &        &     & экватора                                          \\
    mars     & 0      & int & 1: инициализация модели для планеты Марс          \\
    kick     & 1      & int & 0: инициализация без шума (\(p_s = const\))       \\
             &        &     & 1: генерация белого шума                          \\
             &        &     & 2: генерация белого шума симметрично относительно \\
             &        &     & экватора                                          \\
    mars     & 0      & int & 1: инициализация модели для планеты Марс          \\
    kick     & 1      & int & 0: инициализация без шума (\(p_s = const\))       \\
             &        &     & 1: генерация белого шума                          \\
             &        &     & 2: генерация белого шума симметрично относительно \\
             &        &     & экватора                                          \\
    mars     & 0      & int & 1: инициализация модели для планеты Марс          \\
    kick     & 1      & int & 0: инициализация без шума (\(p_s = const\))       \\
             &        &     & 1: генерация белого шума                          \\
             &        &     & 2: генерация белого шума симметрично относительно \\
             &        &     & экватора                                          \\
    mars     & 0      & int & 1: инициализация модели для планеты Марс          \\
\end{longtable}
\normalsize% возвращаем шрифт к нормальному
\endgroup
\section{Использование длинных таблиц с окружением \textit{longtabu}}\label{app:B2a}

В таблице \cref{tab:test-functions} более книжный вариант
длинной таблицы, используя окружение \verb!longtabu! и разнообразные
\verb!toprule! \verb!midrule! \verb!bottomrule! из~пакета
\verb!booktabs!. Чтобы визуально таблица смотрелась лучше, можно
использовать следующие параметры: в самом начале задаётся расстояние
между строчками с~помощью \verb!arraystretch!. Таблица задаётся на
всю ширину, \verb!longtabu! позволяет делить ширину колонок
пропорционально "--- тут три колонки в~пропорции 1.1:1:4 "--- для каждой
колонки первый параметр в~описании \verb!X[]!. Кроме того, в~таблице
убраны отступы слева и справа с~помощью \verb!@{}!
в~преамбуле таблицы. К~первому и~второму столбцу применяется
модификатор

\verb!>{\setlength{\baselineskip}{0.7\baselineskip}}!,

\noindent который уменьшает межстрочный интервал в для текста таблиц (иначе
заголовок второго столбца значительно шире, а двухстрочное имя
сливается с~окружающими). Для первой и второй колонки текст в ячейках
выравниваются по~центру как по~вертикали, так и по горизонтали "---
задаётся буквами \verb!m!~и~\verb!c!~в~описании столбца \verb!X[]!.

Так как формулы большие "--- используется окружение \verb!alignedat!,
чтобы отступ был одинаковый у всех формул "--- он сделан для всех, хотя
для большей части можно было и не использовать.  Чтобы формулы
занимали поменьше места в~каждом столбце формулы (где надо)
используется \verb!\textstyle! "--- он~делает дроби меньше, у~знаков
суммы и произведения "--- индексы сбоку. Иногда формула слишком большая,
сливается со следующей, поэтому после неё ставится небольшой
дополнительный отступ \verb!\vspace*{2ex}!. Для штрафных функций "---
размер фигурных скобок задан вручную \verb!\Big\{!, т.\:к. не~умеет
\verb!alignedat! работать с~\verb!\left! и~\verb!\right! через
несколько строк/колонок.

В примечании к таблице наоборот, окружение \verb!cases! даёт слишком
большие промежутки между вариантами, чтобы их уменьшить, в конце
каждой строчки окружения использовался отрицательный дополнительный
отступ \verb!\\[-0.5em]!.

\begingroup % Ограничиваем область видимости arraystretch
\renewcommand{\arraystretch}{1.6}%% Увеличение расстояния между рядами, для улучшения восприятия.
\begin{longtabu} to \textwidth
    {%
    @{}>{\setlength{\baselineskip}{0.7\baselineskip}}X[1.1mc]%
    >{\setlength{\baselineskip}{0.7\baselineskip}}X[1.1mc]%
    X[4]@{}%
    }
    \caption{Тестовые функции для оптимизации, \(D\) "---
        размерность. Для всех функций значение в точке глобального
        минимума равно нулю.\label{tab:test-functions}}\\% label всегда желательно идти после caption

    \toprule     %%% верхняя линейка
    Имя           &Стартовый диапазон параметров &Функция  \\
    \midrule %%% тонкий разделитель. Отделяет названия столбцов. Обязателен по ГОСТ 2.105 пункт 4.4.5
    \endfirsthead

    \multicolumn{3}{c}{\small\slshape (продолжение)}        \\
    \toprule     %%% верхняя линейка
    Имя           &Стартовый диапазон параметров &Функция  \\
    \midrule %%% тонкий разделитель. Отделяет названия столбцов. Обязателен по ГОСТ 2.105 пункт 4.4.5
    \endhead

    \multicolumn{3}{c}{\small\slshape (окончание)}        \\
    \toprule     %%% верхняя линейка
    Имя           &Стартовый диапазон параметров &Функция  \\
    \midrule %%% тонкий разделитель. Отделяет названия столбцов. Обязателен по ГОСТ 2.105 пункт 4.4.5
    \endlasthead

    \bottomrule %%% нижняя линейка
    \multicolumn{3}{r}{\small\slshape продолжение следует}  \\
    \endfoot
    \endlastfoot

    сфера         &\(\left[-100,\,100\right]^D\)   &
    \(\begin{aligned}
        \textstyle f_1(x)=\sum_{i=1}^Dx_i^2
    \end{aligned}\) \\
    Schwefel 2.22 &\(\left[-10,\,10\right]^D\)     &
    \(\begin{aligned}
        \textstyle f_2(x)=\sum_{i=1}^D|x_i|+\prod_{i=1}^D|x_i|
    \end{aligned}\) \\
    Schwefel 1.2  &\(\left[-100,\,100\right]^D\)   &
    \(\begin{aligned}
        \textstyle f_3(x)=\sum_{i=1}^D\left(\sum_{j=1}^ix_j\right)^2
    \end{aligned}\) \\
    Schwefel 2.21 &\(\left[-100,\,100\right]^D\)   &
    \(\begin{aligned}
        \textstyle f_4(x)=\max_i\!\left\{\left|x_i\right|\right\}
    \end{aligned}\) \\
    Rosenbrock    &\(\left[-30,\,30\right]^D\)     &
    \(\begin{aligned}
        \textstyle f_5(x)=
        \sum_{i=1}^{D-1}
        \left[100\!\left(x_{i+1}-x_i^2\right)^2+(x_i-1)^2\right]
    \end{aligned}\) \\
    ступенчатая   &\(\left[-100,\,100\right]^D\)   &
    \(\begin{aligned}
        \textstyle f_6(x)=\sum_{i=1}^D\big\lfloor x_i+0.5\big\rfloor^2
    \end{aligned}\) \\
    зашумлённая квартическая &\(\left[-1.28,\,1.28\right]^D\) &
    \(\begin{aligned}
        \textstyle f_7(x)=\sum_{i=1}^Dix_i^4+rand[0,1)
    \end{aligned}\)\vspace*{2ex}\\
    Schwefel 2.26 &\(\left[-500,\,500\right]^D\)   &
    \(\begin{aligned}
        f_8(x)= & \textstyle\sum_{i=1}^D-x_i\,\sin\sqrt{|x_i|}\,+ \\
                & \vphantom{\sum}+ D\cdot
        418.98288727243369
    \end{aligned}\)\\
    Rastrigin     &\(\left[-5.12,\,5.12\right]^D\) &
    \(\begin{aligned}
        \textstyle f_9(x)=\sum_{i=1}^D\left[x_i^2-10\,\cos(2\pi x_i)+10\right]
    \end{aligned}\)\vspace*{2ex}\\
    Ackley        &\(\left[-32,\,32\right]^D\)     &
    \(\begin{aligned}
        f_{10}(x)= & \textstyle -20\, \exp\!\left(
        -0.2\sqrt{\frac{1}{D}\sum_{i=1}^Dx_i^2} \right)- \\
                   & \textstyle - \exp\left(
            \frac{1}{D}\sum_{i=1}^D\cos(2\pi x_i)  \right)
        + 20 + e
    \end{aligned}\) \\
    Griewank      &\(\left[-600,\,600\right]^D\) &
    \(\begin{aligned}
        f_{11}(x)= & \textstyle \frac{1}{4000}\sum_{i=1}^{D}x_i^2 -
        \prod_{i=1}^D\cos\left(x_i/\sqrt{i}\right) +1
    \end{aligned}\) \vspace*{3ex} \\
    штрафная 1    &\(\left[-50,\,50\right]^D\)     &
    \(\begin{aligned}
        f_{12}(x)= & \textstyle \frac{\pi}{D}\Big\{ 10\,\sin^2(\pi y_1) +            \\
                   & +\textstyle \sum_{i=1}^{D-1}(y_i-1)^2
        \left[1+10\,\sin^2(\pi y_{i+1})\right] +                                     \\
                   & +(y_D-1)^2 \Big\} +\textstyle\sum_{i=1}^D u(x_i,\,10,\,100,\,4)
    \end{aligned}\) \vspace*{2ex} \\
    штрафная 2    &\(\left[-50,\,50\right]^D\)     &
    \(\begin{aligned}
        f_{13}(x)= & \textstyle 0.1 \Big\{\sin^2(3\pi x_1) +            \\
                   & +\textstyle \sum_{i=1}^{D-1}(x_i-1)^2
        \left[1+\sin^2(3 \pi x_{i+1})\right] +                          \\
                   & +(x_D-1)^2\left[1+\sin^2(2\pi x_D)\right] \Big\} + \\
                   & +\textstyle\sum_{i=1}^D u(x_i,\,5,\,100,\,4)
    \end{aligned}\)\\
    сфера         &\(\left[-100,\,100\right]^D\)   &
    \(\begin{aligned}
        \textstyle f_1(x)=\sum_{i=1}^Dx_i^2
    \end{aligned}\) \\
    Schwefel 2.22 &\(\left[-10,\,10\right]^D\)     &
    \(\begin{aligned}
        \textstyle f_2(x)=\sum_{i=1}^D|x_i|+\prod_{i=1}^D|x_i|
    \end{aligned}\) \\
    Schwefel 1.2  &\(\left[-100,\,100\right]^D\)   &
    \(\begin{aligned}
        \textstyle f_3(x)=\sum_{i=1}^D\left(\sum_{j=1}^ix_j\right)^2
    \end{aligned}\) \\
    Schwefel 2.21 &\(\left[-100,\,100\right]^D\)   &
    \(\begin{aligned}
        \textstyle f_4(x)=\max_i\!\left\{\left|x_i\right|\right\}
    \end{aligned}\) \\
    Rosenbrock    &\(\left[-30,\,30\right]^D\)     &
    \(\begin{aligned}
        \textstyle f_5(x)=
        \sum_{i=1}^{D-1}
        \left[100\!\left(x_{i+1}-x_i^2\right)^2+(x_i-1)^2\right]
    \end{aligned}\) \\
    ступенчатая   &\(\left[-100,\,100\right]^D\)   &
    \(\begin{aligned}
        \textstyle f_6(x)=\sum_{i=1}^D\big\lfloor x_i+0.5\big\rfloor^2
    \end{aligned}\) \\
    зашумлённая квартическая &\(\left[-1.28,\,1.28\right]^D\) &
    \(\begin{aligned}
        \textstyle f_7(x)=\sum_{i=1}^Dix_i^4+rand[0,1)
    \end{aligned}\)\vspace*{2ex}\\
    Schwefel 2.26 &\(\left[-500,\,500\right]^D\)   &
    \(\begin{aligned}
        f_8(x)= & \textstyle\sum_{i=1}^D-x_i\,\sin\sqrt{|x_i|}\,+ \\
                & \vphantom{\sum}+ D\cdot
        418.98288727243369
    \end{aligned}\)\\
    Rastrigin     &\(\left[-5.12,\,5.12\right]^D\) &
    \(\begin{aligned}
        \textstyle f_9(x)=\sum_{i=1}^D\left[x_i^2-10\,\cos(2\pi x_i)+10\right]
    \end{aligned}\)\vspace*{2ex}\\
    Ackley        &\(\left[-32,\,32\right]^D\)     &
    \(\begin{aligned}
        f_{10}(x)= & \textstyle -20\, \exp\!\left(
        -0.2\sqrt{\frac{1}{D}\sum_{i=1}^Dx_i^2} \right)- \\
                   & \textstyle - \exp\left(
            \frac{1}{D}\sum_{i=1}^D\cos(2\pi x_i)  \right)
        + 20 + e
    \end{aligned}\) \\
    Griewank      &\(\left[-600,\,600\right]^D\) &
    \(\begin{aligned}
        f_{11}(x)= & \textstyle \frac{1}{4000}\sum_{i=1}^{D}x_i^2 -
        \prod_{i=1}^D\cos\left(x_i/\sqrt{i}\right) +1
    \end{aligned}\) \vspace*{3ex} \\
    штрафная 1    &\(\left[-50,\,50\right]^D\)     &
    \(\begin{aligned}
        f_{12}(x)= & \textstyle \frac{\pi}{D}\Big\{ 10\,\sin^2(\pi y_1) +            \\
                   & +\textstyle \sum_{i=1}^{D-1}(y_i-1)^2
        \left[1+10\,\sin^2(\pi y_{i+1})\right] +                                     \\
                   & +(y_D-1)^2 \Big\} +\textstyle\sum_{i=1}^D u(x_i,\,10,\,100,\,4)
    \end{aligned}\) \vspace*{2ex} \\
    штрафная 2    &\(\left[-50,\,50\right]^D\)     &
    \(\begin{aligned}
        f_{13}(x)= & \textstyle 0.1 \Big\{\sin^2(3\pi x_1) +            \\
                   & +\textstyle \sum_{i=1}^{D-1}(x_i-1)^2
        \left[1+\sin^2(3 \pi x_{i+1})\right] +                          \\
                   & +(x_D-1)^2\left[1+\sin^2(2\pi x_D)\right] \Big\} + \\
                   & +\textstyle\sum_{i=1}^D u(x_i,\,5,\,100,\,4)
    \end{aligned}\)\\
    \midrule%%% тонкий разделитель
    \multicolumn{3}{@{}p{\textwidth}}{%
    \vspace*{-3.5ex}% этим подтягиваем повыше
    \hspace*{2.5em}% абзацный отступ - требование ГОСТ 2.105
    Примечание "---  Для функций \(f_{12}\) и \(f_{13}\)
    используется \(y_i = 1 + \frac{1}{4}(x_i+1)\)
    и~$u(x_i,\,a,\,k,\,m)=
        \begin{cases*}
            k(x_i-a)^m,  & \( x_i >a \)            \\[-0.5em]
            0,           & \( -a\leq x_i \leq a \) \\[-0.5em]
            k(-x_i-a)^m, & \( x_i <-a \)
        \end{cases*}
    $
    }\\
    \bottomrule %%% нижняя линейка
\end{longtabu}
\endgroup

\section{Форматирование внутри таблиц}\label{app:B3}

В таблице \cref{tab:other-row} пример с чересстрочным
форматированием. В~файле \verb+userstyles.tex+  задаётся счётчик
\verb+\newcounter{rowcnt}+ который увеличивается на~1 после каждой
строчки (как указано в преамбуле таблицы). Кроме того, задаётся
условный макрос \verb+\altshape+ который выдаёт одно
из~двух типов форматирования в~зависимости от чётности счётчика.

В таблице \cref{tab:other-row} каждая чётная строчка "--- синяя,
нечётная "--- с наклоном и~слегка поднята вверх. Визуально это приводит
к тому, что среднее значение и~среднеквадратичное изменение
группируются и хорошо выделяются взглядом в~таблице. Сохраняется
возможность отдельные значения в таблице выделить цветом или
шрифтом. К первому и второму столбцу форматирование не применяется
по~сути таблицы, к шестому общее форматирование не~применяется для
наглядности.

Так как заголовок таблицы тоже считается за строчку, то перед ним (для
первого, промежуточного и финального варианта) счётчик обнуляется,
а~в~\verb+\altshape+ для нулевого значения счётчика форматирования
не~применяется.

\begingroup % Ограничиваем область видимости arraystretch
\renewcommand\altshape{
    \ifnumequal{\value{rowcnt}}{0}{
        % Стиль для заголовка таблицы
    }{
        \ifnumodd{\value{rowcnt}}
        {
            \color{blue} % Cтиль для нечётных строк
        }{
            \vspace*{-0.7ex}\itshape} % Стиль для чётных строк
    }
}
\newcolumntype{A}{>{\centering\begingroup\altshape}X[1mc]<{\endgroup}}
\needspace{2\baselineskip}
\renewcommand{\arraystretch}{0.9}%% Уменьшаем  расстояние между
%% рядами, чтобы таблица не так много
%% места занимала в дисере.
\begin{longtabu} to \textwidth {@{}X[0.27ml]@{}X[0.7mc]@{}A@{}A@{}A@{}X[0.98mc]@{}>{\setlength{\baselineskip}{0.7\baselineskip}}A@{}A<{\stepcounter{rowcnt}}@{}}
    % \begin{longtabu} to \textwidth {@{}X[0.2ml]X[1mc]X[1mc]X[1mc]X[1mc]X[1mc]>{\setlength{\baselineskip}{0.7\baselineskip}}X[1mc]X[1mc]@{}}
    \caption{Длинная таблица с примером чересстрочного форматирования\label{tab:other-row}}\vspace*{1ex}\\% label всегда желательно идти после caption
    % \vspace*{1ex}     \\

    \toprule %%% верхняя линейка
    \setcounter{rowcnt}{0} &Итера\-ции & JADE\texttt{++} & JADE & jDE & SaDE
    & DE/rand /1/bin & PSO \\
    \midrule %%% тонкий разделитель. Отделяет названия столбцов. Обязателен по ГОСТ 2.105 пункт 4.4.5
    \endfirsthead

    \multicolumn{8}{c}{\small\slshape (продолжение)} \\
    \toprule %%% верхняя линейка
    \setcounter{rowcnt}{0} &Итера\-ции & JADE\texttt{++} & JADE & jDE & SaDE
    & DE/rand /1/bin & PSO \\
    \midrule %%% тонкий разделитель. Отделяет названия столбцов. Обязателен по ГОСТ 2.105 пункт 4.4.5
    \endhead

    \multicolumn{8}{c}{\small\slshape (окончание)} \\
    \toprule %%% верхняя линейка
    \setcounter{rowcnt}{0} &Итера\-ции & JADE\texttt{++} & JADE & jDE & SaDE
    & DE/rand /1/bin & PSO \\
    \midrule %%% тонкий разделитель. Отделяет названия столбцов. Обязателен по ГОСТ 2.105 пункт 4.4.5
    \endlasthead

    \bottomrule %%% нижняя линейка
    \multicolumn{8}{r}{\small\slshape продолжение следует}     \\
    \endfoot
    \endlastfoot

    f1  & 1500 & \textbf{1.8E-60}   & 1.3E-54   & 2.5E-28   & 4.5E-20   & 9.8E-14   & 9.6E-42   \\\nopagebreak
    &      & (8.4E-60) & (9.2E-54) & {\color{red}(3.5E-28)} & (6.9E-20) & (8.4E-14) & (2.7E-41) \\
    f2  & 2000 & 1.8E-25   & 3.9E-22   & 1.5E-23   & 1.9E-14   & 1.6E-09   & 9.3E-21   \\\nopagebreak
    &      & (8.8E-25) & (2.7E-21) & (1.0E-23) & (1.1E-14) & (1.1E-09) & (6.3E-20) \\
    f3  & 5000 & 5.7E-61   & 6.0E-87   & 5.2E-14   & {\color{green}9.0E-37}   & 6.6E-11   & 2.5E-19   \\\nopagebreak
    &      & (2.7E-60) & (1.9E-86) & (1.1E-13) & (5.4E-36) & (8.8E-11) & (3.9E-19) \\
    f4  & 5000 & 8.2E-24   & 4.3E-66   & 1.4E-15   & 7.4E-11   & 4.2E-01   & 4.4E-14   \\\nopagebreak
    &      & (4.0E-23) & (1.2E-65) & (1.0E-15) & (1.8E-10) & (1.1E+00) & (9.3E-14) \\
    f5  & 3000 & 8.0E-02   & 3.2E-01   & 1.3E+01   & 2.1E+01   & 2.1E+00   & 2.5E+01   \\\nopagebreak
    &      & (5.6E-01) & (1.1E+00) & (1.4E+01) & (7.8E+00) & (1.5E+00) & (3.2E+01) \\
    f6  & 100  & 2.9E+00   & 5.6E+00   & 1.0E+03   & 9.3E+02   & 4.7E+03   & 4.5E+01   \\\nopagebreak
    &      & (1.2E+00) & (1.6E+00) & (2.2E+02) & (1.8E+02) & (1.1E+03) & (2.4E+01) \\
    f7  & 3000 & 6.4E-04   & 6.8E-04   & 3.3E-03   & 4.8E-03   & 4.7E-03   & 2.5E-03   \\\nopagebreak
    &      & (2.5E-04) & (2.5E-04) & (8.5E-04) & (1.2E-03) & (1.2E-03) & (1.4E-03) \\
    f8  & 1000 & 3.3E-05   & 7.1E+00   & 7.9E-11   & 4.7E+00   & 5.9E+03   & 2.4E+03   \\\nopagebreak
    &      & (2.3E-05) & (2.8E+01) & (1.3E-10) & (3.3E+01) & (1.1E+03) & (6.7E+02) \\
    f9  & 1000 & 1.0E-04   & 1.4E-04   & 1.5E-04   & 1.2E-03   & 1.8E+02   & 5.2E+01   \\\nopagebreak
    &      & (6.0E-05) & (6.5E-05) & (2.0E-04) & (6.5E-04) & (1.3E+01) & (1.6E+01) \\
    f10 & 500  & 8.2E-10   & 3.0E-09   & 3.5E-04   & 2.7E-03   & 1.1E-01   & 4.6E-01   \\\nopagebreak
    &      & (6.9E-10) & (2.2E-09) & (1.0E-04) & (5.1E-04) & (3.9E-02) & (6.6E-01) \\
    f11 & 500  & 9.9E-08   & 2.0E-04   & 1.9E-05   & 7.8E-04  & 2.0E-01   & 1.3E-02   \\\nopagebreak
    &      & (6.0E-07) & (1.4E-03) & (5.8E-05) & (1.2E-03)  & (1.1E-01) & (1.7E-02) \\
    f12 & 500  & 4.6E-17   & 3.8E-16   & 1.6E-07   & 1.9E-05   & 1.2E-02   & 1.9E-01   \\\nopagebreak
    &      & (1.9E-16) & (8.3E-16) & (1.5E-07) & (9.2E-06) & (1.0E-02) & (3.9E-01) \\
    f13 & 500  & 2.0E-16   & 1.2E-15   & 1.5E-06   & 6.1E-05   & 7.5E-02   & 2.9E-03   \\\nopagebreak
    &      & (6.5E-16) & (2.8E-15) & (9.8E-07) & (2.0E-05) & (3.8E-02) & (4.8E-03) \\
    f1  & 1500 & \textbf{1.8E-60}   & 1.3E-54   & 2.5E-28   & 4.5E-20   & 9.8E-14   & 9.6E-42   \\\nopagebreak
    &      & (8.4E-60) & (9.2E-54) & {\color{red}(3.5E-28)} & (6.9E-20) & (8.4E-14) & (2.7E-41) \\
    f2  & 2000 & 1.8E-25   & 3.9E-22   & 1.5E-23   & 1.9E-14   & 1.6E-09   & 9.3E-21   \\\nopagebreak
    &      & (8.8E-25) & (2.7E-21) & (1.0E-23) & (1.1E-14) & (1.1E-09) & (6.3E-20) \\
    f3  & 5000 & 5.7E-61   & 6.0E-87   & 5.2E-14   & 9.0E-37   & 6.6E-11   & 2.5E-19   \\\nopagebreak
    &      & (2.7E-60) & (1.9E-86) & (1.1E-13) & (5.4E-36) & (8.8E-11) & (3.9E-19) \\
    f4  & 5000 & 8.2E-24   & 4.3E-66   & 1.4E-15   & 7.4E-11   & 4.2E-01   & 4.4E-14   \\\nopagebreak
    &      & (4.0E-23) & (1.2E-65) & (1.0E-15) & (1.8E-10) & (1.1E+00) & (9.3E-14) \\
    f5  & 3000 & 8.0E-02   & 3.2E-01   & 1.3E+01   & 2.1E+01   & 2.1E+00   & 2.5E+01   \\\nopagebreak
    &      & (5.6E-01) & (1.1E+00) & (1.4E+01) & (7.8E+00) & (1.5E+00) & (3.2E+01) \\
    f6  & 100  & 2.9E+00   & 5.6E+00   & 1.0E+03   & 9.3E+02   & 4.7E+03   & 4.5E+01   \\\nopagebreak
    &      & (1.2E+00) & (1.6E+00) & (2.2E+02) & (1.8E+02) & (1.1E+03) & (2.4E+01) \\
    f7  & 3000 & 6.4E-04   & 6.8E-04   & 3.3E-03   & 4.8E-03   & 4.7E-03   & 2.5E-03   \\\nopagebreak
    &      & (2.5E-04) & (2.5E-04) & (8.5E-04) & (1.2E-03) & (1.2E-03) & (1.4E-03) \\
    f8  & 1000 & 3.3E-05   & 7.1E+00   & 7.9E-11   & 4.7E+00   & 5.9E+03   & 2.4E+03   \\\nopagebreak
    &      & (2.3E-05) & (2.8E+01) & (1.3E-10) & (3.3E+01) & (1.1E+03) & (6.7E+02) \\
    f9  & 1000 & 1.0E-04   & 1.4E-04   & 1.5E-04   & 1.2E-03   & 1.8E+02   & 5.2E+01   \\\nopagebreak
    &      & (6.0E-05) & (6.5E-05) & (2.0E-04) & (6.5E-04) & (1.3E+01) & (1.6E+01) \\
    f10 & 500  & 8.2E-10   & 3.0E-09   & 3.5E-04   & 2.7E-03   & 1.1E-01   & 4.6E-01   \\\nopagebreak
    &      & (6.9E-10) & (2.2E-09) & (1.0E-04) & (5.1E-04) & (3.9E-02) & (6.6E-01) \\
    f11 & 500  & 9.9E-08   & 2.0E-04   & 1.9E-05   & 7.8E-04  & 2.0E-01   & 1.3E-02   \\\nopagebreak
    &      & (6.0E-07) & (1.4E-03) & (5.8E-05) & (1.2E-03)  & (1.1E-01) & (1.7E-02) \\
    f12 & 500  & 4.6E-17   & 3.8E-16   & 1.6E-07   & 1.9E-05   & 1.2E-02   & 1.9E-01   \\\nopagebreak
    &      & (1.9E-16) & (8.3E-16) & (1.5E-07) & (9.2E-06) & (1.0E-02) & (3.9E-01) \\
    f13 & 500  & 2.0E-16   & 1.2E-15   & 1.5E-06   & 6.1E-05   & 7.5E-02   & 2.9E-03   \\\nopagebreak
    &      & (6.5E-16) & (2.8E-15) & (9.8E-07) & (2.0E-05) & (3.8E-02) & (4.8E-03) \\
    \bottomrule %%% нижняя линейка
\end{longtabu} \endgroup

\section{Стандартные префиксы ссылок}\label{app:B4}

Общепринятым является следующий формат ссылок: \texttt{<prefix>:<label>}.
Например, \verb+\label{fig:knuth}+; \verb+\ref{tab:test1}+; \verb+label={lst:external1}+.
В~таблице \cref{tab:tab_pref} приведены стандартные префиксы для различных
типов ссылок.

\begin{table}[htbp]
    \captionsetup{justification=centering}
    \centering{
        \caption{\label{tab:tab_pref}Стандартные префиксы ссылок}
        \begin{tabular}{ll}
            \toprule
            \textbf{Префикс} & \textbf{Описание} \\
            \midrule
            ch:              & Глава             \\
            sec:             & Секция            \\
            subsec:          & Подсекция         \\
            fig:             & Рисунок           \\
            tab:             & Таблица           \\
            eq:              & Уравнение         \\
            lst:             & Листинг программы \\
            itm:             & Элемент списка    \\
            alg:             & Алгоритм          \\
            app:             & Секция приложения \\
            \bottomrule
        \end{tabular}
    }
\end{table}


Для упорядочивания ссылок можно использовать разделительные символы.
Например, \verb+\label{fig:scheemes/my_scheeme}+ или \\ \verb+\label{lst:dts/linked_list}+.

\section{Очередной подраздел приложения}\label{app:B5}

Нужно больше подразделов приложения!

\section{И ещё один подраздел приложения}\label{app:B6}

Нужно больше подразделов приложения!

\clearpage
\refstepcounter{chapter}
\addcontentsline{toc}{appendix}{\protect\chapternumberline{\thechapter}Чертёж детали}

\includepdf[pages=-]{Dissertation/images/drawing.pdf}
        % Приложения

\setcounter{totalappendix}{\value{chapter}} % Подсчёт количества приложений

\end{document}
